\documentclass[11pt,a4paper]{article}
\usepackage[T1]{fontenc}
\usepackage[utf8]{inputenc}
\usepackage{lmodern,amsmath,amssymb,amsthm,mathtools,microtype,mathrsfs}
\usepackage[margin=25mm,top=23mm,bottom=25mm]{geometry}
\usepackage{booktabs,array,longtable,enumitem,xcolor,fancyhdr}
\usepackage[colorlinks=true,linkcolor=blue!45!black,urlcolor=blue!45!black,citecolor=blue!45!black]{hyperref}
\hypersetup{pdftitle={Reading the Alpoge--Claude RH and Dirichlet-GRH proof: working version 0.51},pdfauthor={ChatGPT},pdfsubject={Partial human-readable reconstruction with explicit arithmetic inputs},pdfkeywords={Riemann hypothesis, Lean, theta lifts, prime sums, Arakelov slopes}}
\setlength{\parindent}{0pt}
\setlength{\parskip}{5pt plus 1pt}
\setlist{topsep=3pt,itemsep=3pt,leftmargin=*}
\setlength{\headheight}{14pt}
\pagestyle{fancy}\fancyhf{}
\fancyhead[L]{\small RH and Dirichlet GRH: a reconstruction}
\fancyhead[R]{\small Working version 0.51}
\fancyfoot[C]{\thepage}
\renewcommand{\headrulewidth}{0.3pt}
\numberwithin{equation}{section}
\newtheorem{theorem}{Theorem}[section]
\newtheorem{proposition}[theorem]{Proposition}
\newtheorem{lemma}[theorem]{Lemma}
\newtheorem{corollary}[theorem]{Corollary}
\theoremstyle{remark}\newtheorem{remark}[theorem]{Remark}
\newcommand{\Rea}{\operatorname{Re}}
\newcommand{\HS}{\mathrm{HS}}
\newcommand{\tr}{\operatorname{tr}}
\newcommand{\codim}{\operatorname{codim}}
\newcommand{\covol}{\operatorname{covol}}
\newcommand{\OK}{\mathcal O_K}
\newcommand{\src}[2]{\href{https://github.com/CocoSgt/riemann-hypothesis-lean/blob/7cbf7a0c32fca1badf05931879a37c3f8a2b03b4/src/#1}{\texttt{#2}}}
\newcommand{\status}[1]{\textbf{Status:} #1}
\begin{document}
\begin{center}
{\LARGE\bfseries RH and Dirichlet GRH: new human proofs}\\[7pt]
{\large Version 0.51 --- 23 November 2026}\\[5pt]
{\large Uniformization, complete frames and Bessel evaluations}
\end{center}
\textbf{Scope.} This edition contains the current audit and all new chapters.
It is a partial human reconstruction of the accepted formal proof.
Earlier named claims are in the cumulative manuscript and preserved components.
Each result below states its original inputs and exact remaining producers.
\setcounter{tocdepth}{1}
\tableofcontents
\clearpage
\begingroup\emergencystretch=3em
\section{Current completion audit: uniformization, complete frames and evaluated periods}
\label{sec:dependency-audit-v51}

\subsection{Status of the human reconstruction}

We accept the pinned Alpoge--Claude formal proof as correct. The task
here is to supply ordinary mathematical proofs and to preserve their
exact interfaces. Version 0.51 adds 80 numbered claims to the
3112 preserved predecessor claims, for a cumulative inventory of
3192. The selected sufficient dependency graph still has 121 nodes
and 161 edges. There remain 24 whole open obligations: 17 class C
and 7 class B. Of these, 17 are shared by RH and unrestricted
Dirichlet GRH and 7 are additional Dirichlet-GRH inputs.
The counts describe the chosen decomposition, not a completion
percentage or a claim that the decomposition is minimal.

\subsection{The five new chapters}

\paragraph{NFL1--NFL22: the transformation laws and actual uniformization.}
The finite quadratic Gauss sum and Gaussian Poisson calculation prove
the exact squared theta multiplier. The lattice congruences and
normalization changes then give the original triangular weight law
and arithmetic invariance. A single witness is valid before every
allowed weight and natural rule is selected.

The holomorphic reading argument uses the actual algebraic frame
charts. Even for an arbitrary coefficient algebra, locally only
finitely many coefficients and nonvanishing denominators occur.
Naturality on a universal full-frame ring makes every rule evaluate
holomorphically. For the normalized section, $n=2m$ gives
\[
 \kappa_b^n=(C\det h_b)^m.
\]
Thus arbitrary pointwise square roots cause no obstruction; no
continuous root is required. These results complete all four
\texttt{Frames\allowbreak OK} fields on the family already constructed in NFR.

An explicit inverse to the period construction gives every actual
framed point. Its fiber calculation recovers the original arithmetic
group and triangular action, including equality of the component
indices. These prove the whole \texttt{Cover\allowbreak Tri} and \texttt{Fib\allowbreak Tri}
statements. The final join verifies every \texttt{Is\allowbreak Unf} field,
including the original section, period lattice and Gaussian mean one.
The same uniformization has principal level
$l.M'=\operatorname{lcm}(l.M,l.A)$ and the all-cusp return.
The supplied \texttt{Moduli}, exact \texttt{Sec\allowbreak Base\allowbreak Change\allowbreak St} and
\texttt{Sec\allowbreak Twist\allowbreak St}, and inherited classical foundations remain
premises. Their global suppliers are still separate work.

\paragraph{LFC1--LFC23: all conditions on one marked arithmetic frame.}
The seven conditions left in MFC are now proved:
\texttt{phi}, \texttt{gen}, \texttt{cl}, \texttt{ec},
\texttt{b1p}, \texttt{b1} and \texttt{h1}.
Kummer theory, class characters and local inertia detection supply
the generator and class conditions. The untwisted cohomology bound
has a shorter proof through the appropriate Kummer eigenspace of
$F(\mu_\ell)$: the CM-unit anti-invariants contribute only the
root-of-unity term, and the class-group term uses the existing
class-rank bound. No new Leopoldt premise is introduced.

For the Euler condition, the proof keeps the actual local subspaces
and continuous cochains. Local degree-two corrections, inflation
and the explicitly retained Poitou--Tate/Euler input give the
source cone estimate. Every result uses the same marked frame,
characters and constant $C_0$. The final assembly has all 28
\texttt{Ok} fields and the unchanged \texttt{Big} condition.
This proves full \texttt{Has\allowbreak Frame\allowbreak W} at the original general
local-ring scope under full \texttt{Gen} and
\texttt{GC.\allowbreak Hypotheses2}. In the original coefficient-DVR case,
the CAF character seed joins the same frame with its modulus
chosen first. The full \texttt{Gen\_\allowbreak gd}/\texttt{Gen\_\allowbreak host}
entry implications and later automorphy producers remain.

\paragraph{BST1--BST11: derivative decay, tensor maps and ordered minors.}
A compact conjugation-averaged smooth approximate identity reproduces
the finite lowest compact type after normalization. Moving right
derivatives onto this compactly supported kernel preserves the
original decay exponent. This proves the full derivative statement
and both actual tensor maps into the smooth $L^2$ representation.

The algebraic part proves alternation of the shifted minors and
equality of their row and column formulas. In the key induction,
every commutator contribution occurs equally often with opposite
sign, so the complete sum cancels. The resulting exact expansions
give the original complementary-minor eigenspaces and both tensor
Casimir shifts. All operator products keep their prescribed order.

\paragraph{BSE1--BSE9: both original Bessel period evaluations.}
The original Haar measures and character give the full equivariance
law. Differentiating the actual period, with the original frame,
produces the framed matrix-unit action. For $r>0$ the product of
its simple character values is
\[
 \kappa=-\frac{T_z}{4^r}\prod_{a=0}^{r-1}|c_a|^2\ne0.
\]
The ordered complementary minor reduces to the compact Capelli
polynomial. Its exact covariance then supplies explicit standard
and dual input vectors. Their periods are respectively
\[
 \kappa\prod_{j<k}(u_j-\lambda_s-\tfrac12)\,\ell(v),
 \qquad
 \kappa\prod_{j<k}(u_j-\lambda_s+\tfrac12)\,\ell(v).
\]
When $r=0$, the respective factors are
$-\sqrt{T_z}$ and $\sqrt{T_z}$ in place of $\kappa$.
These are the literal original \texttt{Eval\allowbreak P} and \texttt{Eval\allowbreak M},
with every free translate, coefficient slot and measure unchanged.

Only the two full quotient statements now remain in this Bessel
step join. They must construct the actual smooth image, shifted
compact type and period injectivity. Given \texttt{Quot\allowbreak P} and
\texttt{Quot\allowbreak M}, the earlier BNV/BSP results yield the full
\texttt{Step\allowbreak Statement} and \texttt{Bessel\allowbreak NVGap}; the setting
conclusion still retains its original \texttt{Forms0} input.

\paragraph{DFD1--DFD15: actual continuous families and canonical data.}
Semisimple character splitting occurs over the fraction field.
Intersecting its complementary subspace with the original integral
lattice gives a saturated free lattice with free quotient.
An integral basis and retraction give continuous integral matrices;
an integral splitting of the character summands is unnecessary.
The exact Frobenius and parity calculations then construct the
original continuous complement family and \texttt{Pt\allowbreak Fam}.

The proof verifies every field of \texttt{Hub\allowbreak Data}, \texttt{Data\allowbreak AD}
and \texttt{Canon}, and proves the entire original
\texttt{exists\_\allowbreak hub\allowbreak Data} at its named Chebotarev, automorphic and
comparison inputs. It supplies the actual point rows, the correctly
oriented Selmer-killing implication with budget $1/\ell^2$, and
the full coefficient-extension theorem from the original integral
Bhatt--Scholze input to crystalline lattices. Uniform large-scale
estimates give admissible parameters and every small-prime bound.

The final family setup still requires coherent Wach comparison
data and the actual Selmer signature comparison for that family.
These remain printed premises. The local rows must then provide
the whole \texttt{Concl\allowbreak F}; DNE already gives its complete
\texttt{Adm\allowbreak F} and \texttt{Link\allowbreak N} consequences with
$K=2000(|C_a|+1)^3$ chosen before the family.

\subsection{The exact frontier}

No additional whole selected obligation is closed in this edition.
Complete subtheorems are distinguished from their remaining
suppliers. In particular the full conditional uniformization does
not supply the global section laws, the complete conditional
marked frame does not supply its entry or automorphy inputs,
and the computed nonzero tensor periods do not construct their
required quotients. The complete v0.49
\texttt{GLOBAL\_\allowbreak LOCAL\_\allowbreak GAIN} result is preserved.

Each whole contract below is preserved from the previous selected
graph. A version-0.51 progress paragraph, where present, states
the current completed subresults and remaining scope. Historical
qualifications inside the unchanged contracts remain for
traceability. Every new claim has an independent mathematical
review with exact source and proof hashes. The graph records
the human argument; no new Lean replay or mechanically extracted
transitive Lean audit is asserted.

\paragraph{1. GLOBAL\_\allowbreak DENOMINATOR — C — Actual LinkN local table and binding producer}

Closed(Des $\mathcal C$ 3) LinkN: the actual etaSum is bounded with exact dZ0 and same-host vvTerm, subject only to the literal ShaFinOdd and summability guards.

\emph{Shared objects and choice order.} one FamP; one class $\mathcal A$ and its target-compatible subclass; one coefficient field per family datum; one x and i at each application; original selected rows, words and levels; thresholds chosen before the common family; actual D.nd; actual etaJ and Aj; sha\_\allowbreak eq

\emph{Completed inputs (version 0.51).} DFD1-DFD15 (full complement/canonical data and parameter/small-prime subtheorems; conditional final setup), DNE1-DNE14 (whole original ConclF -> AdmF -> LinkN consumer)

\emph{Remaining producer scope.} DFD constructs the original continuous complement family, PtFam, every field of HubData, DataAD and Canon, and the whole exists\_\allowbreak hubData theorem under Cheb, adm2, Conclusion41 and the supplied HW input. It supplies GoodPtRow, exact oriented KillAt at budget 1/ell\textasciicircum{}2, the BhattScholzeIntegral-to-CrystallineLattices coefficient theorem, uniform admissible parameters and the full eventual small-prime family bounds. The final exists\_\allowbreak setup\_\allowbreak large implication still retains coherent HW/Fon/Cmp for the actual canonical datum and the actual Selmer SignatureChoice comparison. Construct those common suppliers and then every original local row/table bound to obtain one ConclF family, retaining ShaFinOdd and vvTerm summability. DNE already proves the full ConclF-to-AdmF-to-LinkN consumer with K=2000(|Ca|+1)\textasciicircum{}3 chosen before the family. GLOBAL\_\allowbreak DENOMINATOR therefore remains open until its actual common ConclF producer is supplied.

\emph{Source.} \href{https://github.com/CocoSgt/riemann-hypothesis-lean/blob/7cbf7a0c32fca1badf05931879a37c3f8a2b03b4/src/Formalism/SharpRestriction/GlobalTheoremInputs.lean#L45-L48}{LinkN}.

\emph{Source.} \href{https://github.com/CocoSgt/riemann-hypothesis-lean/blob/7cbf7a0c32fca1badf05931879a37c3f8a2b03b4/src/Formalism/SharpRestriction/GlobalTheoremInputs.lean#L40-L44}{LinkN\_\allowbreak at}.

\emph{Source.} \href{https://github.com/CocoSgt/riemann-hypothesis-lean/blob/7cbf7a0c32fca1badf05931879a37c3f8a2b03b4/src/Development/AutomorphicArithmetic/CharacterBoxes/DivinNOfHypV.lean#L117-L134}{divinN\_\allowbreak of\_\allowbreak hypV}.



\paragraph{2. GLOBAL\_\allowbreak RETAINED\_\allowbreak DEFECT — C — Original retained-row defect bound LinkIRR}

Rows.DelpNonneg and delp$\le$exp(-x\textasciicircum{}(ez/7)) uniformly for every late actual row datum. CURRENT VERSION0.47 QUALIFICATION: Finish the same-D/common-family adelic setting after ASG: normalized nonzero Bessel W and decay, measured-field block, original Siegel-Weil rescaling and measure/character ties, Bessel uniqueness/local ties, then the original AutTie/RowDictTie/LabTie. RowsFinDim is now proved universally by RFD. Supply the actual CountIdempotentsClass theorem or its three sufficient arithmetic cylinder counts, with one simultaneous family threshold; reuse the completed ARO structure. Completed new original producers: ARO1-ARO30, RFD1-RFD9, ASG1-ASG18. Any earlier narrative naming these completed producers as outstanding is historical; this qualification is the current scope authority. CURRENT VERSION0.48 QUALIFICATION: Finish the same-D/common-family adelic setting after ASG and AMD: normalized nonzero Bessel W and decay, original orbital/fiber and Siegel-Weil measure/character ties, Bessel uniqueness/local ties, and the full original transport theorem, then AutTie/RowDictTie/LabTie. The complete MExists measured-data block and unipotent covolume normalization are now supplied, and universal RowsFinDim remains available. Supply the actual CountIdempotentsClass theorem or its three sufficient arithmetic cylinder counts with one simultaneous family threshold. Completed current producers or conditional consumers: ARO1-ARO30, RFD1-RFD9, ASG1-ASG18, AMD1-AMD17. Earlier frontier descriptions are historical; this qualification is the current scope authority. Original whole-contract premises and quantifiers are retained. CURRENT VERSION0.49 QUALIFICATION: Prove the exact original BesselNVGap nonzero integral on its supplied discrete-series, compact-representation and character data. CoefWeightedLoc, WlocExists with its original nonvanishing antecedent, WFromLoc, the normalized global product and all-later-measured-block decay are complete. On the inherited Forms0 scope this leaves BesselNVGap for ConcreteSettingExists. Continue original orbital/fiber and Siegel-Weil measure/character ties, the required Bessel uniqueness/local ties and the full original transport theorem, then AutTie/RowDictTie/LabTie. Supply the actual CountIdempotentsClass theorem or its three sufficient arithmetic cylinder counts with one simultaneous family threshold. Completed current producers or conditional consumers: ARO1-ARO30, RFD1-RFD9, ASG1-ASG18, AMD1-AMD17, BWP1-BWP24 at their exact BesselNVGap/Forms0 scope. Earlier frontier descriptions are historical; this qualification is the current scope authority. Original whole-contract premises and quantifiers are retained.

\emph{Shared objects and choice order.} one FamP; one class $\mathcal A$ and its target-compatible subclass; one coefficient field per family datum; one x and i at each application; original selected rows, words and levels; thresholds chosen before the common family; AllRows A; LabelCountGoodInjective; AllJoinA; CountAllRows

\emph{Completed inputs (version 0.51).} DFD common-family subresults at stated inputs, DNE full separate denominator consumer, WRS1-WRS10 full selected-class local gain, Inherited first DivDef3 and row-pair integrability

\emph{Remaining producer scope.} Complete the original F3 local coefficient bridges and F4 common-overfield readings, with valuation/norm/GoodCoef/iq and the entire coefficient field, and actual complex Hecke standard/no-Jordan plus arithmetic character data on the same form datum. Construct the actual base rows, host/frame lengths, nonzero chains and group ties, MultEq support, eigensystem/carrier/parts, all-embedding logLam, HD/d3b and NEb. Keep the stronger second DivRows3 scope distinct from the completed first DivDef3 and row-pair integrability. The completed conditional common marked frame and CountV character seed do not by themselves install HOST\_\allowbreak NUMBER\_\allowbreak S, HOST\_\allowbreak FRAME\_\allowbreak LENGTH, TASK\_\allowbreak HOST\_\allowbreak VAN\_\allowbreak E\_\allowbreak INSTALL, TASK\_\allowbreak HOST\_\allowbreak XPP\_\allowbreak INSTALL or TASK\_\allowbreak COUNT\_\allowbreak V\_\allowbreak INSTALL. Continue the distinct retained-defect, period-height and normalized spectral producers, and the original compatible second-layer selection on one common arithmetic/geometric family. GLOBAL\_\allowbreak LOCAL\_\allowbreak GAIN remains fully closed at its original v0.49 scope. The complete denominator consumer and the conditional frame/uniformization constructors do not by themselves construct this whole common family. Common design existence, NBUILT compatibility, original true-column records, shared fields/embeddings, Fits and eta agreement and eventual-family thresholds remain explicit. A stronger historical WPt/TSel construction may serve other producers but is not needed to reopen the proved selected-class local-gain estimate.

\emph{Source.} \href{https://github.com/CocoSgt/riemann-hypothesis-lean/blob/7cbf7a0c32fca1badf05931879a37c3f8a2b03b4/src/Formalism/SharpRestriction/GlobalTheoremInputs.lean#L55-L58}{LinkIRR}.

\emph{Source.} \href{https://github.com/CocoSgt/riemann-hypothesis-lean/blob/7cbf7a0c32fca1badf05931879a37c3f8a2b03b4/src/Development/AbelianVarietyTheory/Proofs/GlobalInterpolation/LinkIRROfJoinEx.lean#L17-L38}{linkIRR\_\allowbreak of\_\allowbreak join\_\allowbreak ex}.

\emph{Source.} \href{https://github.com/CocoSgt/riemann-hypothesis-lean/blob/7cbf7a0c32fca1badf05931879a37c3f8a2b03b4/src/Development/AbelianVarietyTheory/Proofs/GlobalInterpolation/LinkIRROfJoinEx.lean#L143-L166}{link\_\allowbreak irr}.



\paragraph{3. GLOBAL\_\allowbreak PERIOD\_\allowbreak HEIGHT — C — Full actual period and height comparison LinkDIVt}

Closed conditional TuSize -> DGamSize -> (LinkEta -> DIVinN -> LinkDIV), with every original LinkDIV\_\allowbreak at clause.

\emph{Shared objects and choice order.} one FamP; one class $\mathcal A$ and its target-compatible subclass; one coefficient field per family datum; one x and i at each application; original selected rows, words and levels; thresholds chosen before the common family; same RR target; same Bu/besS; same Nu and Bk

\emph{Source.} \href{https://github.com/CocoSgt/riemann-hypothesis-lean/blob/7cbf7a0c32fca1badf05931879a37c3f8a2b03b4/src/Formalism/SharpRestriction/GlobalTheoremInputs.lean#L77-L90}{LinkDIV\_\allowbreak at}.

\emph{Source.} \href{https://github.com/CocoSgt/riemann-hypothesis-lean/blob/7cbf7a0c32fca1badf05931879a37c3f8a2b03b4/src/Formalism/SharpRestriction/GlobalTheoremInputs.lean#L94-L96}{LinkDIVt}.

\emph{Source.} \href{https://github.com/CocoSgt/riemann-hypothesis-lean/blob/7cbf7a0c32fca1badf05931879a37c3f8a2b03b4/src/Development/AbelianVarietyTheory/Proofs/GlobalInterpolation/LinkBrest.lean#L7-L25}{GlobalTheory.Top.link\_\allowbreak brest}.

\emph{Source.} \href{https://github.com/CocoSgt/riemann-hypothesis-lean/blob/7cbf7a0c32fca1badf05931879a37c3f8a2b03b4/src/Development/AbelianVarietyTheory/Proofs/GlobalInterpolation/WireSel.lean#L11-L41}{GlobalTheory.Top.wire\_\allowbreak sel}.



\paragraph{4. GLOBAL\_\allowbreak SPECTRAL — C — Original normalized kernel/Gram estimates LinkS}

Closed TuSize -> LinkS with exact Gamma normalization, off-diagonal Hilbert-Schmidt bound, diagonal error and Dall/Dstar estimate, errors exp(-x\textasciicircum{}(epsilon/8)). CURRENT VERSION0.48 QUALIFICATION: Complete the original same-family spectral/Bessel and Siegel-Weil inputs with all measure, character, transport and common-object ties. The complete MExists measured block and prescribed-measure covolume normalization are supplied; these are not the full normalized spectral theorem. Completed current producers or conditional consumers: ASG1-ASG18, AMD1-AMD17, GCF1-GCF18 at their stated scopes. Earlier frontier descriptions are historical; this qualification is the current scope authority. Original whole-contract premises and quantifiers are retained. CURRENT VERSION0.49 QUALIFICATION: Complete the original same-family spectral/Bessel and Siegel-Weil inputs with every measure, character, transport and common-object tie. The normalized packet and decay are now constructed conditional on the exact BesselNVGap, and the analytic setting uses the original supplied Forms0 scope. These do not imply the full normalized spectral theorem. Completed current producers or conditional consumers: ASG1-ASG18, AMD1-AMD17, GCF1-GCF18 at their stated scopes, BWP1-BWP24 at their exact BesselNVGap/Forms0 scope, WRS1-WRS10 original local gain. Earlier frontier descriptions are historical; this qualification is the current scope authority. Original whole-contract premises and quantifiers are retained.

\emph{Shared objects and choice order.} one FamP; one class $\mathcal A$ and its target-compatible subclass; one coefficient field per family datum; one x and i at each application; original selected rows, words and levels; thresholds chosen before the common family; actual KD; same S.Cu spaces; GamOf; Dstar; Dall

\emph{Completed inputs (version 0.51).} BST1-BST11 (full derivative/tensor/minor/eigenspace suppliers), BSE1-BSE8 (full framed period laws and both evaluations), BSE9 conditional on only original QuotP/QuotM, BNV1-BNV23, BSP1-BSP10, BWP1-BWP24 at original nonvanishing/Forms0 scope, AMD1-AMD24

\emph{Remaining producer scope.} Complete the full original QuotPStatement and QuotMStatement on the actual tensor coefficient maps and period kernel. BST supplies DecDeriv, both TMap/TMapD statements, the noncommutative complementary-minor algebra, EigStatement and exact tensor Casimir shifts. BSE supplies full PerEquiv, FrameLaw, CapelliB/WhittakerZero, CapelliCov and both original EvalP/EvalM with explicit witnesses and exact nonzero scalars. The quotient suppliers must still produce literal W\_\allowbreak f and W-prime, the exact image relation, shifted lowest compact type and Casimir, actual compact-group stability and period injectivity. BSE9 then gives full StepStatement and BesselNVGap, with Forms0 still retained for ConcreteSettingExists. The actual orbital/fiber and Siegel-Weil character/measure comparisons, Bessel uniqueness/local ties, full transport, AutTie/RowDictTie/LabTie and simultaneous arithmetic cylinder counts for CountIdempotentsClass remain separate.

\emph{Source.} \href{https://github.com/CocoSgt/riemann-hypothesis-lean/blob/7cbf7a0c32fca1badf05931879a37c3f8a2b03b4/src/Formalism/SharpRestriction/GlobalTheoremInputs.lean#L145-L149}{LinkS\_\allowbreak at}.

\emph{Source.} \href{https://github.com/CocoSgt/riemann-hypothesis-lean/blob/7cbf7a0c32fca1badf05931879a37c3f8a2b03b4/src/Formalism/SharpRestriction/GlobalTheoremInputs.lean#L152-L157}{LinkS}.

\emph{Source.} \href{https://github.com/CocoSgt/riemann-hypothesis-lean/blob/7cbf7a0c32fca1badf05931879a37c3f8a2b03b4/src/Development/AbelianVarietyTheory/Proofs/GlobalInterpolation/LinkAlt.lean#L16-L50}{GlobalTheory.Top.link\_\allowbreak alt}.



\paragraph{5. AUT\_\allowbreak UNRAM\_\allowbreak COEF — C — Actual unramified-coefficient automorphy}

For each actual Dat with coefficient ramification one and fixed iota, choose one marked g1 supplying FinA, AutB and LemW together. CURRENT VERSION0.48 QUALIFICATION: Construct the actual AutomorphyUnramifiedCoefficients package: for every original Dat with coefficient ramification one and fixed iota, choose one marked g1 that supplies FinA, AutB and LemW together. The original InpZ record, full RigII, CNR3, finite PropGK consumer and now the entire OneChar implication are supplied. Continue RH Cases I/V and the independent additional Dirichlet-GRH CaseIII/promodularity, SeedIV, tensor-factor and rank-four-chain tasks. Completed current producers or conditional consumers: OCH1-OCH30 (the full consumer implication), KIZ original InpZ, KTR1-KTR13, RNC16, PGK1-PGK6. Earlier frontier descriptions are historical; this qualification is the current scope authority. Original whole-contract premises and quantifiers are retained. CURRENT VERSION0.49 QUALIFICATION: Construct one common GdEnd for each original datum and fixed iota: the same marked g1 must supply FinEnd, AutEnd3 and AutEnd2. UCE now proves all ordinary local/curve/coefficient consumers and the exact implication to FinA/AutB/LemW, including arbitrary entry maps and every original outer-corner extension. The producer still needs the actual common abelian-CM GdFrameEx and LinkedEndE with both rank-two and rank-three runs, full ModUAlt/PfromAut and the two actual seeds, followed by the original local/automorphic descent. At this pin the alternative upstream route gives GdEndAlt; its additional local/descent fields must be matched to the actual original points, so no name-only GdEndAlt-to-GdEnd identification is allowed. Keep the supplied InpZ, RigII, CNR3, finite PropGK and whole OneChar implication. Continue RH Cases I/V and the independent additional Dirichlet-GRH CaseIII/promodularity, SeedIV, tensor-factor and rank-four-chain tasks. Completed current producers or conditional consumers: UCE1-UCE19 (complete GdEnd consumer), OCH1-OCH30, KIZ original InpZ, KTR1-KTR13, RNC16, PGK1-PGK6. Earlier frontier descriptions are historical; this qualification is the current scope authority. Original whole-contract premises and quantifiers are retained.

\emph{Completed inputs (version 0.51).} LFC1-LFC23 (all 28 Ok fields and full HasFrameW at full Gen/P; original coefficient-DVR SeedClause join), CAF1-CAF11, MFC1-MFC14, UCE1-UCE19 (full original GdEnd consumer), OCH1-OCH30, KIZ original InpZ, KTR1-KTR13, RNC16, PGK1-PGK6

\emph{Remaining producer scope.} LFC completes phi/gen/cl/ec/b1p/b1/h1 and all 28 Ok fields on the exact marked frame, with the same C0, Big, characters and coefficient embedding. It gives full HasFrameW under the original full Gen and GC.Hypotheses2/P, and joins the CAF SeedClause to HasFrame in the original coefficient-DVR case with its modulus chosen first. The full entry implications Gen\_\allowbreak gd and Gen\_\allowbreak host remain, followed by FinEnd, AutEnd3 and AutEnd2, the actual rank-two and rank-three automorphic seeds, full ModUAlt/PfromAut and local/automorphic descent to GdEnd. UCE supplies the complete GdEnd consumer; it does not replace these producers. The alternative GdEndAlt route retains its exact additional local/descent fields. The independent CountV HostData character-seed scope still includes ell=3. Do not weaken it to the final GD.gen guard.

\emph{Source.} \href{https://github.com/CocoSgt/riemann-hypothesis-lean/blob/7cbf7a0c32fca1badf05931879a37c3f8a2b03b4/src/Development/AutomorphicArithmetic/AutomorphicVarieties/AutomorphicCases/AutomorphyUnramifiedCoefficients.lean#L24-L27}{AutomorphyUnramifiedCoefficients}.

\emph{Source.} \href{https://github.com/CocoSgt/riemann-hypothesis-lean/blob/7cbf7a0c32fca1badf05931879a37c3f8a2b03b4/src/Development/AutomorphicArithmetic/AutomorphicVarieties/AutomorphicCases/AutomorphyUnramifiedCoefficients.lean#L12-L13}{FinA}.

\emph{Source.} \href{https://github.com/CocoSgt/riemann-hypothesis-lean/blob/7cbf7a0c32fca1badf05931879a37c3f8a2b03b4/src/Development/AutomorphicArithmetic/AutomorphicVarieties/AutomorphicCases/AutomorphyUnramifiedCoefficients.lean#L15-L17}{AutB}.

\emph{Source.} \href{https://github.com/CocoSgt/riemann-hypothesis-lean/blob/7cbf7a0c32fca1badf05931879a37c3f8a2b03b4/src/Development/AutomorphicArithmetic/AutomorphicVarieties/AutomorphicCases/AutomorphyUnramifiedCoefficients.lean#L19-L22}{LemW}.



\paragraph{6. TASK\_\allowbreak CASE3\_\allowbreak AUT\_\allowbreak JOIN — B — Complete the original Case-III automorphy/image join}

Install the actual Case-III ProModAlt, SeedIVAlt and shared PMExAlt with the inherited CohomologyHyp/Ferrero-Washington/HFit/SeedII facts into the original whole Case-III automorphy/image conclusion.

\emph{Source.} \href{https://github.com/CocoSgt/riemann-hypothesis-lean/blob/7cbf7a0c32fca1badf05931879a37c3f8a2b03b4/src/Development/AutomorphicArithmetic/AutomorphicVarieties/ModularityArgument/ThresholdArguments/DenseThresholdCases/PromodularityCaseThree.lean#L43-L44}{AutomorphyRankFourDense}.

\emph{Source.} \href{https://github.com/CocoSgt/riemann-hypothesis-lean/blob/7cbf7a0c32fca1badf05931879a37c3f8a2b03b4/src/Development/AutomorphicArithmetic/AutomorphicVarieties/ModularityArgument/ThresholdArguments/DenseThresholdCases/CuFinal.lean#L44-L44}{main4\_\allowbreak alt}.



\paragraph{7. TASK\_\allowbreak CASE3\_\allowbreak PROMOD\_\allowbreak JOIN — C — Reconstruct the actual Case-III chain-to-promodularity proof}

Use the original tensor-factor tower, ChainRankFour, ChainDesc4, seeded rank-four chain reduction, actual PMExAlt and SeedIVAlt to construct ProModAlt on every original Ctx3 star point at the same fixed iota, preserving all finite-coefficient, inertia and support maps.

\emph{Source.} \href{https://github.com/CocoSgt/riemann-hypothesis-lean/blob/7cbf7a0c32fca1badf05931879a37c3f8a2b03b4/src/Development/AutomorphicArithmetic/AutomorphicVarieties/ModularityArgument/ThresholdArguments/DenseThresholdCases/PromodularityCaseThreeOfLemmas.lean#L131-L153}{UpperTheory.V.Cu.promodularity\_\allowbreak case\_\allowbreak three\_\allowbreak of\_\allowbreak chain\_\allowbreak reduction}.

\emph{Source.} \href{https://github.com/CocoSgt/riemann-hypothesis-lean/blob/7cbf7a0c32fca1badf05931879a37c3f8a2b03b4/src/Development/AutomorphicArithmetic/AutomorphicVarieties/ModularityArgument/ThresholdArguments/DenseThresholdCases/PromodularityCaseThree.lean#L40-L42}{PromodularityCaseThreeDense}.



\paragraph{8. SEED\_\allowbreak IV\_\allowbreak ALT — C — Actual Case-III seed with prechosen residual datum}

Base Ctx3/Four tower, iota and Amb -> w,s before every later Dies tower -> compatible finite-coefficient seed with exact residual support/root entries.

\emph{Source.} \href{https://github.com/CocoSgt/riemann-hypothesis-lean/blob/7cbf7a0c32fca1badf05931879a37c3f8a2b03b4/src/Development/AutomorphicArithmetic/AutomorphicVarieties/ModularityArgument/ThresholdArguments/DenseThresholdCases/PromodularityCaseThree.lean#L19-L25}{SeedIVAlt}.



\paragraph{9. TENSOR\_\allowbreak FACTOR\_\allowbreak TOWER — B — Actual tensor-factor tower and characters}

Original finite-field, tame, factor and character constraints; source-level construction is sizeable and has not been reconstructed here.

\emph{Source.} \href{https://github.com/CocoSgt/riemann-hypothesis-lean/blob/7cbf7a0c32fca1badf05931879a37c3f8a2b03b4/src/Development/AutomorphicArithmetic/AutomorphicVarieties/ModularityArgument/ThresholdArguments/DenseThresholdCases/PromodularityCaseThreeOfLemmas.lean#L131-L153}{UpperTheory.V.Cu.promodularity\_\allowbreak case\_\allowbreak three\_\allowbreak of\_\allowbreak chain\_\allowbreak reduction}.



\paragraph{10. CHAIN\_\allowbreak RANK\_\allowbreak FOUR — C — Actual SomeIn-to-AllIn propagation}

Same source Space and PM ideal, residual point Good4, original bound 30(C+1)\textasciicircum{}2 <= d; geometry must preserve all identities.

\emph{Source.} \href{https://github.com/CocoSgt/riemann-hypothesis-lean/blob/7cbf7a0c32fca1badf05931879a37c3f8a2b03b4/src/Development/AutomorphicArithmetic/AutomorphicVarieties/ModularityArgument/ThresholdArguments/DenseThresholdCases/PromodularityCaseThree.lean#L26-L32}{ChainRankFour}.



\paragraph{11. CHAIN\_\allowbreak DESC4 — B — Original residual chain descent}

Finite residual descent retaining source Chain4 compatibility.

\emph{Source.} \href{https://github.com/CocoSgt/riemann-hypothesis-lean/blob/7cbf7a0c32fca1badf05931879a37c3f8a2b03b4/src/Development/AutomorphicArithmetic/AutomorphicVarieties/ModularityArgument/ThresholdArguments/DenseThresholdCases/PromodularityCaseThreeOfLemmas.lean#L131-L153}{UpperTheory.V.Cu.promodularity\_\allowbreak case\_\allowbreak three\_\allowbreak of\_\allowbreak chain\_\allowbreak reduction}.



\paragraph{12. RANK4\_\allowbreak CHAIN\_\allowbreak REDUCTION — C — Seeded reduction to the original residual chain}

Select T1 before every later Deep tower; retain one compatible seed and Chain4 for every original residual lift over Q.

\emph{Source.} \href{https://github.com/CocoSgt/riemann-hypothesis-lean/blob/7cbf7a0c32fca1badf05931879a37c3f8a2b03b4/src/Development/AutomorphicArithmetic/AutomorphicVarieties/ModularityArgument/ThresholdArguments/DenseThresholdCases/PromodularityCaseThree.lean#L33-L39}{RankFourChainReductionDense}.

\emph{Source.} \href{https://github.com/CocoSgt/riemann-hypothesis-lean/blob/7cbf7a0c32fca1badf05931879a37c3f8a2b03b4/src/Development/AutomorphicArithmetic/AutomorphicVarieties/ModularityArgument/ThresholdArguments/DenseThresholdCases/CuPChainIV.lean#L16-L82}{UpperTheory.V.Cu.chainIV\_\allowbreak alt}.



\paragraph{13. TASK\_\allowbreak COUNT\_\allowbreak V\_\allowbreak INSTALL — B — Complete the original CountV finite assembly}

From exactly the listed source host/cohomology inputs, complete countV\_\allowbreak of\_\allowbreak finiteness\_\allowbreak duality\_\allowbreak euler\_\allowbreak host, including original host/frame formation, normalized length, the CaseV guard and R1<=2\textasciicircum{}c2.

\emph{Completed inputs (version 0.51).} LFC1-LFC23 (conditional full common frame), CAF9-CAF11 (full CountV HostData character seed including ell=3), MFC1-MFC14, Inherited CountV consumers

\emph{Remaining producer scope.} LFC completes phi/gen/cl/ec/b1p/b1/h1 and all 28 Ok fields on the exact marked frame, with the same C0, Big, characters and coefficient embedding. It gives full HasFrameW under the original full Gen and GC.Hypotheses2/P, and joins the CAF SeedClause to HasFrame in the original coefficient-DVR case with its modulus chosen first. The full entry implications Gen\_\allowbreak gd and Gen\_\allowbreak host remain, followed by FinEnd, AutEnd3 and AutEnd2, the actual rank-two and rank-three automorphic seeds, full ModUAlt/PfromAut and local/automorphic descent to GdEnd. UCE supplies the complete GdEnd consumer; it does not replace these producers. The alternative GdEndAlt route retains its exact additional local/descent fields. The independent CountV HostData character-seed scope still includes ell=3. Do not weaken it to the final GD.gen guard. Complete the original F3 local coefficient bridges and F4 common-overfield readings, with valuation/norm/GoodCoef/iq and the entire coefficient field, and actual complex Hecke standard/no-Jordan plus arithmetic character data on the same form datum. Construct the actual base rows, host/frame lengths, nonzero chains and group ties, MultEq support, eigensystem/carrier/parts, all-embedding logLam, HD/d3b and NEb. Keep the stronger second DivRows3 scope distinct from the completed first DivDef3 and row-pair integrability. The completed conditional common marked frame and CountV character seed do not by themselves install HOST\_\allowbreak NUMBER\_\allowbreak S, HOST\_\allowbreak FRAME\_\allowbreak LENGTH, TASK\_\allowbreak HOST\_\allowbreak VAN\_\allowbreak E\_\allowbreak INSTALL, TASK\_\allowbreak HOST\_\allowbreak XPP\_\allowbreak INSTALL or TASK\_\allowbreak COUNT\_\allowbreak V\_\allowbreak INSTALL.

\emph{Source.} \href{https://github.com/CocoSgt/riemann-hypothesis-lean/blob/7cbf7a0c32fca1badf05931879a37c3f8a2b03b4/src/Development/AutomorphicArithmetic/AutomorphicVarieties/HeightAndVanishing/Counts.lean#L69-L69}{CountV}.

\emph{Source.} \href{https://github.com/CocoSgt/riemann-hypothesis-lean/blob/7cbf7a0c32fca1badf05931879a37c3f8a2b03b4/src/Development/AutomorphicArithmetic/AutomorphicVarieties/ModularityArgument/ThresholdArguments/DenseThresholdCases/CountVOf.lean#L33-L60}{UpperTheory.V.UnitaryThreeInertia.countV\_\allowbreak of\_\allowbreak finiteness\_\allowbreak duality\_\allowbreak euler\_\allowbreak host}.



\paragraph{14. THM\_\allowbreak S6\_\allowbreak S — C — Original unitary-three ordinary support and automorphy}

Original host/frame and fixed iota -> one marked g1 -> every universal R,u -> weight-zero finiteness and every Type-I BPrime\_\allowbreak s point.

\emph{Source.} \href{https://github.com/CocoSgt/riemann-hypothesis-lean/blob/7cbf7a0c32fca1badf05931879a37c3f8a2b03b4/src/Development/AutomorphicArithmetic/AutomorphicVarieties/ModularityArgument/ThresholdArguments/DenseThresholdCases/ThmS6CoreS.lean#L17-L26}{ThmS6\_\allowbreak s}.

\emph{Source.} \href{https://github.com/CocoSgt/riemann-hypothesis-lean/blob/7cbf7a0c32fca1badf05931879a37c3f8a2b03b4/src/Development/AutomorphicArithmetic/AutomorphicVarieties/ModularityArgument/ThresholdArguments/DenseThresholdCases/ThmS6CoreS.lean#L12-L16}{BPrime\_\allowbreak s}.



\paragraph{15. HOST\_\allowbreak FRAME\_\allowbreak LENGTH — B — Original Case-V host/frame and normalized-length installation}

Actual HostData, frame, local-condition equality, abstract count and normalized original Selmer length must all refer to the same objects.

\emph{Completed inputs (version 0.51).} LFC full conditional common-frame construction, CAF and MFC subresults, DNE weight/host-degree dictionary at supplied tables

\emph{Remaining producer scope.} Complete the original F3 local coefficient bridges and F4 common-overfield readings, with valuation/norm/GoodCoef/iq and the entire coefficient field, and actual complex Hecke standard/no-Jordan plus arithmetic character data on the same form datum. Construct the actual base rows, host/frame lengths, nonzero chains and group ties, MultEq support, eigensystem/carrier/parts, all-embedding logLam, HD/d3b and NEb. Keep the stronger second DivRows3 scope distinct from the completed first DivDef3 and row-pair integrability. The completed conditional common marked frame and CountV character seed do not by themselves install HOST\_\allowbreak NUMBER\_\allowbreak S, HOST\_\allowbreak FRAME\_\allowbreak LENGTH, TASK\_\allowbreak HOST\_\allowbreak VAN\_\allowbreak E\_\allowbreak INSTALL, TASK\_\allowbreak HOST\_\allowbreak XPP\_\allowbreak INSTALL or TASK\_\allowbreak COUNT\_\allowbreak V\_\allowbreak INSTALL.

\emph{Source.} \href{https://github.com/CocoSgt/riemann-hypothesis-lean/blob/7cbf7a0c32fca1badf05931879a37c3f8a2b03b4/src/Development/AutomorphicArithmetic/AutomorphicVarieties/ModularityArgument/ThresholdArguments/DenseThresholdCases/CountVOf.lean#L33-L60}{UpperTheory.V.UnitaryThreeInertia.countV\_\allowbreak of\_\allowbreak finiteness\_\allowbreak duality\_\allowbreak euler\_\allowbreak host}.



\paragraph{16. HOST\_\allowbreak NUMBER\_\allowbreak S — C — Uniform original-host rank-three catalogue}

One exponent before all closures, coefficient rings, levels, characters, frames, embeddings and marked points; one separating trace function controls all later finite pairwise nonisomorphic collections.

\emph{Completed inputs (version 0.51).} LFC full conditional common-frame construction, CAF and MFC subresults at original inputs, Inherited host-count consumers

\emph{Remaining producer scope.} Complete the original F3 local coefficient bridges and F4 common-overfield readings, with valuation/norm/GoodCoef/iq and the entire coefficient field, and actual complex Hecke standard/no-Jordan plus arithmetic character data on the same form datum. Construct the actual base rows, host/frame lengths, nonzero chains and group ties, MultEq support, eigensystem/carrier/parts, all-embedding logLam, HD/d3b and NEb. Keep the stronger second DivRows3 scope distinct from the completed first DivDef3 and row-pair integrability. The completed conditional common marked frame and CountV character seed do not by themselves install HOST\_\allowbreak NUMBER\_\allowbreak S, HOST\_\allowbreak FRAME\_\allowbreak LENGTH, TASK\_\allowbreak HOST\_\allowbreak VAN\_\allowbreak E\_\allowbreak INSTALL, TASK\_\allowbreak HOST\_\allowbreak XPP\_\allowbreak INSTALL or TASK\_\allowbreak COUNT\_\allowbreak V\_\allowbreak INSTALL.

\emph{Source.} \href{https://github.com/CocoSgt/riemann-hypothesis-lean/blob/7cbf7a0c32fca1badf05931879a37c3f8a2b03b4/src/Development/AutomorphicArithmetic/AutomorphicVarieties/ModularityArgument/ThresholdArguments/DenseThresholdCases/HostNumberS.lean#L18-L28}{HostNumber\_\allowbreak s}.



\paragraph{17. TASK\_\allowbreak HOST\_\allowbreak VAN\_\allowbreak E\_\allowbreak INSTALL — B — Transport K/duality to literal HostVanE}

Use the already supplied LemmaK/local-global duality and finiteness on the actual CountV host/frame and exact original coefficient modules.

\emph{Source.} \href{https://github.com/CocoSgt/riemann-hypothesis-lean/blob/7cbf7a0c32fca1badf05931879a37c3f8a2b03b4/src/Development/AutomorphicArithmetic/AutomorphicVarieties/UnitaryThreeDeformations/HostNumber.lean#L9-L12}{HostVanE}.

\emph{Source.} \href{https://github.com/CocoSgt/riemann-hypothesis-lean/blob/7cbf7a0c32fca1badf05931879a37c3f8a2b03b4/src/Development/AutomorphicArithmetic/AutomorphicVarieties/LevelCohomology/SoHostVanE.lean#L11-L18}{UpperTheory.V.X.So.hostVanE}.



\paragraph{18. TASK\_\allowbreak HOST\_\allowbreak XPP\_\allowbreak INSTALL — B — Transport CorX to literal HostXpp}

Preserve Psi\textasciicircum{}-2 eta, the prime-to-ell quotient, actual local conditions and the original CountV host while installing BR7.

\emph{Source.} \href{https://github.com/CocoSgt/riemann-hypothesis-lean/blob/7cbf7a0c32fca1badf05931879a37c3f8a2b03b4/src/Development/AutomorphicArithmetic/AutomorphicVarieties/UnitaryThreeDeformations/HostNumber.lean#L14-L17}{HostXpp}.

\emph{Source.} \href{https://github.com/CocoSgt/riemann-hypothesis-lean/blob/7cbf7a0c32fca1badf05931879a37c3f8a2b03b4/src/Development/AutomorphicArithmetic/AutomorphicVarieties/LevelCohomology/Xpp.lean#L11-L84}{UpperTheory.V.X.xpp}.



\paragraph{19. DESIGN\_\allowbreak CHOICES\_\allowbreak COMMON — C — Construct the original common design family and bundle}

For the one target HostExW, construct the original compatible DesignChoices on Cdes=WithTw(WithSm Cprime), where Cprime=WithChN(A intersect target class), A is the one Built construction class, and cT(eps,y)=cTmin(y). Retain every original field, tuple, word, metric and common bundle duty. DSZ10 and RMP11 give size/rank properties of class members, not member existence. The original RecExP/RecAtG and common object/row route remains coupled. Construct ONE original common design family R and its recursive/geometric half with original Inputs/Pins/Holds/ExAt/InGeo and required CoreP. PKT now reconstructs full exists\_\allowbreak pkT, rowPk\_\allowbreak nonempty and exists\_\allowbreak moved with Copy/Coh on those original given inputs and a supplied fitting later stage: one catalogue, both ChT, complete LfT, actual integral lattices and all metric/measure/period constants are retained; one numerical record precedes every eta and V/S/N/G. RDR supplies one Side0 before all readers, ReadPk/pins and full ReadL with exact o.eta=q.eta. The initial actual geometry, Fits/division/central conditions, primitive rational tests and eta agreement must still be produced. Inherited CSR supplies finite columns, one ReadC before later twists, period/evaluation and scalar/error transports, RowSideDetermined and same-R Bundle/BTw/BWit/BunU consumers. RPI supplies exact RowPairInt for actual pairs and finite translates with original mu\_\allowbreak b/fd0; first literal DivDef3 at the same updated R,F,Y,Z,C,sEE follows from its retained original readings, Pins/Holds, stored nonzero IsItW chain, five real-group predicates and group tie. Remaining common fields include host/frame, MultEq nonzero-part support, that original nonzero chain/group tie, eigensystem/carrier/parts and all-embedding logLam, HD/d3b, NEb, actual rational tests and base rows. Retain every simultaneous family hypothesis and eventual threshold. Stronger second DivRows3 remains separate; universal quotient-height CuspRapidDecay and actual common member existence are not inferred from the comparison constructors. Version 0.41 supplies the finite phase-correct rational theta route: exact coefficient extension, split representatives, square-law/PhiIs, Pins from CoversAt, class-sum/point-identity transfer to OneMapOverL, actual-good-projector RatG and both-level LiftTestB. Its original higher covers, all-tuple class-sum uniqueness including phase rescaling, one PointIdentCore before O, base ClassSumsSpan and geometric LatFull remain explicit producers. MLat is available from the original LatFull input. This sharpens the rational-test work without claiming actual common geometry, Fits or eta agreement. Version0.46 proves finite generation of the original whole integral cusp lattice by actual integral-frame evaluation. It also proves the exact identity-cover, analytic unfolding, original-field coefficient descent, global clearing-scalar and both-level lattice-transport implications. Full spanning retains the supplied original IsNab connection and the five explicit principal-chart, level-system, raw-Gaussian, short-rule/dictionary and ramified-jet producers. The higher tuple covers, actual common class sums, PointIdentCore and ClassSumsSpan remain separate. CURRENT VERSION0.47 QUALIFICATION: Construct one original common geometric family and the higher CoversAt/all-tuple class sums and uniqueness with original phase rescalings, one PointIdentCore with common lambda,z, and base ClassSumsSpan. CNC now supplies the full original IsNab; PCL supplies the conditional trace-dual principal-lattice and quantitative normalization. Actual uniformization/HasPrinc and the five principal-chart, finite-level-system, Gaussian/reproduction, short-rule/simultaneous-dictionary and ramified-jet producer groups remain. Keep all original both-level, coefficient-field, Fits and eta agreements. Completed new original producers: ICL1-ICL21, CNC1-CNC16, PCL1-PCL12. Any earlier narrative naming these completed producers as outstanding is historical; this qualification is the current scope authority. CURRENT VERSION0.48 QUALIFICATION: Construct the original finite holomorphic FamilyUnif/N1c, normalized theta families and local lifting, then the common principal charts and finite-level theta systems. PUN now supplies the exact fixed-space models, all-family center lattices, A=1 exponent matching and individual primitive points with the full point-equality criterion; CNC supplies the full canonical connection and PCL the quantitative trace-dual normalization. Complete higher CoversAt and all-tuple class sums/uniqueness with the original phase rescalings, one PointIdentCore with common lambda,z, and base ClassSumsSpan. The Gaussian/reproduction, short-rule/simultaneous-dictionary and ramified-jet producers, and all original both-level, coefficient-field, Fits and eta agreements remain explicit. Completed current producers or conditional consumers: ICL1-ICL21, CNC1-CNC16, PCL1-PCL12, PUN1-PUN24, GCF1-GCF18 at their stated scopes. Earlier frontier descriptions are historical; this qualification is the current scope authority. Original whole-contract premises and quantifiers are retained. CURRENT VERSION0.49 QUALIFICATION: Construct the adjoining simultaneous principal framed points and test-ring dictionary. Full finite holomorphic FamilyUnif/N1c, normalized theta families and local lifts are now supplied by HUF. The next bounded target is allTest\_\allowbreak of\_\allowbreak secBaseChange: represent FrPt by the Lie-eigenframe bundles and theta-section frame line, use the full n=1 SecBaseChangeSt on every nonflat coefficient base change, and prove finite-generated subalgebra descent with the exact TRing/FromTest transport. Then supply the full FramesOK/CoverTri/FibTri and IsUnf/hasPrinc package, including adapted mod-eight symplectic basis, normalized theta multiplier/weight and arithmetic invariance. Continue higher CoversAt, all-tuple class sums and uniqueness with original phase rescalings, one PointIdentCore with common lambda,z, base ClassSumsSpan, Gaussian/reproduction, short-rule dictionaries and ramified jets, retaining both-level, coefficient-field, Fits and eta agreements. Common-family nonemptiness remains distinct from compatibility of an already supplied member. Completed current producers or conditional consumers: ICL1-ICL21, CNC1-CNC16, PCL1-PCL12, PUN1-PUN24, HUF1-HUF21 (full original FamilyUnif/N1c), GCF1-GCF18 and WRS1-WRS10 at their stated scopes. Earlier frontier descriptions are historical; this qualification is the current scope authority. Original whole-contract premises and quantifiers are retained.

\emph{Completed inputs (version 0.51).} NFL1-NFL22 (full FramesOK, CoverTri/FibTri and conditional IsUnf/principal return), TDR1-TDR13, NFR1-NFR13, HUF1-HUF21, PUN1-PUN24, ICL1-ICL21, CNC1-CNC16, PCL1-PCL12, GCF/WRS at original supplied-member scope

\emph{Remaining producer scope.} NFL now proves all four original FramesOK fields, CoverTri, FibTri and full IsUnf, with the principal-level/all-cusp return on the same NFR family. It retains the supplied Moduli and the exact global SecBaseChangeSt and SecTwistSt laws and inherited HUF/PUN classical foundations; their global suppliers remain explicit. Do not reopen the normalized theta multiplier, holomorphic readings, actual covering/fiber construction or arbitrary-ring test descent. Next supply the original global section laws and compatible moduli inputs, then the finite level system, Gaussian/reproduction, short-rule/simultaneous dictionaries and ramified jets. Higher CoversAt, all-tuple class sums and uniqueness with phase rescalings, one PointIdentCore with common lambda and z, and base ClassSumsSpan remain. Preserve both levels, the entire coefficient field, Fits and eta agreements, and the original common-family quantifier order. Common design nonemptiness is still required.



\paragraph{20. LOCAL\_\allowbreak NBUILT\_\allowbreak COMMON — C — Actual local specialization on the common selected family}

Construct, with the required closed-class property on the same Des(C,3), one eventual family of actual NBuilt certificates for the original Fm.D(x,i), retaining the certificate before forgetting to TSel. HLP now reconstructs full HeckeLift0, all-index commutation, HeckeStd0 and HeckeIntTw from the original MtInput and complex Hecke facts, in the full ambient GenAt range k<=nV. Every maximal ideal is an actual supported residual character class; literal MtSub is the O-span of projections of the original MtLat generators, is finite free, and localizes onto the exact NCov.part. Literal Psi has its exact range and off-space projection error, both original span assertions hold, and ActEx/ActHK are constructed on the same lattice under original semisimplicity. Original MjLat relation quotients, universal coefficient-map relation fields, selected tame generators, coefficient embeddings and all identity transports give the complete nonempty\_\allowbreak spec constructor on its full printed input list. A single compositum combines the four original coefficient requirements with big fixed beforehand. Whole ndSpec\_\allowbreak exists still needs actual common geometric/choice MtInput, complex HeckeStdPlain and the no-Jordan producer, original F3 bridges with valuation/GoodCoef/iq, original F4 common-overfield reading of the entire coefficient field, and exact arithmetic character/reciprocity/admissibility suppliers. Original AutSpec, Galois specialization and the remaining NBuilt fields on that same installation also remain. Do not replace the actual projected lattice or maps by abstract objects with the same dimensions.

\emph{Completed inputs (version 0.51).} NFL full conditional complex uniformization, LFC full conditional common frame, TDR/NFR and CAF/MFC, WRS selected built-class witness extraction at membership

\emph{Remaining producer scope.} Continue the distinct retained-defect, period-height and normalized spectral producers, and the original compatible second-layer selection on one common arithmetic/geometric family. GLOBAL\_\allowbreak LOCAL\_\allowbreak GAIN remains fully closed at its original v0.49 scope. The complete denominator consumer and the conditional frame/uniformization constructors do not by themselves construct this whole common family. Common design existence, NBUILT compatibility, original true-column records, shared fields/embeddings, Fits and eta agreement and eventual-family thresholds remain explicit. A stronger historical WPt/TSel construction may serve other producers but is not needed to reopen the proved selected-class local-gain estimate. NFL now proves all four original FramesOK fields, CoverTri, FibTri and full IsUnf, with the principal-level/all-cusp return on the same NFR family. It retains the supplied Moduli and the exact global SecBaseChangeSt and SecTwistSt laws and inherited HUF/PUN classical foundations; their global suppliers remain explicit. Do not reopen the normalized theta multiplier, holomorphic readings, actual covering/fiber construction or arbitrary-ring test descent. Next supply the original global section laws and compatible moduli inputs, then the finite level system, Gaussian/reproduction, short-rule/simultaneous dictionaries and ramified jets. Higher CoversAt, all-tuple class sums and uniqueness with phase rescalings, one PointIdentCore with common lambda and z, and base ClassSumsSpan remain. Preserve both levels, the entire coefficient field, Fits and eta agreements, and the original common-family quantifier order.



\paragraph{21. SECOND\_\allowbreak LAYER\_\allowbreak COMPATIBLE\_\allowbreak SELECTION — C — Original second-layer selector and class duties}

For the same A, supply Duty on A and a Closed common original TSel family whose selected records are the original true-column rec7 records (or have the entire dependent tuple transported), with their real witness and original coefficient-field root data. Existence/compatibility remains arithmetic; merely abstractly isomorphic fields or independently selected columns do not suffice. CSR21–22 now construct the same-R Bundle/BTw/BWit/BunU outputs from the exact selected realization/row-side packet and NEb, with the original EvP/ParEv quantifier order. Actual common selection, the five stated realization packets and NEb remain upstream; none is replaced by an independently chosen compatible record. Version 0.38 closes the selected analytic extension and reconstructs the first updated DivDef3 consumer; it leaves the original common selection, geometric readings, rational tests, base rows, host/parts/NEb and exact retained DI premises upstream, with unchanged EvP/ParEv thresholds. Version 0.39 supplies the original given-data reader constructors and exact RowPairInt/first DivDef3 on original premises. It leaves one actual common geometric/comparison/test member, eta agreement, compatible selection, rational tests, base rows, host/parts/NEb and original EvP/ParEv thresholds upstream. Version 0.40 also supplies original comparison/RowPk/fitting-stage Moved algorithms and exact common data transport. The actual initial geometric member, fitting tests, eta agreement and compatible second-layer selector with original eventual thresholds remain upstream. Version 0.41 supplies the finite phase-correct rational theta route: exact coefficient extension, split representatives, square-law/PhiIs, Pins from CoversAt, class-sum/point-identity transfer to OneMapOverL, actual-good-projector RatG and both-level LiftTestB. Its original higher covers, all-tuple class-sum uniqueness including phase rescaling, one PointIdentCore before O, base ClassSumsSpan and geometric LatFull remain explicit producers. MLat is available from the original LatFull input. This sharpens the rational-test work without claiming actual common geometry, Fits or eta agreement. CURRENT VERSION0.49 QUALIFICATION: Continue the distinct global denominator, period/height, normalized spectral and second-layer compatible-selection producers, together with existence of the original common arithmetic/geometric family. GLOBAL\_\allowbreak LOCAL\_\allowbreak GAIN and its actual GroupCompatibilityFamily supplier are now complete on the unchanged selected built class; no additional WPt/TSel realization is needed for that estimate. Do not reopen local first-growth, mean, auxiliary admission, selected D3 tops/Adm2F, the single GIs2 stage or the exact comparison/kernel transport. A full RecT or TSel may remain useful for other tasks, but is not claimed by this local-gain closure. Completed current producers or conditional consumers: GCF1-GCF18 at their stated scopes, WRS1-WRS10 original selected-class compatibility/local gain. Earlier frontier descriptions are historical; this qualification is the current scope authority. Original whole-contract premises and quantifiers are retained.

\emph{Completed inputs (version 0.51).} WRS1-WRS10 original selected-class compatibility/local gain, DNE full denominator consumer, DFD canonical-data subresults, NFL and LFC conditional constructors, Inherited exact same-record realization/transport consumers

\emph{Remaining producer scope.} Continue the distinct retained-defect, period-height and normalized spectral producers, and the original compatible second-layer selection on one common arithmetic/geometric family. GLOBAL\_\allowbreak LOCAL\_\allowbreak GAIN remains fully closed at its original v0.49 scope. The complete denominator consumer and the conditional frame/uniformization constructors do not by themselves construct this whole common family. Common design existence, NBUILT compatibility, original true-column records, shared fields/embeddings, Fits and eta agreement and eventual-family thresholds remain explicit. A stronger historical WPt/TSel construction may serve other producers but is not needed to reopen the proved selected-class local-gain estimate.



\paragraph{22. BASE\_\allowbreak THETA\_\allowbreak READINGS — C — Actual base reading on the common bundle}

For every family in the original class, eventually on the SAME selected Bundle and for each i<m, supply the actual R,F,Read,ReadL, weak models/Weil data and phase-one pins. Identify their LiftTie.theta, mtRead, mjRead and ThDef3 basis with that Bundle.P.A3, P.Zs, P.S0 and its stored maps. Exact field, embeddings, metrics, measures, row family and weak choices are retained. This remains one part of the coupled BTH12 producer. PKT now supplies the complete original comparison packet, RowPk and fitting-stage Moved with Copy/Coh on original given geometry, Inputs/Pins/Holds/ExAt/InGeo and Fits. Its ChT/LfT/Reps, actual lattice carriers, original Haar and norm constants, period maps, representatives and all-place RrT are constructed together; its one numerical record precedes every eta/V/S/N/G. RDR gives the original level-zero side, moved-Read and ReadPk/pin constructors on those objects, with one Side0 before every reader and all twelve Read ties. Full ReadL retains the exact o.eta(j,v)=q.eta(j,v) guard for every 1<=j<=m and place v; arbitrary Over does not imply it. Inherited CSR fixes one ReadC before all TwArgs, transports both metrics and all thirteen reading maps, and gives actual period/evaluation with kappa=(c*nP)\textasciicircum{}(-1/2). RPI supplies exact row-pair integrability, so first updated DivDef3 follows from the retained original Pins/Holds, nonzero chain, five real-group predicates and group tie. Initial theta extension remains complete. Remaining work is actual initial common geometry and its original predicates, Fits/primitive rational tests, eta agreement and outer family/eventual/same-Bundle duties, reFace compatibility, MultEq support, original chain/group data, actual base rows and any stronger DivRows3 theta square. No new compatible reading, family threshold or Haar normalization is substituted. Version 0.41 supplies the finite phase-correct rational theta route: exact coefficient extension, split representatives, square-law/PhiIs, Pins from CoversAt, class-sum/point-identity transfer to OneMapOverL, actual-good-projector RatG and both-level LiftTestB. Its original higher covers, all-tuple class-sum uniqueness including phase rescaling, one PointIdentCore before O, base ClassSumsSpan and geometric LatFull remain explicit producers. MLat is available from the original LatFull input. This sharpens the rational-test work without claiming actual common geometry, Fits or eta agreement.

\emph{Completed inputs (version 0.51).} NFL1-NFL22, TDR1-TDR13, NFR1-NFR13, HUF1-HUF21, PUN1-PUN24, CNC1-CNC16, Inherited actual-reading consumers at original inputs

\emph{Remaining producer scope.} NFL now proves all four original FramesOK fields, CoverTri, FibTri and full IsUnf, with the principal-level/all-cusp return on the same NFR family. It retains the supplied Moduli and the exact global SecBaseChangeSt and SecTwistSt laws and inherited HUF/PUN classical foundations; their global suppliers remain explicit. Do not reopen the normalized theta multiplier, holomorphic readings, actual covering/fiber construction or arbitrary-ring test descent. Next supply the original global section laws and compatible moduli inputs, then the finite level system, Gaussian/reproduction, short-rule/simultaneous dictionaries and ramified jets. Higher CoversAt, all-tuple class sums and uniqueness with phase rescalings, one PointIdentCore with common lambda and z, and base ClassSumsSpan remain. Preserve both levels, the entire coefficient field, Fits and eta agreements, and the original common-family quantifier order. Complete the original F3 local coefficient bridges and F4 common-overfield readings, with valuation/norm/GoodCoef/iq and the entire coefficient field, and actual complex Hecke standard/no-Jordan plus arithmetic character data on the same form datum. Construct the actual base rows, host/frame lengths, nonzero chains and group ties, MultEq support, eigensystem/carrier/parts, all-embedding logLam, HD/d3b and NEb. Keep the stronger second DivRows3 scope distinct from the completed first DivDef3 and row-pair integrability. The completed conditional common marked frame and CountV character seed do not by themselves install HOST\_\allowbreak NUMBER\_\allowbreak S, HOST\_\allowbreak FRAME\_\allowbreak LENGTH, TASK\_\allowbreak HOST\_\allowbreak VAN\_\allowbreak E\_\allowbreak INSTALL, TASK\_\allowbreak HOST\_\allowbreak XPP\_\allowbreak INSTALL or TASK\_\allowbreak COUNT\_\allowbreak V\_\allowbreak INSTALL.



\paragraph{23. BASE\_\allowbreak THETA\_\allowbreak RATIONAL\_\allowbreak TESTS — C — Simultaneous original base rational tests}

On the SAME original Bundle/reading/theta map, for every i<m supply one finite separating complex-linear test family whose values lie in the fixed embedded coefficient field on theta of EVERY original rational source vector and on EVERY stored target-coordinate basis vector. BTH8/BTH11 use this same LiftTestB for rational adjoint image and matrix coordinates. RTG now reconstructs prescribed-field coefficient extension, original split class representatives/RepsOK, the self-dual split-lattice square law and finite residue PhiIs; every Pins field follows from the actual CoversAt. With the base coefficient record literally fixed, exact class-sum existence/uniqueness including phase-rescaled O, PointIdentCore with common lambda,z before every O, and base ClassSumsSpan imply all-tuple finite lift span and OneMapOverL. The phase is exactly cj(twAt embLL A3.A.v0)(R.phs i). RatG uses the actual good projector eOp(y), and full LiftTestB follows at BOTH original levels. These finite implication proofs are complete. Remaining original producers are the prescribed higher covers, those actual class sums/uniqueness, PointIdentCore, base ClassSumsSpan and the base/common geometric datum; exact MLat is supplied only from original LatFull, whose geometric production remains separate. Finite dimension or an unrelated rational theta model cannot replace these same-object identities. Original outer eventual-family/Bundle quantifiers remain. Version0.46 proves finite generation of the original whole integral cusp lattice by actual integral-frame evaluation. It also proves the exact identity-cover, analytic unfolding, original-field coefficient descent, global clearing-scalar and both-level lattice-transport implications. Full spanning retains the supplied original IsNab connection and the five explicit principal-chart, level-system, raw-Gaussian, short-rule/dictionary and ramified-jet producers. The higher tuple covers, actual common class sums, PointIdentCore and ClassSumsSpan remain separate. CURRENT VERSION0.47 QUALIFICATION: Construct one original common geometric family and the higher CoversAt/all-tuple class sums and uniqueness with original phase rescalings, one PointIdentCore with common lambda,z, and base ClassSumsSpan. CNC now supplies the full original IsNab; PCL supplies the conditional trace-dual principal-lattice and quantitative normalization. Actual uniformization/HasPrinc and the five principal-chart, finite-level-system, Gaussian/reproduction, short-rule/simultaneous-dictionary and ramified-jet producer groups remain. Keep all original both-level, coefficient-field, Fits and eta agreements. Completed new original producers: ICL1-ICL21, CNC1-CNC16, PCL1-PCL12. Any earlier narrative naming these completed producers as outstanding is historical; this qualification is the current scope authority. CURRENT VERSION0.48 QUALIFICATION: Construct the original finite holomorphic FamilyUnif/N1c, normalized theta families and local lifting, then the common principal charts and finite-level theta systems. PUN now supplies the exact fixed-space models, all-family center lattices, A=1 exponent matching and individual primitive points with the full point-equality criterion; CNC supplies the full canonical connection and PCL the quantitative trace-dual normalization. Complete higher CoversAt and all-tuple class sums/uniqueness with the original phase rescalings, one PointIdentCore with common lambda,z, and base ClassSumsSpan. The Gaussian/reproduction, short-rule/simultaneous-dictionary and ramified-jet producers, and all original both-level, coefficient-field, Fits and eta agreements remain explicit. Completed current producers or conditional consumers: ICL1-ICL21, CNC1-CNC16, PCL1-PCL12, PUN1-PUN24. Earlier frontier descriptions are historical; this qualification is the current scope authority. Original whole-contract premises and quantifiers are retained. CURRENT VERSION0.49 QUALIFICATION: Construct the adjoining simultaneous principal framed points and test-ring dictionary. Full finite holomorphic FamilyUnif/N1c, normalized theta families and local lifts are now supplied by HUF. The next bounded target is allTest\_\allowbreak of\_\allowbreak secBaseChange: represent FrPt by the Lie-eigenframe bundles and theta-section frame line, use the full n=1 SecBaseChangeSt on every nonflat coefficient base change, and prove finite-generated subalgebra descent with the exact TRing/FromTest transport. Then supply the full FramesOK/CoverTri/FibTri and IsUnf/hasPrinc package, including adapted mod-eight symplectic basis, normalized theta multiplier/weight and arithmetic invariance. Continue higher CoversAt, all-tuple class sums and uniqueness with original phase rescalings, one PointIdentCore with common lambda,z, base ClassSumsSpan, Gaussian/reproduction, short-rule dictionaries and ramified jets, retaining both-level, coefficient-field, Fits and eta agreements. Completed current producers or conditional consumers: ICL1-ICL21, CNC1-CNC16, PCL1-PCL12, PUN1-PUN24, HUF1-HUF21 (full original FamilyUnif/N1c and theta family). Earlier frontier descriptions are historical; this qualification is the current scope authority. Original whole-contract premises and quantifiers are retained.

\emph{Completed inputs (version 0.51).} NFL1-NFL22, TDR1-TDR13, NFR1-NFR13, HUF1-HUF21, PUN1-PUN24, ICL1-ICL21, CNC1-CNC16, PCL1-PCL12, Inherited both-level RatG/LiftTestB consumers

\emph{Remaining producer scope.} NFL now proves all four original FramesOK fields, CoverTri, FibTri and full IsUnf, with the principal-level/all-cusp return on the same NFR family. It retains the supplied Moduli and the exact global SecBaseChangeSt and SecTwistSt laws and inherited HUF/PUN classical foundations; their global suppliers remain explicit. Do not reopen the normalized theta multiplier, holomorphic readings, actual covering/fiber construction or arbitrary-ring test descent. Next supply the original global section laws and compatible moduli inputs, then the finite level system, Gaussian/reproduction, short-rule/simultaneous dictionaries and ramified jets. Higher CoversAt, all-tuple class sums and uniqueness with phase rescalings, one PointIdentCore with common lambda and z, and base ClassSumsSpan remain. Preserve both levels, the entire coefficient field, Fits and eta agreements, and the original common-family quantifier order.



\paragraph{24. SLR\_\allowbreak BASE\_\allowbreak ROWS\_\allowbreak PRODUCER — C — Actual base row identities on that same bundle}

For the SAME selected Bundle, common original fields and readings used by the other three BTH12 inputs, supply actual RowsHold0\_\allowbreak 3 uniformly eventually for every family. Its original primitive row realization, spectrum/type and basis identifications must be retained. BTH supplies a theta consumer only; it does not construct the base row identity. No separate witness may be chosen merely to satisfy this row requirement.



\clearpage
\section{The transformation laws and the actual complex uniformization}
\label{sec:normalized-frame-laws-v051}

\status{NFL1--NFL22 prove the remaining three fields of
\texttt{Frames\allowbreak OK}, the original \texttt{Cover\allowbreak Tri} and \texttt{Fib\allowbreak Tri},
and the full \texttt{Is\allowbreak Unf} and principal/cusp return on the same family
constructed in NFR13.  The supplied \texttt{Moduli} record and the exact
\texttt{Sec\allowbreak Base\allowbreak Change\allowbreak St} and \texttt{Sec\allowbreak Twist\allowbreak St} laws remain premises.
No whole selected RH/GRH obligation is declared closed.  All source
identifications refer to commit
\texttt{7cbf7a0c32fca1badf05931879a37c3f8a2b03b4}.}

We retain all the original data and conventions of NFR1--NFR13:
$E,H,D_w,l,X_o,c$, $\sigma,g,e,F,r,M$, with \texttt{Good\allowbreak Emb},
\texttt{Model\allowbreak Base} and \texttt{Neg\allowbreak Xi}.  Thus $F$ is the actual
\texttt{Family\allowbreak Unif} of the supplied moduli problem, and $M$ its model
on the prescribed Hermitian space.  Fix a component $k$ temporarily.
Write $\Lambda=F.D_i(k).\mathrm{La}$, $G=M.X_k.G_m$,
$\xi_0=M.X_k.\xi_0$, $\phi=M.\phi_k$, and

\[
 E_\Lambda(v,w)=\operatorname{Tr}_{E/\mathbb Q}
       \bigl(-\xi_0\operatorname{herm}_G(\phi v,\phi w)\bigr).
 \tag{NFL.1}
\]

Hermitian forms are linear in the first argument.  The first $q$
columns of a frame have positive sign, and the remaining $r-q$
columns negative sign.  Let $\mathcal D$ be the actual open flag
domain, $T$ the \texttt{Theta\allowbreak Data} of NFR4, and
$b=T.\mathrm{co}^{-1}x$ its linear coordinates.  Its period map,
normalizer, Siegel matrix and rigidified local frames are denoted
$P_b,h_b,\tau_b,fr_j(b,z)$.  Its two fixed half-characteristics are
$\alpha,\beta\in\{0,1/2\}^g$.  The actual normalized point
$u_k(x)$ has section

\[
 s_b(z)=\kappa_b\theta[\alpha,\beta](\tau_b,h_bz)fr_j(b,z),
 \qquad \kappa_b^2=C\det h_b,
 \quad C=\operatorname{normConst}(\mathrm{ty},\xi_0)\ne0.
 \tag{NFL.2}
\]

The square root is chosen separately at each point.  No continuity
of $b\mapsto\kappa_b$ is assumed.  NFR12 already proves every field
of \texttt{Frame\allowbreak OK} for these exact points, including its algebraic
section and mean-one condition.

Here are the retained laws with their actual scope.  The
\texttt{Sec\allowbreak Base\allowbreak Change\allowbreak St} law starts with a locally Noetherian scheme
$S$, a polarized abelian variety $P/S$, and $n\ge1$.  Near each
point it supplies $n^{\dim P}$ sections whose pullbacks form a basis
after \emph{every} affine base change $\operatorname{Spec}R\to S$;
$R$ need not be reduced, Noetherian, or flat over the original base.
We use the $n=1$ instance of the actual principally polarized object
attached to $X_o$.  The \texttt{Sec\allowbreak Twist\allowbreak St} law, for every $S,P,q\ge1$
and every local ring $R$ with $t:\operatorname{Spec}R\to S$ and
$c\in P(R)$, supplies a basis of $q^{\dim P}$ sections of
$L^{q-1}\otimes t_c^*L$ and lifts every residue-field section.
The arguments below use its original complex-fiber instances.
Neither law is replaced by an assertion only for flat extensions.

We use HUF21 for the actual holomorphic family, its exponential,
lattice fibers and theta frames; PUN11 for its models; and TDR1--TDR13
for literal dual-number eigenframes, section generators, universal
frame rings and descent.  Their already stated classical foundations
are retained.  The additional analytic ingredients below are ordinary
Gaussian integration, Poisson summation for Schwartz functions,
Cauchy's theorem, Liouville's theorem and local analytic coordinates
on smooth complex algebraic varieties.  The needed theta multiplier
is proved, rather than made an extra input.

\subsection{The square of the theta multiplier}

% BEGIN NFL1
\paragraph{NFL1: the finite quadratic sum with its exact sign.}
\label{claim:NFL1}
Let $\gamma=\left(\begin{smallmatrix}A&B\\ C&D\end{smallmatrix}\right)
\in\operatorname{Sp}_{2g}(\mathbb Z)$, with $B,C$ even and $\det D$
odd.  On $Q=\mathbb Z^g/D\mathbb Z^g$ put
\[
 S(B,D)=\sum_{v\in Q}\exp\bigl(\pi i v^tBD^{-1}v\bigr).
\]
Then the summand is well defined and
\[
 S(B,D)^2=\chi_4(\det D)\det D,
 \qquad
 \chi_4(d)=\begin{cases}1&d\equiv1\pmod4,\\-1&d\equiv3\pmod4.
 \end{cases}
 \tag{NFL.3}
\]
The assertion includes $g=0$, with all determinants and sums equal
to one.

\begin{proof}
Symplecticity gives $BD^{-1}=(BD^{-1})^t$ and $D^tB$ symmetric.
Replacing $v$ by $v+Dn$ changes $v^tBD^{-1}v$ by
$2n^tB^tv+n^tD^tBn$, an even integer.  The resulting quadratic
function $q(v)=\tfrac12v^tBD^{-1}v\pmod{\mathbb Z}$ has associated
pairing $b(v,w)=v^tBD^{-1}w\pmod{\mathbb Z}$.
If this pairing kills $w$, then $BD^{-1}w$ is integral.  Both
$B(D^{-1}w)$ and $D(D^{-1}w)$ are integral, so
$(A^tD-C^tB)D^{-1}w=D^{-1}w$ is integral.  Hence $w=0$ in $Q$.
Thus $q$ is nondegenerate and $|Q|=|\det D|$ is odd.

For completeness, an odd finite abelian group with a nondegenerate
symmetric pairing into $\mathbb Q/\mathbb Z$ splits orthogonally
into cyclic prime-power groups with nondegenerate self-pairing.
First split its primary parts.  On a $p$-part choose an element of
largest order $p^e$.  Nondegeneracy pairs some two elements with
denominator $p^e$.  If neither has self-pairing with that denominator,
their sum does: its self-pairing contains twice their pairing, and
$p$ is odd.  The cyclic subgroup it generates is nondegenerate.
Subtracting the appropriate multiple of that generator projects
every element onto its orthogonal complement.  Induction proves
the assertion.  Since $2$ is invertible on an odd group, $q$ is
determined by $b(v,v)/2$; thus the Gauss sum factors across this
orthogonal decomposition.

On a cyclic summand it is
$G(a,p^e)=\sum_{x\bmod p^e}\exp(2\pi i ax^2/p^e)$ with $p\nmid a$.
For $e\ge2$, write $x=y+tp^{e-1}$ and sum over $t\bmod p$.
The inner sum vanishes unless $p\mid y$, and gives
$G(a,p^e)=pG(a,p^{e-2})$ otherwise.  For $e=1$, if $\chi$ is the
quadratic character of $\mathbb F_p$, counting square roots gives
$G(a,p)=\chi(a)\tau$, where
$\tau=\sum_x\chi(x)e^{2\pi ix/p}$.
Character orthogonality gives $|\tau|^2=p$: expand the square,
put $x=ty$ for $y\ne0$, and use
$\sum_{y\ne0}e^{2\pi i(t-1)y/p}=p-1$ for $t=1$ and $-1$ otherwise.
Also $\overline\tau=\chi(-1)\tau$.  Therefore
$G(a,p)^2=\chi(-1)p$, and the recurrence yields
$G(a,p^e)^2=\chi_4(p)^e p^e$.
Multiplication over all summands gives
$S^2=\chi_4(|\det D|)|\det D|$.
This equals the right side of (NFL.3) also when $\det D<0$,
because $\chi_4(-1)=-1$.
\end{proof}
% END NFL1

% BEGIN NFL2
\paragraph{NFL2: the exact theta-square transformation.}
\label{claim:NFL2}
For a Siegel matrix $\tau$, put
\[
 J=C\tau+D,\qquad \tau'=(A\tau+B)J^{-1},
 \qquad Q_\gamma(z)=z^tJ^{-1}Cz.
\]
If $B,C$ are divisible by four, then, for integer vectors $k,k',m,m'$
with $m\equiv A^tk\pmod2$ and $m'\equiv D^tk'\pmod2$,
\[
 \begin{split}
 \theta[k/2,k'/2](\tau',J^{-t}z)^2
   ={}&\chi_4(\det D)\det J\,e^{2\pi iQ_\gamma(z)}\\
     &{}\cdot\theta[m/2,m'/2](\tau,z)^2.
 \end{split}
 \tag{NFL.4}
\]
In particular this proves the original zero-characteristic
\texttt{theta\_\allowbreak sq\_\allowbreak transform\_\allowbreak level\_\allowbreak eight} law.

\begin{proof}
We give the multiplier computation.  For $\operatorname{Im}T>0$ let
$\mathcal G(T)=\int_{\mathbb R^g}e^{\pi it^tTt}\,dt$.
Real Gaussian integration and holomorphic continuation on the
connected Siegel space give
$\mathcal G(T)^2\det T=i^g$.
The Fourier transform of $e^{\pi it^tTt+2\pi it^tz}$ is
$\mathcal G(T)e^{-\pi i(z-m)^tT^{-1}(z-m)}$ at $m\in\mathbb Z^g$.
Poisson summation therefore gives
\[
 \theta[0,0](\tau,z)=\mathcal G(\tau)
 e^{-\pi iz^t\tau^{-1}z}
 \theta[0,0](-\tau^{-1},\tau^{-1}z).
 \tag{NFL.5}
\]
The Gaussian decay permits the Fourier transform, summation and
holomorphic continuation just used; derivatives on compact subsets
are bounded by a polynomial times a decaying real Gaussian.

Since $AD^t-BC^t=I$ and $B,C$ are even, $D$ is invertible modulo
two, so $\det D$ is odd and nonzero.  Set
$N=\tau^{-1}+D^{-1}C$.  The real symmetric matrix $D^{-1}C$
does not change the negative definite imaginary part of $\tau^{-1}$;
thus $N^{-1}$ is Siegel.  Block multiplication gives
\[
 J=DN\tau,\quad
 \tau'=D^{-t}N^{-1}D^{-1}+BD^{-1},\quad
 J^{-t}z=D^{-t}N^{-1}\tau^{-1}z.
 \tag{NFL.6}
\]
Decompose the series defining the theta function at $\tau'$ into
the cosets $v+D\mathbb Z^g$.  Complete the square in each coset
and use (NFL.5).  The finite Fourier sum that occurs is
\[
 \sum_{v\bmod D\mathbb Z^g}
 e^{\pi iv^tBD^{-1}v+2\pi im^tD^{-1}v}
   =S(B,D)e^{\pi im^tD^{-1}Cm}.
 \tag{NFL.7}
\]
Here replacing $v$ by $v-Cm$ proves the identity: use
$BC^t=AD^t-I$, the symmetries of $BD^{-1}$ and $D^{-1}C$,
and that $m^tA^tCm$ is even.  All remaining terms are integral
multiples of $2\pi i$.  Substitution into the coset sum yields
\[
 \theta[0,0](\tau',J^{-t}z)e^{-\pi iQ_\gamma(z)}
   =\lambda\theta[0,0](\tau,z),\qquad
 \lambda=S(B,D)\frac{\mathcal G(N^{-1})}{\mathcal G(\tau)}.
 \tag{NFL.8}
\]
Equations (NFL.3), (NFL.6) and the squared Gaussian formula give
$\lambda^2=\chi_4(\det D)\det J$.  This computes the sign without
choosing a square root of $\det J$.

To include the characteristics, use the series identity
\[
 \theta[a,b](\tau,z)=
 e^{\pi ia^t\tau a+2\pi ia^t(z+b)}
 \theta[0,0](\tau,z+\tau a+b).
 \tag{NFL.9}
\]
Put $n=A^tk+C^tk'$ and $n'=B^tk+D^tk'$.
Insert (NFL.9) into (NFL.8) with the translated argument and
expand its quadratic exponent.  The terms depending on $z$ give
$e^{\pi iQ_\gamma(z)}$; the remaining characteristic factor before
squaring is $e^{\pi i(k\cdot k'-n\cdot n')/4}$.
Indeed the required shift equality follows by multiplying by
$J^t$ and using $A^tD-C^tB=I$; expansion of the constant
quadratic terms then leaves $(k\cdot k'-n\cdot n')/4$ in the
exponent $\pi i$.  The symplectic identity also gives
\[
 k\cdot k'=(A^tk)\cdot(D^tk')-(B^tk)\cdot(C^tk').
\]
Consequently $k\cdot k'-n\cdot n'$ is divisible by four when
$B,C$ are.  Its squared exponential is one.  Changing an integer
representative of a half-characteristic multiplies theta by a
sign, so leaves its square unchanged.  Since
$n\equiv m$ and $n'\equiv m'$ modulo two, (NFL.4) follows.
\end{proof}
% END NFL2

% BEGIN NFL3
\paragraph{NFL3: a rigidified unit fixes the characteristic and multiplier square.}
\label{claim:NFL3}
Suppose $B,C$ are divisible by eight, $\chi_4(\det D)=1$, and
$\rho:\mathbb C^g\to\mathbb C^*$ is entire with $\rho(0)=1$.
If $\lambda\ne0$ and
\[
 \theta[\alpha,\beta](\tau',J^{-t}z)\rho(z)
      =\lambda\theta[\alpha,\beta](\tau,z),
 \tag{NFL.10}
\]
then $A^t(2\alpha)\equiv2\alpha\pmod2$,
$D^t(2\beta)\equiv2\beta\pmod2$, and $\lambda^2=\det J$.
This holds even when the theta function vanishes at zero.

\begin{proof}
Write $n=2\alpha,n'=2\beta$ and
$r=\lambda/\rho$.  This is an entire nowhere-zero function and
$\theta(\tau',J^{-t}z)=r(z)\theta(\tau,z)$.
For $\ell=\tau m+m'$ set
$p=Dm-Cm'$ and $p'=Am'-Bm$; then
$J^{-t}\ell=\tau'p+p'$.
The translation formula proved in NFR5, applied on the nonzero
set of theta and then extended by continuity, shows that
$s(z)=r(z)e^{-\pi iQ_\gamma(z)}$ satisfies
\[
 s(z+\ell)=(-1)^{L(m,m')}s(z),\quad
 L=n\cdot p'-n'\cdot p-n\cdot m'+n'\cdot m
          -m\cdot m'+p\cdot p'.
 \tag{NFL.11}
\]
One can check the quadratic part directly from
\[
 \begin{split}
 Q_\gamma(z+\ell)-Q_\gamma(z)
   ={}&2m\cdot z+m\cdot\tau m
       -2p\cdot J^{-t}z-p\cdot\tau'p\\
     &{}+m\cdot m'-p\cdot p'.
 \end{split}
\]
It follows that $s^2$ is periodic for the full lattice
$\mathbb Z^g+\tau\mathbb Z^g$.  It is bounded on a compact
fundamental parallelepiped, hence everywhere.  Liouville's theorem
on complex lines makes it constant.  Since $s$ never vanishes,
it too is constant (its image lies in the two square roots of a
nonzero constant).  Thus every $L(m,m')$ is even.
Taking $m=0,m'=e_i$, and then $m=e_i,m'=0$, and using that
$B,C$ are even, gives the two asserted parity conditions.

NFL2 now applies with these unchanged half-characteristics.
Squaring (NFL.10) gives
\[
 [\det J\,e^{2\pi iQ_\gamma(z)}\rho(z)^2-\lambda^2]
       \theta[\alpha,\beta](\tau,z)^2=0.
\]
The theta function is not identically zero by NFR5.  The bracket
therefore vanishes on a nonempty open set and hence everywhere
by the identity theorem.  Evaluating the bracket, not a quotient
of theta functions, at zero gives $\lambda^2=\det J$.
\end{proof}
% END NFL3

\subsection{Level structure and comparison of the actual sections}

% BEGIN NFL4
\paragraph{NFL4: the adapted lattice forces the required level-eight sign.}
\label{claim:NFL4}
Let $a:\Lambda\simeq_{\mathcal O_E}\Lambda$ preserve $E_\Lambda$
and satisfy the original two level conditions: it is the identity
modulo $l.M$ on $\overline{l.M_b}\Lambda$, and the identity modulo
$l.A$ on all of $\Lambda$.  In the adapted symplectic basis of
NFR3, write its integral matrix as
$N=(N_{ij})_{1\le i,j\le2}$.  Then
\[
 N_{12}\equiv N_{21}\equiv0\pmod8,\qquad
 \det N_{11}\equiv1\pmod8.
 \tag{NFL.12}
\]
These conclusions also hold for $a^{-1}$.

\begin{proof}
PUN6 upgrades the original half-level congruence to
$(a-I)\Lambda\subset l.M\Lambda$; it uses the complementary
host ideals and the perfect trace pairing, so does not assume
that the lattice is free over $\mathcal O_E$.
Extend $a$ to $V=\Lambda\otimes_{\mathcal O_E}E$.
Its $\mathcal O_E$-linearity and preservation of (NFL.1) imply
preservation of $\operatorname{herm}_G$: replace one variable
by all $s\in E$ and use nondegeneracy of the trace pairing.
Thus $\delta=\det_Ea$ has absolute value one at every complex
embedding.  The endomorphism $(a-I)/l.M$ stabilizes the finite
projective lattice, so its characteristic coefficients are
integral.  Expanding $\det(I+l.M\,T)$ shows
$(\delta-1)/l.M\in\mathcal O_E$.  Since $l.M\ge3$, every
conjugate of this algebraic integer has absolute value at most
$2/l.M<1$.  If it were nonzero, its nonzero integer norm would
have absolute value less than one.  Hence $\delta=1$.

Let $e_0=w_0^4\in\mathcal O_E$ be the host element used in NFR3.
It satisfies $e_0^2\equiv e_0\pmod8$ and acts in that basis by
$\operatorname{diag}(I_g,0)$ modulo eight.
The map $a$ commutes with $e_0$, so its off-diagonal blocks are
zero modulo eight.  Consider $T=e_0a+(1-e_0)I$ on $V$.
Modulo $8\mathcal O_E$, the idempotent $e_0$ splits the base ring
into its two factors.  On those factors $T$ is respectively
$a$ and $I$, whence $\det_ET\equiv1\pmod8$.
This determinant calculation is valid on the finite projective
lattice: localize where it is free, compute the determinant
polynomial there, and glue its integral coefficients.
Taking $E/\mathbb Q$ norms gives
$\det_{\mathbb Q}T\equiv1\pmod8$.
Its integral matrix on $\Lambda/8\Lambda$ is
$\operatorname{diag}(N_{11},I_g)$, so its determinant is
$\det N_{11}$ modulo eight.  This proves (NFL.12).
The inverse preserves the same levels and form, and the same
argument applies to it.
\end{proof}
% END NFL4

% BEGIN NFL5
\paragraph{NFL5: the precise change of period normalization.}
\label{claim:NFL5}
Suppose $P_{b'}(a v)=H_0P_b(v)$ for the lattice isometry in NFL4,
where $H_0\in\operatorname{GL}_g(\mathbb C)$.  Let $N$ be the
matrix of $a^{-1}$ in the adapted basis, and set
\[
 \gamma=\begin{pmatrix}N_{22}^t&N_{12}^t\\N_{21}^t&N_{11}^t\end{pmatrix},
 \quad J=N_{21}^t\tau_b+N_{11}^t.
\]
Then $\gamma$ is integral symplectic, its off-diagonal blocks are
divisible by eight, and $\chi_4(\det\gamma_{22})=1$.  Moreover
\[
 \tau_{b'}=\gamma\tau_b,\qquad
 h_{b'}H_0=J^{-t}h_b,\qquad
 \det h_b=\det J\det h_{b'}\det H_0.
 \tag{NFL.13}
\]

\begin{proof}
Invertibility of $H_0$ follows already from the period relation
if it is initially specified only as a square matrix: $a$ is
surjective and $P_{b'}(\Lambda)$ spans the underlying real
vector space, so $H_0$ is surjective and hence invertible.
Preservation of the integral alternating form gives
$N^t\left(\begin{smallmatrix}0&I\\-I&0\end{smallmatrix}\right)N
=\left(\begin{smallmatrix}0&I\\-I&0\end{smallmatrix}\right)$.
Transposing its four block equations proves symplecticity of
the displayed $\gamma$.  NFL4 gives its congruences and sign.
Normalize $P_b$ to the matrix $[I\ \tau_b]$ in the fixed basis.
The periods of $P_ba^{-1}$ are then
$[A_0\ B_0]$, where
$A_0=N_{11}+\tau_bN_{21}$ and
$B_0=N_{12}+\tau_bN_{22}$.
Its normalizer is $A_0^{-1}h_b$.  Since
$P_ba^{-1}=H_0^{-1}P_{b'}$, this is $h_{b'}H_0$.
Now $A_0^t=J$ and the normalized second block is
$A_0^{-1}B_0=B_0^t(A_0^t)^{-1}$, by symmetry of the normalized
Siegel matrix.  These are precisely (NFL.13).
\end{proof}
% END NFL5

% BEGIN NFL6
\paragraph{NFL6: the determinant identity for rigidified theta sections.}
\label{claim:NFL6}
In NFL5 the two actual base points are equal, and their
exponentials satisfy $\exp_{b'}(H_0z)=\exp_b(z)$.
On every chart of the actual universal line, put
$\Theta_b(z)=\theta[\alpha,\beta](\tau_b,h_bz)fr_j(b,z)$.
Then
\[
 \det h_b\,\Theta_b(z)^2
   =\det h_{b'}\det H_0\,\Theta_{b'}(H_0z)^2.
 \tag{NFL.14}
\]
The assertion compares values in the same rigidified line, not
sections of two merely isomorphic abstract theta bundles.

\begin{proof}
The lattice isometry, its two level conditions, and the equality
of periods meet exactly the \texttt{Family\allowbreak Unif} identification
clause reconstructed in HUF21.  It gives equality of the
moduli points and the displayed exponential identity.
On a chart containing $\exp_bz$ form
$\rho_0(z)=fr_j(b',H_0z)/fr_j(b,z)$.
The actual cocycle equations for the two frames make these
quotients agree on overlaps.  Hence they glue to an entire
nowhere-zero function on $\mathbb C^g$.  The rigidification at
the zero section gives $\rho_0(0)=1$.
The two actual algebraic theta sections are nonzero in the same
one-dimensional section space.  This follows from the retained
complex-fiber section law, or directly from NFR5's Fourier
classification for this fiber.  Consequently for one constant
$\lambda\ne0$,
\[
 \theta[\alpha,\beta](\tau_{b'},h_{b'}H_0z)\rho_0(z)
       =\lambda\theta[\alpha,\beta](\tau_b,h_bz).
\]
Replace $z$ by $h_b^{-1}w$.  The resulting unit $\rho(w)$ has
value one at zero, and NFL5 puts this identity in the form of
NFL3.  Thus $\lambda^2=\det J$.
Since $\Theta_{b'}(H_0z)=\lambda\Theta_b(z)$, the determinant
relation in NFL5 proves (NFL.14).  No division by a theta value
at zero, or anywhere on its divisor, has been used.
\end{proof}
% END NFL6

% BEGIN NFL7
\paragraph{NFL7: change of an actual theta point, with its literal section map.}
\label{claim:NFL7}
Let $y,y'$ be the original \texttt{Is\allowbreak Theta\allowbreak Pt} at $b,b'$ with
nonzero scalars $\kappa,\kappa'$, and suppose the periods are
related as in NFL5.  Then their base points agree, their eigenframes
satisfy
\[
 u'_j=d\exp_b(H_0^{-1}e_j),
\]
and there is a unique $t_0\in\mathbb C^*$ whose actual section
base-change relation is $s'=t_0s$.  It obeys
\[
 t_0^2\kappa^2\det h_{b'}\det H_0
       =\kappa'^2\det h_b.
 \tag{NFL.15}
\]
Here $s'=t_0s$ means the literal \texttt{Sec\allowbreak BC} over the identity:
the canonical map of the two fiber products respects both
projections and pulls the scaled section to the indicated section
on every original line chart.

\begin{proof}
Differentiate $\exp_{b'}H_0=\exp_b$ using the original
dual-number interpretation from HUF1 and TDR1.  This gives the
eigenframe formula, including its prescribed CM type.
Both sections generate the same rank-one section module, so
$s'=t_0s$ for a unique unit.  The canonical fiber-product map
over the equal base point supplies the two projection equations
and chart pullback equations in \texttt{Sec\allowbreak BC}; uniqueness of
the generator makes the scalar independent of the chart.
Use $s=\kappa\Theta_b$ and $s'=\kappa'\Theta_{b'}$, square their
relation and substitute (NFL.14).  Choose any $z$ where theta
is nonzero and a chart containing its exponential image.
The frame value is nonzero there, so cancellation gives
(NFL.15).  Such a point exists by NFR5.
\end{proof}
% END NFL7

\subsection{The two exact transformation laws, for every allowed weight}

% BEGIN NFL8
\paragraph{NFL8: the period matrix of an upper triangular frame move.}
\label{claim:NFL8}
Let $x\in\mathcal D$, let $b\in\operatorname{Bor}(\Psi,r)$, and
put $x'=xb$.  There is an invertible matrix $H_0$ with
$P_{x'}=H_0P_x$.  In the original coordinates indexed by
$j\leftrightarrow(\rho,i)$, both $H_0$ and $H_0^{-1}$ have support
only when the two $\rho$'s agree and the row index $i$ is at most
the column index $i'$.  Write $d_j=(H_0^{-1})_{jj}$.
If $t_0^2\det H_0=1$, then, for every original $\operatorname{Dom}(a,n)$,
\[
 t_0^n\prod_jd_j^{a((ej)_1,(ej)_2)}
       =\prod_{\tau\in\Psi}\prod_{s=0}^{r-1}
                       b_{ss}(\tau)^{\lambda_{\tau,s}},
 \tag{NFL.16}
\]
where $n=2m$ and, using the NFR1 relabelling,
\[
 \lambda_{\tau,s}=
 \begin{cases}
 -a(\operatorname{relab}^{-1}\bar\tau,q-1-s)-m,&s<q,\\
 a(\operatorname{relab}^{-1}\tau,s-q)+m,&s\ge q.
 \end{cases}
 \tag{NFL.17}
\]
These are exactly the source's $\operatorname{laU}$ and $bchar$.

\begin{proof}
Use the literal period formulas of NFR1.  For $s<q$, a period
coordinate is $\operatorname{herm}_{G_\tau}(\operatorname{col}_sX_\tau,
\tau(\phi v))$.  Replacing $X$ by $Xb$ multiplies this early
block by $b_{++}^t$.  For $s\ge q$ it is
$(X_\tau^{-1}\tau(\phi v))_s$, so the late block is multiplied
by $(b_{--})^{-1}$.  The other possible block vanishes because
$b$ is upper triangular.  Thus, in unreversed column order,
\[
 H_0=\bigoplus_{\tau\in\Psi}
            \operatorname{diag}(b_{++}^t,b_{--}^{-1}).
 \tag{NFL.18}
\]
The early columns are reversed by the NFR1 indexing, converting
the first transpose into the indicated upper triangular support.
Inverting preserves that support.  Their diagonals are
$b_{ss}^{-1}$ for $s<q$ and $b_{ss}$ for $s\ge q$.
Also $\det H_0=\prod_{\tau,s<q}b_{ss}
\prod_{\tau,s\ge q}b_{ss}^{-1}$.
Since $n=2m$, $t_0^n=(\det H_0)^{-m}$.  Inserting these three
identities gives (NFL.16) one diagonal entry at a time.
The antitone condition in \texttt{Dom} is retained, although this
particular identity only uses that $n$ is even.  No bound on
$a,n$, and no one-dimensional restriction on the eigenblocks,
has been made.
\end{proof}
% END NFL8

% BEGIN NFL9
\paragraph{NFL9: the complete normalized Borel law.}
\label{claim:NFL9}
For every original $k,x,b$ with $x\in\mathcal D$,
$b\in\operatorname{Bor}(\Psi,r)$ and $xb\in\mathcal D$, there is
\emph{one} $t:\operatorname{Ix}(D_w)\to\mathbb C^*$ such that
\[
 \operatorname{FrBor}(t,u_k(x),u_k(xb))
\]
and, simultaneously for every $a,n$ with $\operatorname{Dom}(a,n)$,
\[
 \operatorname{wchar}(\operatorname{ex}(D_w,a,n))(t)
      =\operatorname{bchar}(\operatorname{laU}(a,n,k))(b).
 \tag{NFL.19}
\]
This is the full original \texttt{norm\allowbreak Un\_\allowbreak bor} statement.

\begin{proof}
Apply NFL7 to the identity lattice isometry and the matrix $H_0$
of NFL8.  The two points have equal base point, and the new
eigenvector $j$ is the derivative of $H_0^{-1}e_j$.
Its coefficient at old coordinate $j'$ is
$(H_0^{-1})_{j'j}$, zero unless its original embedding is the
same and its original index is at most that of $j$.
Set $t(\operatorname{some}(\rho,i))=(H_0^{-1})_{jj}$ and
$t(\operatorname{none})=t_0$, the actual section scalar from NFL7.
The diagonal entries are units.  Thus the derivative coefficients,
the equal base point and the literal section pullback of NFL7
give every clause of \texttt{Fr\allowbreak Bor}.
Substitute the two identities $\kappa_x^2=C\det h_x$ and
$\kappa_{xb}^2=C\det h_{xb}$ into (NFL.15).
All three determinants and $C$ are nonzero, so it reduces to
$t_0^2\det H_0=1$.  NFL8 now gives (NFL.19) for every allowed
weight using the same $t$.  In particular its choice is made
before $a,n$ and before any section to be evaluated.
\end{proof}
% END NFL9

% BEGIN NFL10
\paragraph{NFL10: the complete arithmetic invariance law.}
\label{claim:NFL10}
For every $\gamma\in(M.X_k).\Gamma$ and every $x\in\mathcal D$,
with the original domain guard on $\gamma x$, there is one $t$
such that
\[
 \operatorname{FrBor}(t,u_k(x),u_k(\gamma x)),\qquad
 \operatorname{wchar}(\operatorname{ex}(D_w,a,n))(t)=1
 \tag{NFL.20}
\]
for all $\operatorname{Dom}(a,n)$.  This is the full original
\texttt{norm\allowbreak Un\_\allowbreak gamma} statement, for the actual arithmetic group.

\begin{proof}
The exact \texttt{Model\allowbreak Base.\allowbreak gam} condition says that $\gamma$
is unitary, preserves the actual $\mathcal O_E$-lattice, and
satisfies the half-level $l.M$ and whole-lattice $l.A$ congruences.
It therefore induces
$a=\phi^{-1}\gamma\phi:\Lambda\simeq_{\mathcal O_E}\Lambda$
with precisely the assumptions of NFL4.  The period formulas
give
\[
 P_{\gamma x}(av)=P_x(v).
 \tag{NFL.21}
\]
For early columns this is the unitary identity for
$\operatorname{herm}_G(\gamma\operatorname{col}_sX,\gamma v)$;
for late columns it is $(\gamma X)^{-1}\gamma v=X^{-1}v$.
Thus NFL7 applies with $H_0=I$.  The base points and eigenframes
are literally the same.  Normalization in (NFL.15) gives
$t_0^2=1$.  Take every eigenframe entry of $t$ equal to one and
its section entry equal to $t_0$.  This is the required
\texttt{Fr\allowbreak Bor}; its character is $t_0^n=1$ because $n$ is even.
The argument keeps the full original level conditions; it does
not replace $\Gamma$ by a convenient subgroup.
\end{proof}
% END NFL10

\subsection{Algebraic evaluations and holomorphic readings}

The holomorphy claims NFL12--NFL16 use only their supplied
$F,M,T$ and the hypotheses stated there, together with the two
retained section laws where indicated.  They do not require
\texttt{Good\allowbreak Emb}, \texttt{Model\allowbreak Base} or \texttt{Neg\allowbreak Xi}; those
conditions enter the construction and transformation laws.

% BEGIN NFL11
\paragraph{NFL11: actual evaluation is well defined and has the required weight law.}
\label{claim:NFL11}
At the retained \texttt{Sec\allowbreak Base\allowbreak Change\allowbreak St} premise, the original
\texttt{Eval\allowbreak WD} holds: every natural rule on \texttt{TRing}s
has the same value through any two test-ring presentations of
one framed point, over \emph{any} coefficient algebra $S$.
Moreover, for every reduced $c.F$-algebra $S$, every $a,n$,
every $f\in\operatorname{sec}(X_o,c,a,n)$ and every actual
$\operatorname{FrBor}(t,y,y')$ over $S$,
\[
 \operatorname{evR}(f,y')
   =t(\operatorname{none})^n
       \prod_{\rho,i}t(\operatorname{some}(\rho,i))^{a(\rho,i)}
          \operatorname{evR}(f,y).
 \tag{NFL.22}
\]
No reducedness is required in the first assertion, and no
weight-domination or parity condition is required in the second.

\begin{proof}
TDR5--TDR6 give, on each affine chart of the actual moduli scheme,
the universal frame ring
\[
 A=R[\{g_{\rho,ij}\},s,
           (s\prod_\rho\det g_\rho)^{-1}]
 \tag{NFL.23}
\]
and its actual universal framed point.  This is a reduced
finite-type algebra and hence a literal \texttt{TRing} by TDR12.
Its classifier is unique over every target ring, since it records
the base point, each eigenframe entry and the section generator.
Fix a maximal ideal $\mathfrak m\subset S$ and a framed point $w$.
Choose a chart through its underlying point over $S_\mathfrak m$;
all frames and the section generator give a unique map
$A\to S_\mathfrak m$ classifying $w_\mathfrak m$.

If $w$ was presented by a test point over $p.R$, the preimage
of this chart contains the point induced by $\mathfrak m$.
A principal neighborhood $p.R_g$ suffices.  After further such
localization, the finitely many entries of the eigenframes and
the section generator give its classifier $A\to p.R_g$.
This is exactly the finite localization capture of TDR8 applied
to the universal frame chart; $p.R_g$ remains reduced and of
finite type.  Naturality of the rule identifies its value from
$p$ with the image of its value at the universal point of $A$.
For a second test presentation the classifier into $S_\mathfrak m$
is the same by uniqueness.  Thus the two values agree in every
$S_\mathfrak m$ and consequently in $S$: a nonzero difference
has an annihilator contained in some maximal ideal and remains
nonzero in that localization.  This proves the full
\texttt{Eval\allowbreak WD}, even for nonreduced $S$.  Notice that the
argument used local universal charts, not a potentially
nonreduced common quotient as a new test ring.

Over (NFL.23) adjoin the entries of lower triangular change
matrices in the original eigenframe indices, their diagonal
inverses, and a section-scaling unit and its inverse.  The resulting
ring $B$ is again reduced and of finite type.  It carries two
actual framed points and their universal \texttt{Fr\allowbreak Bor} relation.
The defining weight axiom of $f$ gives (NFL.22) in $B$.
Any specified triangular change over $S_\mathfrak m$ specializes
this universal change, including its diagonal and section units.
The preceding evaluation argument identifies its two evaluations
with those of $y,y'$.  Hence (NFL.22) holds locally at every
maximal ideal and therefore in $S$, whenever both points have
test presentations.  TDR13 supplies those presentations for
every reduced $S$.  All pullbacks here are the literal
\texttt{Fr\allowbreak Map} and \texttt{Sec\allowbreak BC} of TDR3; their composition
and cancellation follow from the two fiber-product projections
and the chart section equalities.
\end{proof}
% END NFL11

% BEGIN NFL12
\paragraph{NFL12: local algebraic coordinates for the actual period family.}
\label{claim:NFL12}
At each $b_0$ of the open period-coordinate domain there are a
neighborhood $B$, a literal test ring $R$, an actual framed point
$y_0$ over $R$, and algebra maps $\psi_b:R\to\mathbb C$ for
$b\in B$ such that the base change $z_b$ of $y_0$ lies above
the actual point $F.u(P_b)$, and every $b\mapsto\psi_b(r)$ is
holomorphic.  The ring and frame are fixed on $B$ before $b$
and before any natural rule are chosen.

\begin{proof}
HUF21 says $b\mapsto F.u(P_b)$ is holomorphic into the actual
moduli scheme over the fixed coefficient embedding $\sigma$.
Choose one of the affine moduli charts with eigenbasis and
theta generator furnished by TDR5 through its value at $b_0$.
The preimage is open; shrink to a neighborhood $B$ lying in it.
The morphism into this affine chart gives the unique evaluation
homomorphism $\psi_b:R\to\mathbb C$.  Every $r\in R$ is a
regular function on the chart and hence its composition with
the holomorphic period family is holomorphic.
Base change of its fixed framed point is legitimate for every
$\psi_b$ by TDR3, using the retained section base-change law.
It gives $z_b$ with the original point, dual-number eigenframe,
and section-chart values, not just an analytic trivialization.
The finite-type reduced ring is a \texttt{TRing} by TDR12.
\end{proof}
% END NFL12

% BEGIN NFL13
\paragraph{NFL13: the change of eigenframes is holomorphic.}
\label{claim:NFL13}
More generally let $R$ be any commutative coefficient algebra,
$y_0$ a framed point, and $\psi_b:R\to\mathbb C$ such that all
its scalar evaluations are holomorphic on an open $B$.
Let $z_b$ be its actual base change above $F.u(P_b)$, and let
$y_b$ be a theta point whose eigenframe is the derivative of
$\exp_b$ at the standard basis.  The unique eigenframe change
matrices in $\operatorname{FrChg}(G_b,a_b,z_b,y_b)$ are
holomorphic functions of $b$.  No assumption on the section
factor $a_b$ is needed for this conclusion.

\begin{proof}
Only finitely many coefficients from $R$ are involved near a
fixed $b_0$: the algebraic group is of finite presentation over
$R$, and the chosen regular functions and differentials have
finite polynomial expressions on local charts.  Choose the
finitely many chart denominators whose images at $b_0$ are
nonzero and shrink $B$ until their images remain nonzero.
Evaluation of each resulting fraction is holomorphic because
its numerator and denominator are evaluations of elements of
$R$.  No complex analytic space structure on
$\operatorname{Spec}R$ is asserted or needed.
Near any chosen base point, smoothness of the actual abelian
scheme gives $g$ regular functions whose relative differentials
are a basis at its zero section.  One may choose them by
selecting a nonzero minor in a finite presentation of the
relative cotangent module.  Shrink the neighborhood so that
this minor stays invertible.
Compose these functions with the HUF holomorphic exponential
$(b,z)\mapsto\exp_bz$.  Their derivatives in $z$ at zero form
a holomorphic invertible Jacobian matrix.  Holomorphy of the
derivatives follows from the Cauchy integral formula on a
fixed small polydisc after shrinking $B$.
Evaluating the same differentials on the algebraic frame of
$y_0$ gives elements of $R$; their images under $\psi_b$ are
holomorphic by hypothesis.  Multiplication by the inverse
Jacobian therefore expresses $z_b$'s tangent vectors in
exponential coordinates by a holomorphic invertible matrix.
Its entries outside each CM eigenspace are zero, since both
sets of vectors have the prescribed $\mathcal O_E$ action.
Inverting each block gives $G_b$ (or its transpose, according
to whether the frame vectors are listed as rows or columns).
In the source convention, the rows of $G_b$ are precisely the
coefficients of $y_b.u_i$ in the frame $z_b.u_j$.
Every entry is holomorphic, since its determinant is nowhere
zero.  This also proves that these matrices are the literal
dual-number combination matrices; TDR1 identifies their
coefficients with the differentials just evaluated.
\end{proof}
% END NFL13

% BEGIN NFL14
\paragraph{NFL14: the actual theta section scaling factor is holomorphic.}
\label{claim:NFL14}
In NFL13 suppose $y_b$ is an actual theta point with scalar
$\kappa(b)$ holomorphic on $B$.  Then the section factor $a_b$
in $\operatorname{FrChg}(G_b,a_b,z_b,y_b)$ is holomorphic.
This holds for the same arbitrary $R$ and all holomorphic
evaluation homomorphisms, with the original section charts.

\begin{proof}
Fix $b_0$.  The pulled-back section $s_{0,b_0}$ generates the
nonzero one-dimensional complex section space.  Hence it does
not vanish everywhere on the complex fiber: this follows
either from its actual local regular functions on the reduced
complex variety or from the complex-fiber section theorem.
Choose a point where its value is nonzero and an original
line chart containing that point.  Surjectivity of the actual
exponential gives a $w_0\in\mathbb C^g$ representing it.
For $b$ near $b_0$, $\exp_bw_0$ remains in the same chart.
Its lift to the family pulled back from $y_0$ is holomorphic:
in finite local charts its coordinates are the evaluations
$\psi_b$ and the HUF holomorphic exponential coordinates,
and their defining base equations agree.  More explicitly,
the fixed algebraic section on this chart is a finite
polynomial or rational expression in those coordinates and
finitely many coefficients from $R$.  Localize at its finitely
many denominators and shrink $B$ so that their evaluations
stay nonzero, as in NFL13.  This proves holomorphy directly
without analytifying the arbitrary ring $R$.
Thus the value $v(b)$ of the fixed algebraic section
$s_0$ at this point is holomorphic, and $v(b_0)\ne0$.
Shrink the neighborhood until $v$ is nowhere zero.

The theta section at that point has the holomorphic value
\[
 \kappa(b)\theta[\alpha,\beta](\tau_b,h_bw_0)fr_j(b,w_0).
\]
Normal convergence from NFR5 gives holomorphy of the theta
factor; the other factors are the actual holomorphic data.
The chart equality in \texttt{Sec\allowbreak BC} says this value is
$a_bv(b)$.  Dividing by $v(b)$ proves local holomorphy of $a_b$.
Since $b_0$ was arbitrary, it proves the assertion on $B$.
The chosen evaluation point is allowed to depend on $b_0$;
there is no assumption that theta is nonzero at the zero section.
\end{proof}
% END NFL14

% BEGIN NFL15
\paragraph{NFL15: natural rules read holomorphically on algebraic frame charts.}
\label{claim:NFL15}
More generally, let $V$ be any complex normed vector space,
$B\subset V$ any set, $R$ any literal test ring, $y_0$ an
actual framed point over $R$, and $\psi_b:R\to\mathbb C$
evaluation homomorphisms with actual base changes $z_b$.
Suppose every scalar evaluation is differentiable on $B$ in
the source's within-set sense, and suppose an actual
framed point $y_b$ is related to $z_b$ by
$\operatorname{FrChg}(G_b,a_b,z_b,y_b)$, with every entry of
$G_b$ and $a_b$ differentiable on $B$.  Then for every natural
rule $f$ on the original test rings,
$b\mapsto\operatorname{evR}(f,y_b)$ is differentiable on $B$.
In the open finite-dimensional setting of NFL12 this is
holomorphy.  No weight condition is necessary here.

\begin{proof}
Use the universal full frame ring (NFL.23) over this $R$.
Its base-changed frame with matrices $G_b$ and section factor
$a_b$ is classified by one map $\eta_b:A\to\mathbb C$.
The entries of $G_b$ and $a_b$ are differentiable on $B$ by hypothesis;
their determinants and $a_b$ are nonzero because both sides
are actual bases and section generators.  Their inverses
are therefore differentiable on $B$.  Every element of $A$ is a
polynomial in the entries, the base elements of $R$, and these
inverses.  Consequently every $b\mapsto\eta_b(v)$ is differentiable on $B$.
Let $y_A$ be its universal framed point.  NFL11 and naturality
give the literal evaluation identity
\[
 \operatorname{evR}(f,y_b)=\eta_b\bigl(f_A(y_A)\bigr).
\]
The element on the right inside parentheses is fixed before
$b$ varies.  This proves the assertion.  It is the universal
frame representing property, including the section factor,
that makes the result hold for every rule simultaneously.
\end{proof}
% END NFL15

% BEGIN NFL16
\paragraph{NFL16: all readings of the normalized points are holomorphic.}
\label{claim:NFL16}
For every $a,n$ with $\operatorname{Dom}(a,n)$, every
$f\in\operatorname{sec}(X_o,c,a,n)$ and every component $k$,
the original reading $\operatorname{rdU}(f,k)$ is holomorphic
on $(M.X_k).\mathrm{dom}$.  This is the complete original
\texttt{differentiable\allowbreak On\_\allowbreak rd\allowbreak U} conclusion for any family satisfying
\texttt{Is\allowbreak Norm\allowbreak Un}, with its unrestricted pointwise square roots.

\begin{proof}
At each period coordinate $b$, take the actual theta point
$y_b^1$ with scalar $1$, whose existence is NFR6.
On a neighborhood from NFL12 compare it with the base change
$z_b$ of the fixed algebraic framed point.
The eigenframe change is holomorphic by NFL13 and the section
factor by NFL14, since the scalar $1$ is holomorphic.
NFL15 therefore proves that
$b\mapsto\operatorname{evR}(f,y_b^1)$ is holomorphic.

The normalized point $u_k(T.\mathrm{co}(b))$ has the same
base point and eigenframes as $y_b^1$; its actual section is
$\kappa_b$ times the section of $y_b^1$.  This is the
\texttt{Fr\allowbreak Bor} with eigenframe diagonal entries $1$ and section
entry $\kappa_b$.  Apply NFL11 over $\mathbb C$ and write $n=2m$:
\[
 \begin{split}
 \operatorname{evR}(f,u_k(T.\mathrm{co}(b)))
   &=\kappa_b^n\operatorname{evR}(f,y_b^1)\\
   &=(C\det h_b)^m\operatorname{evR}(f,y_b^1).
 \end{split}
 \tag{NFL.24}
\]
The last expression is holomorphic and is independent of the
chosen square root at every point.  The coordinate map
$T.\mathrm{co}$ and its inverse are complex linear; transferring
back to the open frame domain proves the assertion.
The source extends $\operatorname{rdU}$ by zero outside the
domain; this has no bearing on holomorphy within the stated
domain and no boundary regularity is claimed.
\end{proof}
% END NFL16

% BEGIN NFL17
\paragraph{NFL17: one family satisfies every field of the original FramesOK.}
\label{claim:NFL17}
For the original supplied $F,M$ and all their original data,
there is one family
\[
 u:\prod_k(M.X_k).\mathrm{dom}\longrightarrow
       \operatorname{FrPt}(X_o,c,\mathbb C_\sigma)
\]
satisfying the complete original \texttt{Frames\allowbreak OK} record.
Every point is also \texttt{From\allowbreak Test}.
The family is chosen before the components, frames, arithmetic
elements, triangular moves, weights and natural rules that occur
in its universal assertions.

\begin{proof}
Choose exactly the family of NFR13, with one theta datum per
component and the pointwise square roots of NFR6.
NFR12 supplies its first field, the full \texttt{Frame\allowbreak OK} at
every point.  NFL9 supplies the second field, with its single
triangular change witness valid for all allowed weights.
NFL10 supplies its third field, the arithmetic invariance
witness valid for all allowed weights.  NFL16 supplies its
fourth field, holomorphic readings for every allowed section.
These are the four fields in the literal definition, with no
replacement of the chosen family at any join.
TDR13 gives \texttt{From\allowbreak Test} at every point because
$\mathbb C_\sigma$ is reduced.  The only additional law premises
are precisely the section base-change and twist laws stated
at the beginning of this chapter.
\end{proof}
% END NFL17

\subsection{Recovering frames and identifying their fibers}

% BEGIN NFL18
\paragraph{NFL18: every original period map is a literal frame period map.}
\label{claim:NFL18}
For every $P:\Lambda\to\mathbb C^g$ satisfying the original
\texttt{Is\allowbreak Period} condition there is $x\in\mathcal D$ whose
literal NFR1 period map is $P$.  Together with NFR2 this gives
the entire original \texttt{exists\_\allowbreak frame\allowbreak Periods}, including
its previously unused surjectivity clause.

\begin{proof}
The lattice and its prescribed model identify
$\Lambda\otimes\mathbb R$ with
$\bigoplus_{\tau\in\Psi}\mathbb C^r$.
The full lattice condition in \texttt{Is\allowbreak Period} makes the real
extension of $P$ an isomorphism.  The $\mathcal O_E$ eigenvalue
condition extends first to $E$ and then to $E\otimes\mathbb R$.
In the unreversed column dictionary, set
\[
 u_{\tau,s}=\begin{cases}
   \overline{P_{\tau,s}},&s<q,\\ P_{\tau,s},&s\ge q.
   \end{cases}
\]
Each block $u_\tau$ is $\tau$-linear in the model coordinate
$v_\tau$, so $u_\tau=M_\tau v_\tau$ for an invertible complex
matrix $M_\tau$.  Put $Y_\tau=M_\tau^{-1}$ and
$S_\tau=Y_\tau^tG_\tau\overline{Y_\tau}$.

Let $K$ be the positive definite Hermitian form provided by
\texttt{Is\allowbreak Period}, so $\operatorname{Im}K(Pv,Pw)=E_\Lambda(v,w)$.
The identity $E_\Lambda(av,w)=E_\Lambda(v,\bar a w)$, first
on the lattice and then over $\mathbb R$, implies
$K(az,w)=K(z,\bar a w)$: their imaginary parts agree, and
replacing $z$ by $iz$ identifies their real parts.
Thus distinct eigenvalue types are orthogonal for $K$.
In particular its early and late blocks at each $\tau$ are
separate, since their types are $\bar\tau$ and $\tau$.
Write these positive definite blocks as $K_+,K_-$, and put
$a_\tau=\operatorname{Im}\tau(-\xi_0)>0$.
The trace identity (NFL.1), on vectors in this real factor, is
$E_\Lambda=-2a_\tau\operatorname{Im}\operatorname{herm}_{G_\tau}$.
In the $u$ coordinates, comparison of imaginary parts gives
\[
 S_\tau=
 \begin{pmatrix}
  \overline{K_+}/(2a_\tau)&0\\
  0&-K_-/(2a_\tau)
 \end{pmatrix}.
 \tag{NFL.25}
\]
The imaginary part determines a first-linear Hermitian form,
since replacing one input by $i$ recovers its real part; hence
this is an equality of matrices, not just of their imaginary
parts on the diagonal.

Let $S_+$ be its early block and define
\[
 B_\tau=\operatorname{diag}((S_+^{-1})^t,I_{r-q}),
 \qquad X_\tau=Y_\tau B_\tau.
 \tag{NFL.26}
\]
For $s<q$, the $s$th row of $B_\tau^tS_\tau$ is the $s$th
standard row, so
\[
 \operatorname{herm}_{G_\tau}(\operatorname{col}_sX_\tau,v_\tau)
      =\overline{u_{\tau,s}}=P_{\tau,s}.
\]
For $s\ge q$, $(X_\tau^{-1}v_\tau)_s=u_{\tau,s}=P_{\tau,s}$.
These are exactly the two original period formulas, including
the reversal of the early indices.  The Gram matrix of $X_\tau$
is $B_\tau^tS_\tau\overline{B_\tau}$, with positive definite
early block and negative definite late block.  Every leading
minor consequently has the prescribed sign: before $q$ it
is positive, and after $q$ it is $(-1)^{s-q}$ times a positive
number.  Hence $X$ lies in the actual domain $\mathcal D$.
Empty blocks use the empty identity matrix and determinant one,
so the argument includes $q=0$ and $q=r$.
\end{proof}
% END NFL18

% BEGIN NFL19
\paragraph{NFL19: the full original triangular covering theorem.}
\label{claim:NFL19}
For every actual complex framed point $y$, there are $k$,
$x\in(M.X_k).\mathrm{dom}$ and the literal period map $P_x$
such that $F.u_k(P_x)=y.x$ and the standard exponential
derivatives equal every prescribed vector $y.u_{\rho,i}$.
Consequently the original \texttt{Cover\allowbreak Tri} holds, with unit
diagonal coefficients and zero off-diagonal coefficients.
This statement concerns every derivative realization in the
source's quantifier order, not just a chosen realization.

\begin{proof}
The onto field of HUF21 supplies a component $k$ and an admissible
period map $P$ above $y.x$.  Its exponential derivative is a
complex-linear isomorphism onto the actual Lie algebra.
Pull back each $y.u_j$ to a vector $w_j\in\mathbb C^g$.
The actual eigenbasis property of $y$ implies that these vectors
form a basis and that their columns have support only in the
corresponding CM eigenspace.  Thus the matrix $W$ with columns
$w_j$ is invertible and block diagonal by original embedding.
Put $V=W^{-1}$ and $P'=VP$.  The image is again a full lattice,
the types are unchanged, and a witnessing positive Hermitian
form is $K'(z,w)=K(V^{-1}z,V^{-1}w)$.  Hence $P'$ satisfies
the exact \texttt{Is\allowbreak Period} condition.

Use the identity lattice isometry and the linear change $V$ in
the HUF family identification.  Its level conditions are
trivial, and it gives
$F.u(P')=F.u(P)=y.x$ and $\exp_{P'}(Vz)=\exp_P(z)$.
Since $Vw_j=e_j$, the standard derivative for $P'$ is exactly
$y.u_j$.  NFL18 now supplies a frame $x$ with $P_x=P'$.
For any other vectors $v_{\rho,i}$ satisfying the source's
specified standard derivative relation, uniqueness of the
dual-number derivative gives $v_{\rho,i}=y.u_{\rho,i}$.
Therefore the coefficient vector $cf_j=\delta_{ij}$ witnesses
every required \texttt{Is\allowbreak Comb\allowbreak At}; its diagonal is $1$ and
the higher coefficients vanish.  This proves the full
\texttt{Cover\allowbreak Tri} quantifiers without using the section of $y$.
\end{proof}
% END NFL19

% BEGIN NFL20
\paragraph{NFL20: the full original triangular fiber theorem.}
\label{claim:NFL20}
Take arbitrary components $k,k'$, frames $x,x'$ in their actual
domains, their literal period maps $P,P'$, and any tangent
families $v,v'$ realizing the standard exponential derivatives.
Suppose the base points are equal and each $v'_{\rho,i}$ is a
linear combination of the $v_{\rho,j}$ with $j\le i$, with
unit coefficient at $i$.  Then there exist
$\gamma\in(M.X_k).\Gamma$ and $b\in\operatorname{Bor}(\Psi,r)$
such that
\[
 (k',x')=(k,\gamma x b)
 \tag{NFL.27}
\]
as an equality in the original dependent union of component
frame spaces.  This is the full original \texttt{Fib\allowbreak Tri}.

\begin{proof}
NFR2 makes $P,P'$ admissible.  The disjointness field of HUF21
first gives $k'=k$; its fiber field then gives an actual
$\mathcal O_E$-lattice isometry $a$ preserving $E_\Lambda$,
the two original level conditions, and a complex-linear
isomorphism $H_0$ with $P'a=H_0P$.
Extend $a$ to the model $E$-space via $\phi$.  This gives
$\gamma\in\operatorname{GL}_r(E)$ on the same full lattice.
As in NFL4, trace nondegeneracy proves that $\gamma$ is unitary.
The two level conditions transfer through $\phi$ unchanged,
so the exact \texttt{Model\allowbreak Base.\allowbreak gam} equivalence gives
$\gamma\in(M.X_k).\Gamma$.  In particular this is a witness
in the original group, not just in its rational unitary group.

By NFL10's period computation, $P_{\gamma x}a=P_x$.
Since $a$ is surjective, $P_{x'}=H_0P_{\gamma x}$.
The family identification gives
$\exp_{x'}(H_0z)=\exp_xz$ on the common fiber.
Let $C$ have as its $i$th column the coefficients expressing
$v'_i$ in the vectors $v_j$; its blocks are upper triangular
in the original indices.  Differentiating the last exponential
identity and using injectivity of the derivative gives
$H_0C=I$.  Hence $H_0=C^{-1}$ has the same block triangular
support, including the separation by original embedding.

Put $b_\tau=(\gamma_\tau X_\tau)^{-1}X'_\tau$.
We explain why it is upper triangular, rather than merely
invertible.  Equality of period maps on the full lattice
extends to every real model vector.  Comparing the early
period coordinates expresses each early column of $X'$ as a
linear combination of the early columns of $\gamma X$:
nondegeneracy of $G_\tau$ permits cancellation against all
vectors.  Thus the lower-left block of $b_\tau$ is zero.
The early-to-early formula is $H_{0,++}=b_{++}^t$ in unreversed
column coordinates.  Because those indices are reversed in
the original eigenframe order, the triangular support of
$H_0$ says $b_{++}$ is upper triangular.  The late formula is
$H_{0,--}=b_{--}^{-1}$; its indices retain their order, so
$b_{--}^{-1}$ and therefore $b_{--}$ are upper triangular.
Both diagonal blocks are invertible.  Together with the
vanishing lower-left block this makes the whole $b_\tau$
upper triangular and invertible.  Its possible upper-right
block is allowed.  Thus $b\in\operatorname{Bor}$ and (NFL.27)
holds.  All steps used only the supplied derivative relations,
so the conclusion is valid for the arbitrary $v,v'$ of the
original statement.
\end{proof}
% END NFL20

\subsection{The complete uniformization in the retained original scope}

% BEGIN NFL21
\paragraph{NFL21: the actual family is a full complex uniformization.}
\label{claim:NFL21}
Let $U=\operatorname{unfOf}(e,F,M,u)$ for any $u$ satisfying
the complete original \texttt{Frames\allowbreak OK}.  At the original
\texttt{Good\allowbreak Emb}, \texttt{Model\allowbreak Base}, \texttt{Neg\allowbreak Xi} and
\texttt{Sec\allowbreak Base\allowbreak Change\allowbreak St} premises, $U$ satisfies the full
\texttt{Is\allowbreak Unf}.  In particular this applies to the single
family of NFL17.  The additional twist law is needed for
that construction, not as an additional premise to this join.

\begin{proof}
We verify the clauses, including the data that accompanies the
orbit assertions.  The rank $r=D_w.p+D_w.q$, the value
$(U.X_k).p_0=D_w.q$, and the exact host CM type are the rank
and type clauses of \texttt{Model\allowbreak Base}.
For an arbitrary framed point $y$, NFL19 supplies $k,x$ whose
base point and standard derivatives are $y$'s.  The
\texttt{Frame\allowbreak OK} field for $u_k(x)$ has the same literal
period map and derivatives.  Its section and $y.s$ are both
generators of the same line, so their ratio is a unit.
Together with the identity eigenframe coefficients this gives
$\operatorname{FrBor}(t,u_k(x),y)$.  This is the covering clause.

If $\operatorname{FrBor}(t,u_k(x),u_{k'}(x'))$, its base equality
and triangular derivative clauses meet NFL20 exactly, giving
$(k',x')=(k,\gamma xb)$ for original $\gamma,b$.
Conversely, for such $\gamma,b$, the arithmetic and Borel
fields of \texttt{Frames\allowbreak OK} give successive actual
\texttt{Fr\allowbreak Bor} relations.  Multiplying their triangular
matrices and their section scalars gives their composite
\texttt{Fr\allowbreak Bor}.  This proves both directions of the exact
dependent-component fiber equivalence.

For $f\in\operatorname{sec}(a,n)$ with $\operatorname{Dom}(a,n)$,
its reading is holomorphic by the fourth \texttt{Frames\allowbreak OK}
field.  The domain is stable under both actions: unitary
left multiplication preserves its Gram matrices, and upper
triangular right multiplication multiplies each signed leading
minor by the positive number $|\det b_{[0,s)}|^2$.
Apply NFL11 to the two \texttt{Fr\allowbreak Bor} relations of the second
and third fields.  Their character equalities give
\[
 \operatorname{rdU}(f,k)(xb)=bchar(\operatorname{laU}(a,n,k))(b)
                              \operatorname{rdU}(f,k)(x),
 \qquad
 \operatorname{rdU}(f,k)(\gamma x)=\operatorname{rdU}(f,k)(x).
 \tag{NFL.28}
\]
The latter is precisely the required arithmetic character law,
because \texttt{Model\allowbreak Base.\allowbreak chi} makes that character identically
one for every weight.  These are all clauses of the original
$\operatorname{IsSec}$ statement.

Finally take the witness $g,e,F$, component map $k\mapsto k$,
and the very same maps $M.\phi_k$ in the last existential field
of \texttt{Is\allowbreak Unf}.  Its trace identity is (NFL.1).  Each
\texttt{Frame\allowbreak OK} supplies the literal period formula,
\texttt{Is\allowbreak Period}, equal base point and all standard derivative
identities.  At every frame with unit absolute diagonal minors,
it also supplies the positive Hermitian matrix, actual
holomorphic nowhere-zero local frames and their cocycle and
rigidification, the unitary lattice multiplier, the actual
section-value identity, the additive lattice fundamental
domain and the Gaussian mean exactly one.  These are exactly
the remaining conjuncts of \texttt{Is\allowbreak Unf}, already proved in
NFR12, with the same section and same period lattice.  Thus no
existential witness has been replaced at the last join.
\end{proof}
% END NFL21

% BEGIN NFL22
\paragraph{NFL22: principal level and every cusp for the prescribed Hermitian space.}
\label{claim:NFL22}
Fix arbitrary original $E,H,D_w,l,X_o,c,\sigma,r,\xi_0,G$ and
$\phi_0:D_w.\mathrm{La}\to_{\mathcal O_E}E^r$ with
$1\le D_w.p+D_w.q$, \texttt{Good\allowbreak Emb}, $\bar\xi_0=-\xi_0$,
$\operatorname{Im}\tau(\xi_0)<0$ on the host type,
$\phi_0$ injective with full $E$-span, $G$ Hermitian and invertible,
and
\[
 D_w.E_F(v,w)=\operatorname{Tr}_{E/\mathbb Q}
       (-\xi_0\operatorname{herm}_G(\phi_0v,\phi_0w)).
\]
At the supplied \texttt{Moduli}, the exact retained
\texttt{Sec\allowbreak Base\allowbreak Change\allowbreak St} and \texttt{Sec\allowbreak Twist\allowbreak St}, and the
inherited HUF/PUN classical-foundation scope, there exists one
$U:\operatorname{Unf}(X_o,c,\sigma,r)$ with all the following:
\[
 \begin{gathered}
 \operatorname{IsUnf}(U),\qquad
 \operatorname{HasPrinc}(U,l.M'),\qquad U.\mathrm{Comp}\ne\varnothing,\\
 \forall k,\lambda,\gamma,\quad (U.X_k).\chi(\lambda,\gamma)=1,\qquad
 \forall k,\quad (U.X_k).G_m=G,\ (U.X_k).\xi_0=\xi_0,\\
 \forall k,x,\quad\operatorname{FromTest}(U.u_k(x)),\\
 \forall\rho,\quad\operatorname{FJ.LatZ}(U.X,\rho),\qquad
 l.A=1\Longrightarrow\operatorname{FJ.ExpoMatch}(U.X,1,\rho,l.M).
 \end{gathered}
 \tag{NFL.29}
\]
Here $l.M'=\operatorname{lcm}(l.M,l.A)$ is the literal original
modulus.  Every cusp $\rho$ and every frame point uses this
same $U$, chosen beforehand.

\begin{proof}
HUF21 supplies the full actual family $F$ for this moduli
problem; its component count is positive by the stipulated
positive rank.  PUN11 supplies $M$ with the exact
\texttt{Model\allowbreak Base} and \texttt{Neg\allowbreak Xi}, with $G$ and $\xi_0$
equal to the prescribed values on every component.  NFL17
supplies one full \texttt{Frames\allowbreak OK} family and NFL21 gives
$\operatorname{IsUnf}$ for its $U$.

To verify the literal \texttt{Has\allowbreak Princ}, take a unitary lattice
stabilizer $\gamma$ congruent to $I$ on the whole lattice
modulo $N=l.M'$.  Since $l.M\mid N$, its difference from
the identity lies in $l.M\Lambda$ for every lattice vector,
in particular on the host half-ideal sublattice.  Since
$l.A\mid N$, the second required congruence also holds.
These implications mean actual vectors: if
$(\gamma-I)v=Nw$ and $N=md$, use the vector $dw$ as the
$m$-divisibility witness.  The exact \texttt{Model\allowbreak Base.\allowbreak gam}
equivalence therefore puts $\gamma$ in $\Gamma$.
This proves \texttt{Has\allowbreak Princ} without imposing any extra
congruence on $U$.  PUN6 in fact also identifies the original
group with the principal subgroup at this modulus.

Nonempty components and trivial characters are the already
specified fields of $F$ and $M$.  The equality of Hermitian
data is the prescribed-space conclusion of PUN11.  TDR13
gives the literal \texttt{From\allowbreak Test} witnesses for all its
complex framed points.  PUN12--PUN14 give the full center
lattice for every cusp of these actual arithmetic groups;
PUN15--PUN19 identify the original integral Fourier exponent
lattice when $l.A=1$, producing the asserted
\texttt{Expo\allowbreak Match} at level $l.M$.  Their statements are
uniform in the cusp and use these same lattices and groups,
so they may be joined after $U$ is chosen.

This reconstructs the mathematical return of the source's
\texttt{exists\_\allowbreak unf\_\allowbreak has\allowbreak Princ} at the two explicitly retained
section laws.  The source obtains those laws from its global
law packages; their global suppliers have not been proved in
this packet.  Accordingly (NFL.29) is a complete theorem
with these stated premises, not an assertion that the
unconditional source wrapper or a selected whole obligation
has been independently reconstructed.
\end{proof}
% END NFL22

\subsection{Source and inherited-proof map}

\begingroup\raggedright
The metadata records every complete pinned file, exact range,
SHA256, and the proof-block hashes of all claims and inherited
inputs.  The following identifies the mathematical source of
each group.  Paths in this paragraph are relative to
\texttt{src/\\allowbreak{}Development/\\allowbreak{}Abelian\allowbreak{}Variety\allowbreak{}Theory/\\allowbreak{}Proofs/\\allowbreak{}}.
NFL1--NFL2 reconstruct the finite Gauss sum and Poisson argument
in \texttt{Homogeneous\allowbreak Ideals/\allowbreak }\\
\texttt{From\allowbreak{}Blocks\allowbreak{}Map\allowbreak{}Int\allowbreak{}Cast\allowbreak{}Mem\allowbreak{}Symplectic\allowbreak{}Group.lean}, together
with its finite-form diagonalization supplier.  NFL3 reconstructs
\texttt{Unitary\allowbreak{}Existence/\\allowbreak{}Two\allowbreak{}Dvd\allowbreak{}Multiplier\allowbreak{}Of\allowbreak{}Theta\allowbreak{}ZAct\allowbreak{}Eq\allowbreak{}Unit\allowbreak{}Mul.lean}
and \texttt{Ratio\allowbreak Sq\allowbreak Eq\allowbreak Det\allowbreak JAut.\allowbreak lean}.
NFL4--NFL7 correspond to \texttt{Cm\allowbreak{}Unit\allowbreak{}Pow\allowbreak{}Four\allowbreak{}Idempotent\allowbreak{}Mod\allowbreak{}Eight.lean},
\texttt{Change\allowbreak Of\allowbreak Normalisation.\allowbreak lean},
\texttt{Exists\allowbreak Sec\allowbreak Generator\allowbreak Fibre.\allowbreak lean},
\texttt{Theta\allowbreak{}Frame\allowbreak{}Det\allowbreak{}Mul\allowbreak{}Sq\allowbreak{}Invariant.lean} and
\texttt{Theta\allowbreak{}Points\allowbreak{}Change\allowbreak{}Level\allowbreak{}Isometry\allowbreak{}Linear\allowbreak{}Map.lean} in that
same unitary-existence directory.

NFL8--NFL10 reconstruct \texttt{Triangular\allowbreak{}Frame\allowbreak{}Moves\allowbreak{}Period\allowbreak{}Block\allowbreak{}Triangular.lean},
\texttt{Normalised\allowbreak{}Framed\allowbreak{}Point\allowbreak{}Triangular\allowbreak{}Frame\allowbreak{}Weight\allowbreak{}Rule.lean},
\texttt{Per\allowbreak Fr\allowbreak Lact.\allowbreak lean} and
\texttt{Normalised\allowbreak{}Framed\allowbreak{}Point\allowbreak{}Arithmetic\allowbreak{}Group\allowbreak{}Invariant\allowbreak{}Sig.lean}.
NFL11 is the evaluation and weight portion of
\texttt{Framed\allowbreak{}Moduli\allowbreak{}Algebraic\allowbreak{}Laws\allowbreak{}Chart\allowbreak{}Weight\allowbreak{}Integrality.lean},
using the actual universal charts already proved in TDR;
it makes no claim here to that file's separate integral-lattice
or unit-independence conjuncts.
NFL12--NFL16 reconstruct, respectively,
\texttt{Exists\allowbreak{}Chart\allowbreak{}Holomorphic\allowbreak{}Coordinates\allowbreak{}Near\allowbreak{}Frame.lean},
\texttt{Differentiable\allowbreak{}On\allowbreak{}Exponential\allowbreak{}Frame\allowbreak{}Matrix.lean},
\texttt{Differentiable\allowbreak{}On\allowbreak{}Theta\allowbreak{}Algebraic\allowbreak{}Generator\allowbreak{}Factor.lean},
\texttt{Differentiable\allowbreak{}On\allowbreak{}Natural\allowbreak{}Rule\allowbreak{}Holomorphic\allowbreak{}Frames.lean} and
\texttt{Differentiable\allowbreak{}On\allowbreak{}Readings\allowbreak{}Normalised\allowbreak{}Theta\allowbreak{}Points.lean}.

NFL18--NFL20 reconstruct the onto clause of
\texttt{Exists\allowbreak Frame\allowbreak Periods.\allowbreak lean} through
\texttt{Exists\allowbreak Frame\allowbreak Of\allowbreak Is\allowbreak Period.\allowbreak lean}, then
\texttt{Exists\allowbreak Inverse\allowbreak Block.\allowbreak lean} and
\texttt{Block\allowbreak Triangular\allowbreak Of\allowbreak Mul\allowbreak Eq\allowbreak One.\allowbreak lean}, with their exact period
and derivative identifications.
NFL17 and NFL21--NFL22 perform the original joins of
\texttt{Exists\allowbreak{}Framed\allowbreak{}Points\allowbreak{}All\allowbreak{}Normalisations\allowbreak{}Family\allowbreak{}Unif.lean},
\texttt{Family\allowbreak{}Uniformisation\allowbreak{}Is\allowbreak{}Complex\allowbreak{}Uniformisation.lean} and
\texttt{Congr\allowbreak Of\allowbreak Dvd\allowbreak Modulus.\allowbreak lean}.  The exact source
\texttt{Norm\allowbreak Const}, \texttt{Univ\allowbreak Frame}, \texttt{Is\allowbreak Comb\allowbreak At},
\texttt{Algebra\allowbreak Map\allowbreak CAt}, \texttt{Un\allowbreak Free} and \texttt{Base\allowbreak Norm}
definitions are included in the manifest, so the source names
in these endpoints refer to their original records and
quantifier order.

The inherited human proofs are NFR1--NFR13, HUF1 and HUF21,
TDR1, TDR3, TDR5--TDR6, TDR8, TDR12--TDR13, PUN6,
PUN11--PUN19.  Each is bound to its exact earlier proof block
and keeps its earlier premises.  The present chapter has no
new Fourier-measure normalization, replacement section, hidden
holomorphic square-root choice, or altered coefficient field.
\par\endgroup


\section{Completing the seven cohomological fields on one linked frame}
\label{sec:linked-frame-completion}

\paragraph{The endpoint.}
This chapter proves all seven fields left open by MFC14:
\texttt{phi}, \texttt{gen}, \texttt{cl}, \texttt{ec},
\texttt{b1p}, \texttt{b1}, and \texttt{h1}.
The proofs apply to the actual linked frame already constructed there.
They neither replace its quotient nor enlarge its class-rank constant.
Consequently the original \texttt{frame\_\allowbreak ex} theorem, with its full
\texttt{Gen} and \(\texttt{GC.\allowbreak Hypotheses}_{2}\) antecedents, is obtained.
For the original coefficient DVR, CAF7 then installs the character
\texttt{Seed\allowbreak Clause} and gives \texttt{Has\allowbreak Frame}.
The separate \texttt{gen\_\allowbreak gd} and \texttt{gen\_\allowbreak host} entry theorems and
the two automorphic seed representations are not asserted here.

\paragraph{Retained mathematical inputs.}
Write \(P:\texttt{GC.\allowbreak Hypotheses}_{2}\) for the original published
cohomological package.  We use its local finiteness, local Tate
duality, local Euler characteristic, global finiteness, global Euler
characteristic and normalized Poitou--Tate fields at their stated
finite-module instances.  MFC's two strict classes retain the
Greenberg--Wiles input used there.  Class-group finiteness, the
Dirichlet unit theorem and finite class-field reciprocity are the
named classical number-theoretic inputs below; the last is the
same class-character dictionary already used in MFC5 and TW1.
We give the needed Kummer, coefficient, quotient, inertia and cone
arguments explicitly.  In particular the specialized source
\texttt{Character\allowbreak H1Galois\allowbreak Quotient} theorem is replaced, for its
trivial-character application here, by a CM-unit argument.

\subsection{The linked data and literal cohomology spaces}

Until the final assembly let \(F\) be a totally real number field,
\(\Omega\) its specified algebraic closure, and
\(G_F=\operatorname{Gal}(\Omega/F)\).  Let \(\ell\) be odd and
\(k\) a field of characteristic \(\ell\).  The given frame
\(\Phi\) has a profinite group \(\Gamma\), a sign character
\(\delta:\Gamma\to\{\pm1\}\), an element \(c\) with
\(\delta(c)=-1\), and \(\Gamma_0=\ker\delta\).
It has finite sets \(V,\mathrm{Inf}\), profinite groups \(G_v\),
maps
\[
 \operatorname{loc}_{0,v}:G_v\longrightarrow\Gamma_0,\qquad
 \operatorname{loc}_v:G_v\longrightarrow\Gamma,
 \qquad I_v\subset G_v,\qquad d_v\in\mathbb N,
\]
the cyclotomic character \(\epsilon:\Gamma\to\mathbb Z_\ell^\times\),
and a residual character \(\Psi:\Gamma\to k^\times\).
Put
\[
 \operatorname{loc}_{2,(v,0)}=\operatorname{loc}_{0,v},\qquad
 \operatorname{loc}_{2,(v,1)}(s)
       =c\,\operatorname{loc}_{0,v}(s)c^{-1},\qquad
 N_v=I_v\cap\ker(\epsilon\operatorname{loc}_v).
\]
The element in the second formula is regarded as an element of
\(\Gamma_0\).  For any characteristic-\(\ell\) field \(E\), denote
the reduction of \(\epsilon\) in \(E^\times\) by \(\omega_E\),
and the sign character in \(E^\times\) by \(\delta_E\).

The full \texttt{FLink} gives a finite set \(S\) consisting
\emph{exactly} of the places of \(F\) above \(\ell\), a normal open
subgroup \(K\subset G_F\) of index prime to \(\ell\), and a
continuous surjection
\[
 \pi:G_F\twoheadrightarrow\Gamma,\qquad
 \ker\pi=
 \overline{\langle I_A\cap K:
                  A\mid v,\ v\notin S\rangle^{\mathrm{normal}}}.
                                                        \tag{LFC.1}
\]
There are bijectively indexed places \(v\mapsto\mathfrak p_v\in S\),
chosen prolongations \(A_v\), and continuous isomorphisms
\(\iota_v:G_v\simeq D_{A_v}\).  Under these isomorphisms,
\(\operatorname{loc}_v=\pi\iota_v\), \(I_v\) is the actual inertia
group, and \(d_v=[F_{\mathfrak p_v}:\mathbb Q_\ell]\).
The set \(\mathrm{Inf}\) indexes all real places, and \(c_w\)
is the image of a complex conjugation at that place.
Furthermore
\[
 \epsilon\pi=\epsilon_F,\qquad
 K\subset\ker(\delta\pi)\cap\ker(\Psi\pi),\qquad
 \ker\delta\ \hbox{is open}.                              \tag{LFC.2}
\]
No assertion of \texttt{Ok} is included in this input.
When only \texttt{GLink} is needed, the hypotheses are (LFC.1),
\(\ell\nmid[G_F:K]\), and \(S\) containing the \(\ell\)-places.

All cochains below are continuous, with discrete coefficient
modules.  For a representation \(\rho\) on \(W\), write
\[
 (d x)(g)=\rho(g)x-x,\qquad
 (d z)(g,h)=z(g)+\rho(g)z(h)-z(gh).
\]
For stable local subspaces \(b_v\subset W\), the source's
\(H^1_C(\operatorname{loc},W,b)\) consists of pairs
\((z,(x_v))\), where \(z\) is a cocycle and
\[
 z(\operatorname{loc}_v s)
       -\bigl(\rho(\operatorname{loc}_v s)x_v-x_v\bigr)\in b_v,
                                                        \tag{LFC.3}
\]
modulo the pairs \((dX,(X+y_v))\), \(y_v\in b_v\).
This is the literal \texttt{Z1C}/\texttt{continuous\allowbreak Coboundaries1}
quotient, including its local primitives.  In particular local
condition \(0\) retains those primitives, whereas local condition
\(W\) makes them irrelevant.

\subsection{Finite characters and the uniform class exponent}

% BEGIN LFC1
\paragraph{LFC1. Kummer classes with bounded ramification are finite.}
\label{claim:LFC1}
Let \(E\) be a number field containing \(\mu_\ell\), and let
\(T\) be a finite set of its finite places.  Put
\[
 V_T(E)=\{a\in E^\times:\ell\mid v(a)\text{ for every }v\notin T\}.
\]
There is an exact sequence of finite-dimensional
\(\mathbb F_\ell\)-spaces
\[
 0\longrightarrow
   \mathcal O_{E,T}^{\times}/(\mathcal O_{E,T}^{\times})^\ell
 \longrightarrow V_T(E)/E^{\times\ell}
 \longrightarrow \operatorname{Cl}(\mathcal O_{E,T})[\ell]
 \longrightarrow0.                                     \tag{LFC.4}
\]
Consequently there are only finitely many continuous additive
characters \(G_E\to\mathbb F_\ell\) unramified outside \(T\).

\begin{proof}
In the ring with the primes \(T\) inverted, write
\((a)=\mathfrak a^\ell\).  Unique ideal factorization makes
\(\mathfrak a\) unique, and sends the class of \(a\) to
\([\mathfrak a]\).  Multiplying \(a\) by an \(\ell\)-th power
changes \(\mathfrak a\) by a principal ideal.  Conversely an ideal
class killed by \(\ell\) has a representative with
\(\mathfrak a^\ell=(a)\), giving surjectivity.  The class is zero
precisely when \(a=u b^\ell\), with \(u\) a \(T\)-unit.
If a \(T\)-unit is an \(\ell\)-th power in \(E^\times\),
its root is again a \(T\)-unit by valuations.  This proves
exactness.  The unit theorem and class-group finiteness prove
the stated finiteness.

Fix a primitive root \(\zeta\in E\).  A nonzero character
\(\chi:G_E\to\mathbb F_\ell\) cuts out a cyclic extension of
degree \(\ell\).  If \(\tau\) generates its Galois group with
\(\chi(\tau)=1\), the operator
\(\sum_{j=0}^{\ell-1}\zeta^{-j}\tau^j\) is nonzero by linear
independence of distinct field embeddings.  Choose \(x\) on
which it is nonzero and call its value \(\beta\).
Then \(\tau\beta=\zeta\beta\), \(\beta^\ell=a\in E^\times\),
and \(g\beta=\zeta^{\chi(g)}\beta\).
The zero character corresponds to \(a=1\).  If two such
\(a\)'s differ by an \(\ell\)-th power, the corresponding
\(\beta\)'s differ by an element of \(E^\times\) times a power
of \(\zeta\), so the two characters are equal.

At a place outside \(T\), the extension cut out by \(\chi\)
is unramified.  Its normalized valuation therefore takes integral
values on \(\beta\), whence \(v(a)=\ell v(\beta)\) is divisible
by \(\ell\).  The characters inject into the finite middle
space of (LFC.4).  This proves the final assertion without
assuming the converse local ramification criterion.
\end{proof}
% END LFC1

% BEGIN LFC2
\paragraph{LFC2. The full \texttt{Phi\allowbreak Cond} on every open subgroup.}
\label{claim:LFC2}
For any \texttt{GLink} over any number field \(F\), without
a total-reality assumption, every open subgroup \(H\subset\Gamma\)
has only finitely many homomorphisms \(H\to(\mathbb F_\ell,+)\)
with open kernel.  Thus \(\operatorname{PhiCond}(\Gamma,\ell)\)
holds, including its quantifier over all open \(H\).

\begin{proof}
Let
\[
 J=\pi^{-1}(H)\cap K\cap\ker(\bar\epsilon_F),\qquad
 E=\Omega^J.
\]
This is an open subgroup, so \(E/F\) is finite, and \(E\)
contains \(\mu_\ell\).  The map \(G_E=J\to H\) induced by
\(\pi\) has image of finite index: representatives of the
finitely many cosets of \(J\) in \(G_F\) also cover its image
cosets in \(\Gamma\), and hence in \(H\).
Choose a nonzero rational integer \(n\) divisible by the residue
characteristic of every place of \(S\).
At a place of \(E\) not dividing \(n\), the underlying place of
\(F\) is outside \(S\).  Its inertia in \(G_E\) is contained
in \(I_A\cap K\), and is killed by (LFC.1).

Consequently restriction sends each character of \(H\) to one
of the finitely many characters supplied by LFC1, with \(T\)
the places of \(E\) dividing \(n\).  Fix finitely many coset
representatives for the image of \(G_E\) in \(H\).  A character
is determined by its restriction and its value on these
representatives, so there are only finitely many possibilities.
Normality of this image is not needed for this argument.
For a homomorphism to a finite discrete group, an open kernel
is equivalent to continuity, since every fiber is a coset of
the kernel.  This is exactly the source's definition of
\texttt{Phi\allowbreak Cond}, not a statement only about \(\Gamma\) itself.
\end{proof}
% END LFC2

% BEGIN LFC3
\paragraph{LFC3. One nonzero exponent for all unramified characters.}
\label{claim:LFC3}
For a full \texttt{FLink}, there is \(h>0\), chosen before
the finite abelian \(\ell\)-group \(B\) and the character
\(\chi:\Gamma_0\to B\), such that every open-kernel \(\chi\)
satisfying
\[
 \chi(\operatorname{loc}_{2,(v,b)}s)=1
 \quad(v\in V,\ b\in\{0,1\},\ s\in I_v)
                                                        \tag{LFC.5}
\]
satisfies \(\chi(g)^h=1\) for every \(g\in\Gamma_0\).
This is the full \texttt{cl} field.

\begin{proof}
Put \(U=\pi^{-1}(\Gamma_0)\), \(E=\Omega^U\), and choose
\[
                  h=2\,|\operatorname{Cl}(\mathcal O_E)|.
                                                        \tag{LFC.6}
\]
This positive integer depends only on the linked data.
The pullback of \(\chi\) to \(G_E=U\) is unramified at every
finite place.  Indeed, above \(v\in S\), conjugate the
prolongation to the chosen \(A_v\).
If the image \(a\) of the conjugating element has
\(\delta(a)=1\), the character sees a conjugate inside
\(\Gamma_0\) of \(\operatorname{loc}_{0,v}(I_v)\).
If \(\delta(a)=-1\), write \(a=(ac^{-1})c\);
it sees a conjugate inside \(\Gamma_0\) of the second branch
\(c\operatorname{loc}_{0,v}(I_v)c^{-1}\).
Both are killed by (LFC.5), since \(B\) is abelian.
This uses the actual inertia groups, including their
intersections with \(G_E\).

At a finite place outside \(S\), let \(t\) be such an inertia
element and put \(r=[G_F:K]\).  Normality of \(K\) gives
\(t^r\in K\), so (LFC.1) gives \(\pi(t)^r=1\).
The order of \(\chi(\pi(t))\) is an \(\ell\)-power and divides
the prime-to-\(\ell\) integer \(r\); hence it is one.
Finally the square of the pullback character kills every
real complex conjugation, since these have order at most two.
Finite class-field reciprocity therefore factors that square
through the ordinary finite ideal class group of \(E\).
It is annihilated by the class number, giving (LFC.6).
The map \(G_E\to\Gamma_0\) is onto, by surjectivity of \(\pi\),
so this proves the assertion for every \(g\).
No bound on \(B\), or choice depending on \(\chi\), entered
the choice of \(h\).
\end{proof}
% END LFC3

\subsection{Local generators with the required cyclotomic exception}

% BEGIN LFC4
\paragraph{LFC4. The elementary procyclic cohomology calculation.}
\label{claim:LFC4}
In this claim \(\ell\) may be any prime and the base number
field need not be totally real.
Let \(D\) be a local decomposition group, \(I\) its inertia,
and \(\varphi\in D\) a lift of arithmetic Frobenius, so that
\(D/I\simeq\widehat{\mathbb Z}\), with \(\varphi\) mapping to
one.  If \(J\triangleleft D\) is closed, contained in \(I\),
and \([I:J]\) is finite and prime to \(\ell\), then for every
finite-dimensional finite-field characteristic-\(\ell\)
continuous representation \(W\) on which \(J\) acts trivially,
\[
 H^2(D/J,W)=0,\qquad
 \dim H^1(D/J,W)=\dim H^0(D/J,W).                         \tag{LFC.7}
\]

\begin{proof}
For a procyclic group with generator \(t\), a continuous
one-cocycle is determined by \(x=z(t)\).
Every \(x\in W\) occurs: if the action has period \(a\),
the cocycle on the integers defined by its partial sums
factors through the cyclic group of order \(a\ell\), because
the sum over \(a\ell\) terms is zero.
Coboundaries correspond to \((t-1)W\).  Thus
\(H^1=W/(t-1)W\), whose dimension is \(\dim\ker(t-1)\).
A continuous two-cocycle factors through some finite cyclic
quotient.  For a cyclic quotient of order \(a\) acting through
\(t\), the periodic resolution gives
\(H^2=W^t/(1+t+\cdots+t^{a-1})W\): this follows by lifting
the cyclic generator in the extension defined by the cocycle,
whose \(a\)-th power is the displayed invariant.
Inflation to the quotient of order \(a\ell\) multiplies that
invariant by \(\ell\), and hence kills the class.
The same conclusion holds in continuous cohomology, since
the new primitive is a cochain on a finite quotient.

Now \(A=I/J\) is finite of order invertible in the coefficient
field.  Averaging over \(A\) projects \(W\) onto \(W^A\).
It also contracts the positive-degree homogeneous cochains
of \(A\): insert an initial variable \(a\), sum over \(a\in A\),
and divide by \(|A|\).  The cochain identity is \(ds+sd=1\)
in positive degrees.  Apply this contraction to the
normal-subgroup double complex for
\(1\to A\to D/J\to\widehat{\mathbb Z}\to1\);
all positive \(A\)-rows disappear and the remaining row is
the continuous cochain complex of
\(\widehat{\mathbb Z}\) with coefficients \(W^A\).
This proves (LFC.7) by the preceding calculation.
All maps are continuous: the \(A\)-sums are finite and every
cochain used here factors through a finite quotient.
\end{proof}
% END LFC4

% BEGIN LFC5
\paragraph{LFC5. Exactly \(d\) universal inertia detectors.}
\label{claim:LFC5}
Let \(v\mid\ell\), \(D=D_A\), \(I=I_A\), and
\(d=[F_v:\mathbb Q_\ell]>0\).  Assume the local fields of
\(P\), and assume \(\bar\epsilon_F|_D\ne1\).
For any normal open \(K\subset G_F\) of prime-to-\(\ell\)
index there are \(u_1,\ldots,u_d\in I\cap K\), with a
distinguished index \(a_0\), such that
\[
 \epsilon_F(u_a)=1\quad(a\ne a_0),                       \tag{LFC.8}
\]
and, for every finite abelian \(\ell\)-group \(B\), every
open-kernel homomorphism \(\chi:D\to B\) which kills all
\(u_a\) kills all of \(I\).

\begin{proof}
The residue-field Galois group is \(\widehat{\mathbb Z}\):
on its finite subfields Frobenius has precisely their degrees
as orders, and taking the inverse limit gives this statement.
Choose a Frobenius lift \(\varphi\in D\).  Then \(I\) and
\(\varphi\) topologically generate \(D\), and for every
integer \(n>0\) there is a continuous degree map
\(\deg_n:D\to\mathbb Z/n\) killing \(I\) and taking
\(\varphi\) to one.

Let \(X=\operatorname{Hom}_{\rm cont}(D,\mathbb F_\ell)\).
Local finiteness makes \(X\) finite-dimensional.
Local Tate duality gives \(H^2(D,\mathbb F_\ell)=0\):
its dual is the invariant subspace of the nontrivial
cyclotomic character.  The local Euler characteristic therefore
gives \(\dim X=d+1\).
Evaluation at \(\varphi\) is onto, witnessed by \(\deg_\ell\).
Thus \(Y=\{f\in X:f(\varphi)=0\}\) has dimension \(d\).
Put \(J=I\cap K\), \(J_0=J\cap\ker\epsilon_F\), and
\[
                   W_0=\{f\in Y:f|_{J_0}=0\}.
 \]
Evaluation on \(J\) separates \(Y\).  Indeed, if \(f|_J=0\)
and \(r=[G_F:K]\), then \(t^r\in J\) for \(t\in I\);
as \(r\) is invertible modulo \(\ell\), \(f|_I=0\).
Together with \(f(\varphi)=0\) and continuity this gives
\(f=0\).

If \(W_0=0\), successively choose points of \(J_0\) which
reduce the common evaluation kernel.  After at most \(d\)
choices that kernel is zero.  Pad the list by identities.
If \(W_0\ne0\), the same procedure chooses at most \(d-1\)
points of \(J_0\) with common kernel \(W_0\).
Let \(N\) be the common kernel in \(J\) of the finitely many
characters in \(W_0\).  It is open in the compact group \(J\)
and contains \(\ker(\epsilon_F|_J)\).
Hence \(J/N\) is a finite continuous quotient of the closed
subgroup \(\epsilon_F(J)\subset\mathbb Z_\ell^\times\).
Such a quotient is cyclic when \(\ell\) is odd.
For completeness, compactness and openness show that its
kernel contains the intersection with \(1+\ell^a\mathbb Z_\ell\)
for some \(a\): otherwise a nested sequence outside the
kernel would have a limit in the kernel with cyclotomic value
one.  The quotient is therefore a quotient of a subgroup of
\((\mathbb Z/\ell^a)^\times\), a cyclic group for odd \(\ell\).
Choose \(u_{a_0}\in J\) mapping to a generator of \(J/N\).
Vanishing at this one point forces a member of \(W_0\) to
vanish on \(J\), and hence to be zero.  Pad with identities
to obtain \(d\) points.  In both cases (LFC.8) holds, and
a character in \(X\) killed by the \(u_a\)'s and \(\varphi\)
is zero.

To obtain the assertion for every \(B\), put \(n=|B|\) and
consider the finite abelian \(\ell\)-group
\[
             C=\operatorname{im}(\chi,\deg_n)
                       \subset B\times\mathbb Z/n.
 \]
It is generated by the image of \(\varphi\).
Otherwise its quotient by that cyclic subgroup is a nonzero
finite abelian \(\ell\)-group and has a nonzero homomorphism
to \(\mathbb F_\ell\) (reduce one cyclic factor modulo \(\ell\)).
Composing with \(D\to C\) contradicts the preceding detection
statement: it kills all \(u_a\) and \(\varphi\).
For \(t\in I\), write its image in \(C\) as the \(b\)-th
power of the image of \(\varphi\).  Its degree is zero, so
\(n\mid b\); since \(n\) annihilates \(B\), \(\chi(t)=1\).
This also handles \(B=1\).  The same \(u_a\)'s work for all
groups \(B\) and all characters \(\chi\).
\end{proof}
% END LFC5

% BEGIN LFC6
\paragraph{LFC6. The exact local \texttt{gen} field.}
\label{claim:LFC6}
Assume \(P\), a full \texttt{FLink}, and
\(\omega_k\operatorname{loc}_v\ne1\) for every \(v\).
Then \(\Phi\) satisfies the full \texttt{gen} field:
for every \(v\) there are \(u:\operatorname{Fin}(d_v)\to G_v\)
and \(a_0\in\operatorname{Fin}(d_v)\), with
\[
 u_a\in I_v,\quad \Psi(\operatorname{loc}_v u_a)=1,\quad
 a\ne a_0\Longrightarrow u_a\in N_v,
                                                        \tag{LFC.9}
\]
which detect the restriction to inertia of every continuous
character to every finite abelian \(\ell\)-group.

\begin{proof}
The inclusion \(\mathbb F_\ell\hookrightarrow k\) is injective.
The compatibility \(\epsilon\pi=\epsilon_F\) thus identifies
the stated nontriviality with
\(\bar\epsilon_F|_{D_{A_v}}\ne1\), including the source's
modulo-\(\ell^1\) convention.
Apply LFC5 at \(A_v\) and pull its points back through the
continuous isomorphism \(\iota_v:G_v\simeq D_{A_v}\).
Every point lies in \(I_{A_v}\cap K\), so its image has
\(\Psi=1\) by (LFC.2).  All but the distinguished point
have cyclotomic value one, giving membership in \(N_v\).
The local degree is exactly \(d_v\), not an upper bound.
A character of \(G_v\) transfers through \(\iota_v^{-1}\)
to a character of \(D_{A_v}\) with open kernel, so LFC5
proves its universal inertia assertion verbatim.
\end{proof}
% END LFC6

\subsection{The actual strict-cohomology dictionary and rank bounds}

% BEGIN LFC7
\paragraph{LFC7. Literal class characters and extension of coefficients.}
\label{claim:LFC7}
Here \(F\) may be any number field and \(\ell\) any prime.
For an open subgroup \(J\subset G_F\), let
\[
 X_E(J)=\{f:J\to E:\ f\text{ locally constant and additive},
                   \ f(I_A\cap J)=0\text{ for every }A\}.
 \]
This is exactly \(\texttt{Kum.\allowbreak XDual}(1,J,E,E)\).
For every characteristic-\(\ell\) field \(E\),
\[
 \dim_E X_E(J)=\dim_{\mathbb F_\ell}X_{\mathbb F_\ell}(J)
                  =\operatorname{clRk}_\ell(J).           \tag{LFC.10}
\]
In particular the literal \texttt{cl\allowbreak Rk}, not a substitute
class number, bounds these coefficient spaces.

\begin{proof}
The extra decomposition condition in \texttt{XDual} is
imposed only when \(A(1)<1\), which never happens.
The other three conditions in its definition are exactly
those displayed above.  Identify \(J=G_{E_0}\) for the finite
fixed field \(E_0\).  Inertia under this identification is
the intersection with \(J\), by its definition as the
kernel of the action on the residue field.
Finite class-field reciprocity identifies the prime-field
space with the characters of the narrow ideal class group.
For odd \(\ell\) these are also the characters of the
ordinary class group, since the extra sign quotient is
a group of order a power of two.
Thus this prime-field space is finite-dimensional.

Any \(f\in X_E(J)\) has finite image, because \(J\) is
compact and \(f\) is locally constant.  Its image spans a
finite-dimensional \(\mathbb F_\ell\)-subspace of \(E\).
Take coordinates in a basis of that span.  Each coordinate
is an additive, locally constant function killing every
inertia intersection, and hence belongs to
\(X_{\mathbb F_\ell}(J)\).
It follows that a prime-field basis of that space spans
\(X_E(J)\) over \(E\).  It remains independent: finitely
many evaluations give a nonsingular minor over
\(\mathbb F_\ell\), which remains nonsingular over \(E\).
This proves (LFC.10), without assuming that \(E\) is finite.
\end{proof}
% END LFC7

% BEGIN LFC8
\paragraph{LFC8. The full strict-generator estimate.}
\label{claim:LFC8}
Here again \(F\) is any number field and \(\ell\) any prime.
Let \(U\subset G_F\), let \(J\triangleleft G_F\) be open with
\(J\subset U\) and \(\ell\nmid[U:J]\), and let
\(\pi_U:U\twoheadrightarrow\Delta\) be continuous.
Here \(\Delta\) is an abstract group equipped with any topology:
compactness, Hausdorffness, and continuity of its group operations
are not required.  The cochains in the conclusion are locally
constant in that topology.
Let \(\kappa:\Delta\to E^\times\) be trivial on \(\pi_U(J)\).
Suppose that the finite local family \(a_v:G_v\to\Delta\)
has the following inertia coverage: for every
\(t\in J\cap I_A\), either \(\pi_U(t)=1\), or
\[
                    \pi_U(t)=g\,a_v(s)g^{-1}
                                                        \tag{LFC.11}
\]
for some \(v,s,g\).  Then the literal strict cone satisfies
\[
 \operatorname{rank}_E H^1_C(a,E(\kappa),0)
 \le \operatorname{clRk}_\ell(J)
           +\#\{v:\kappa a_v=1\}.                       \tag{LFC.12}
\]
This is the original \texttt{str\allowbreak Gen}, including its local
primitive correction and arbitrary coefficient field.

\begin{proof}
Restrict a representative \((z,(x_v))\) along \(\pi_U\) to \(J\).
Because \(\kappa\pi_U=1\) there, the restriction is additive.
It is locally constant and is unchanged by adding a
global coboundary.  It kills inertia.  Indeed if (LFC.11)
holds, then \(\kappa(a_v(s))=1\), while strictness says
\(z(a_v(s))=(\kappa(a_v(s))-1)x_v=0\).
The cocycle identity gives the exact formula
\[
 z(gag^{-1})=\kappa(g)z(a)+(1-\kappa(a))z(g),             \tag{LFC.13}
\]
so \(z(\pi_U(t))=0\).  The case \(\pi_U(t)=1\) is immediate.
We have constructed a linear map from the strict cone
space to the \emph{actual} \(X_E(J)\) of LFC7.

To bound its kernel, put \(N=\pi_U(J)\).  This is normal
of finite index \(n\) dividing \([U:J]\), hence
\(\ell\nmid n\).  If the restricted character is zero,
then \(z\) vanishes on \(N\) and descends to
the finite group \(\Delta/N\), as does \(\kappa\).
Set
\[
           X=-n^{-1}\sum_{q\in\Delta/N}z(q).
\]
The cocycle identity and permutation of the summation
indices give \(z(g)=\kappa(g)X-X\).
Subtract the coboundary pair \((dX,(X)_v)\).
The remaining representative has zero global cocycle,
and its local primitive \(x_v-X\) satisfies
\((\kappa(a_v(s))-1)(x_v-X)=0\) for every \(s\).
It is zero whenever \(\kappa a_v\ne1\).
The kernel is therefore an image of
\(E^{\{v:\kappa a_v=1\}}\); it is not silently identified
with zero when a local invariant exists.
Rank--nullity, or the same dimension inequality for image
and kernel, now gives (LFC.12) using LFC7.
\end{proof}
% END LFC8

% BEGIN LFC9
\paragraph{LFC9. The original \texttt{b1p} bound.}
\label{claim:LFC9}
For a full \texttt{FLink}, assume
\[
 \forall v\ \exists s:\ \Psi(\operatorname{loc}_v s)^2\ne1,
 \qquad
 \operatorname{clRk}_\ell
       (K\cap\ker\bar\epsilon_F)\le C_0.                 \tag{LFC.14}
\]
Then
\[
 \operatorname{rank}_k
 H^1_C\bigl(\operatorname{loc},
               k(\Psi^{-2}\delta_k),0\bigr)\le C_0.
                                                        \tag{LFC.15}
\]

\begin{proof}
Apply LFC8 with \(U=G_F\),
\(J=K\cap\ker\bar\epsilon_F\), \(\Delta=\Gamma\), and
\(\kappa=\Psi^{-2}\delta_k\).
The subgroup \(J\) is normal open and its index divides
\([G_F:K](\ell-1)\), so its index is prime to \(\ell\).
Equation (LFC.2) makes \(\kappa\pi\) trivial on \(J\).
For an inertia element at a place outside \(S\), (LFC.1)
kills its image.  At a place in \(S\), conjugacy of
prolongations sends it to the chosen inertia subgroup,
which lies in \(D_{A_v}=\iota_v(G_v)\).  Applying \(\pi\)
gives precisely (LFC.11).
The trivial valuation has trivial inertia and causes
no extra case.

Because \(\delta\operatorname{loc}_v=1\), local triviality
of \(\kappa\) would imply
\(\Psi(\operatorname{loc}_v s)^2=1\) for all \(s\),
contradicting (LFC.14).
The local correction term in (LFC.12) is therefore zero.
Its remaining term is bounded by the same \(C_0\).
\end{proof}
% END LFC9

% BEGIN LFC10
\paragraph{LFC10. The original doubled-local \texttt{b1} bound.}
\label{claim:LFC10}
Under the same hypotheses (LFC.14),
\[
 \operatorname{rank}_k H^1_C
   \bigl(\operatorname{loc}_2,
         k(\Psi^{-1}|_{\Gamma_0}),0\bigr)\le C_0.
                                                        \tag{LFC.16}
\]

\begin{proof}
Use LFC8 with \(U=\pi^{-1}\Gamma_0\),
\(J=K\cap\ker\bar\epsilon_F\), \(\Delta=\Gamma_0\), and
the surjection induced by \(\pi\).  The same index argument
works for \([U:J]\), and \(K\subset U\) by (LFC.2).
The character is trivial on \(J\).
Outside \(S\) the required inertia images vanish.
At \(v\in S\), a conjugating image \(a\in\Gamma\) either
belongs to \(\Gamma_0\), giving the first local branch,
or is \((ac^{-1})c\), giving the second branch with
conjugator \(ac^{-1}\in\Gamma_0\).
Thus the full inertia coverage holds for the actual
family indexed by \(V\times\{0,1\}\).

Characters are conjugation invariant, so
\[
 \Psi(\operatorname{loc}_{2,(v,b)}s)
                  =\Psi(\operatorname{loc}_v s)
                     \quad(b=0,1).
\]
A trivial local \(\Psi^{-1}\) would again contradict the
nonquadratic witness in (LFC.14).
Both branches have zero local correction in (LFC.12),
which proves (LFC.16).  No identification of this doubled
cone with an undoubled cone has been made.
\end{proof}
% END LFC10

\subsection{A shorter proof of the untwisted bound}

% BEGIN LFC11
\paragraph{LFC11. The cyclotomic eigenspace bound over a CM field.}
\label{claim:LFC11}
Let \(F\) be totally real and let \(\ell\) be odd.
Suppose \(\bar\epsilon_F\) is nontrivial on every
decomposition group above \(\ell\).
Put \(E=F(\mu_\ell)\) and let \(r_\ell(E)\) be the
\(\ell\)-rank of its ideal class group.  Then
\[
 \dim_{\mathbb F_\ell}
   \operatorname{Hom}_{\rm cont}
      (G_{F,S},\mathbb F_\ell)\le 1+r_\ell(E),
                                                        \tag{LFC.17}
\]
where \(S\) is the set of \(\ell\)-places.

\begin{proof}
Let \(\Delta=\operatorname{Gal}(E/F)\) and
\(\omega:\Delta\hookrightarrow\mathbb F_\ell^\times\)
be its cyclotomic character.  Its order is prime to \(\ell\).
Restriction of additive characters from \(G_F\) to \(G_E\)
is injective: a character vanishing on \(G_E\) factors
through \(\Delta\), whose order is prime to \(\ell\).
The Kummer construction in LFC1 sends such a restriction
to a class \([a]\in E^\times/E^{\times\ell}\).
For \(d\in\Delta\), its class satisfies
\[
                         d[a]=\omega(d)[a]               \tag{LFC.18}
\]
in additive notation for this \(\mathbb F_\ell\)-space.
Indeed, if \(g\beta/\beta=\zeta^{\chi(g)}\), a lift of \(d\)
sends \(\zeta\) to \(\zeta^{\omega(d)}\), while
\(\chi(d^{-1}gd)=\chi(g)\).  Thus \(d\beta\) is a
Kummer root for the character \(\omega(d)\chi\),
proving (LFC.18) by the injectivity argument of LFC1.
Unramifiedness outside \(\ell\) places the class in
\(V_T(E)/E^{\times\ell}\), with \(T\) the places above \(\ell\).

Taking \(\omega\)-eigenspaces in (LFC.4) is exact:
the projector is
\[
 e_\omega=|\Delta|^{-1}
           \sum_{d\in\Delta}\omega(d)^{-1}d.
 \]
The class-group term has dimension at most \(r_\ell(E)\).
To see this also for the localized class group, it is a
quotient of \(\operatorname{Cl}(\mathcal O_E)\);
the minimal number of generators of a finite abelian
\(\ell\)-group, equivalently the dimension of its
\(\ell\)-torsion, cannot increase on taking a quotient.

We claim that the \(\omega\)-part of the \(T\)-unit term
has dimension at most one.  The valuation map gives
\(\mathcal O_{E,T}^\times/\mathcal O_E^\times\) as a
subgroup \(\Lambda\) of the permutation lattice
\(\mathbb Z[T]\).  The quotient by the units is torsion-free,
so reducing
\(0\to\mathcal O_E^\times\to
 \mathcal O_{E,T}^\times\to\Lambda\to0\)
modulo \(\ell\) is right exact with the left injection.
For a prime \(w\in T\), its stabilizer in \(\Delta\) is
the image of the local decomposition group in
\(\operatorname{Gal}(E/F)\).  By hypothesis \(\omega\)
is nontrivial on this stabilizer.  A vector of the
permutation representation transforming by \(\omega\)
must therefore have zero coefficient at \(w\), and hence
at every prime in its orbit.

Here it is essential to handle the lattice, rather than
assume that reduction preserves an arbitrary inclusion.
Since \(|\Delta|\) is prime to \(\ell\), the idempotent
\(e_\omega\) lifts to \(\mathbb Z_\ell[\Delta]\) using
the Teichmuller lift of \(\omega\).
Its action on \(\mathbb Q_\ell[T]\) is zero by the same
stabilizer argument.  Therefore its action on
\(\Lambda\otimes\mathbb Z_\ell\) is zero, and reduction
shows that \((\Lambda/\ell\Lambda)_\omega=0\).
Thus the \(\omega\)-part of the \(T\)-units modulo
\(\ell\)-th powers comes entirely from ordinary units.

The field \(E=F(\mu_\ell)\) is CM.  Complex conjugation
\(j\in\Delta\) acts by \(\omega(j)=-1\).
For an ordinary unit \(u\), every complex absolute value
of \(u/j(u)\) is one.  This is an algebraic integer unit
all of whose conjugates have absolute value one, hence
a root of unity.  One elementary proof of this last fact
observes that the monic polynomials having the conjugates
of its positive powers as roots have integral coefficients
bounded uniformly by binomial coefficients.  There are
only finitely many such polynomials and thus finitely many
possible powers; two powers agree.
Consequently \((1-j)\mathcal O_E^\times\subset\mu(E)\).
Modulo \(\ell\), every \(j\)-anti-invariant unit class is
half of its image under \(1-j\), and so lies in the image
of \(\mu(E)/\mu(E)^\ell\).
The roots of unity in a number field form a finite cyclic
group, whose quotient by \(\ell\)-th powers has dimension
at most one.  This proves the claim about units.
Combining it with the class-group bound proves (LFC.17).
\end{proof}
% END LFC11

% BEGIN LFC12
\paragraph{LFC12. The full original \texttt{h1} field.}
\label{claim:LFC12}
For a full \texttt{FLink}, suppose \(k\) is finite,
\(\omega_k\operatorname{loc}_v\ne1\) for every \(v\), and
\(\operatorname{clRk}_\ell(K\cap\ker\bar\epsilon_F)\le C_0\).
Then
\[
 \operatorname{rank}_k
       H^1_C(\operatorname{loc},k(1),k)\le1+C_0.          \tag{LFC.19}
\]
In particular the source's theorem \texttt{GC.\allowbreak h1}, with
its additional retained premise \(P\), follows.

\begin{proof}
With the full local subspaces \(k\), the local primitives
in (LFC.3) are quotiented out.  Since the action is trivial,
global coboundaries are zero.  Thus the space in (LFC.19)
is literally \(\operatorname{Hom}_{\rm cont}(\Gamma,k)\).
For an additive \(\mathbb F_\ell\)-character of \(\Gamma\),
its pullback to \(G_F\) kills every inertia outside \(S\):
for \(t\in I_A\), the prime-to-\(\ell\) power
\(t^{[G_F:K]}\) belongs to \(I_A\cap K\), and (LFC.1)
kills its image.  Hence LFC11 bounds the prime-field
character dimension by \(1+r_\ell(F(\mu_\ell))\).

Put \(J=K\cap\ker\bar\epsilon_F\) and \(E_1=\Omega^J\).
Then \(E_1\supset E=F(\mu_\ell)\) and \([E_1:E]\)
is prime to \(\ell\).
Restriction of unramified additive characters
\(G_E\to\mathbb F_\ell\) to \(G_{E_1}\) is injective.
Indeed, for a character vanishing on the latter subgroup,
the \([E_1:E]\)-th power of every element lies in it
(these are normal subgroups here), and that degree is
invertible modulo \(\ell\).
Restriction remains unramified because inertia is
intersected with \(G_{E_1}\).
Finite class-field reciprocity and LFC7 therefore give
\[
 r_\ell(E)\le\dim_{\mathbb F_\ell}X_{\mathbb F_\ell}(J)
             =\operatorname{clRk}_\ell(J)\le C_0.
 \]
Finally a basis of the finite extension \(k/\mathbb F_\ell\)
shows, by taking coordinates of additive functions, that
\(\operatorname{Hom}_{\rm cont}(\Gamma,k)\) is obtained
from its prime-field character space by scalar extension.
Its dimension over \(k\) is the prime-field dimension.
This gives (LFC.19) with the original \(C_0\).
\end{proof}
% END LFC12

\subsection{Euler estimates through the actual inertia quotient}

% BEGIN LFC13
\paragraph{LFC13. A finite ramification level with its exact kernel.}
\label{claim:LFC13}
For a \texttt{GLink} over any number field there is a
finite set \(S'\supset S\)
and a continuous surjection
\[
 p:G_{F,S'}\twoheadrightarrow\Gamma,\qquad p\,q=\pi,
                                                        \tag{LFC.20}
\]
where \(q:G_F\to G_{F,S'}\), such that
\[
 \ker p=\overline{\langle J_v:v\in S'\setminus S
                                      \rangle^{\rm normal}},
 \qquad J_v=q(I_{A_v}\cap K).                            \tag{LFC.21}
\]
One chosen prolongation at each place suffices in this formula.

\begin{proof}
The finite extension fixed by \(K\) ramifies at finitely
many finite places, say \(T\); this follows by inverting
the nonzero discriminant of an integral basis.
At \(v\notin T\) all inertia is contained in \(K\).
Take \(S'=S\cup T\).
For \(v\notin S'\), (LFC.1) therefore kills all of \(I_A\),
so \(\pi\) factors through the maximal quotient
unramified outside \(S'\), giving (LFC.20).
The quotient topology gives continuity, and surjectivity
comes from that of \(\pi\).

Under \(q\), all generators of (LFC.1) at places outside
\(S'\) become one.  At a place of \(S'\setminus S\),
any prolongation is conjugate to the chosen one, and
normality of \(K\) carries its inertia intersection to
the chosen intersection.  Thus their images are precisely
the normal generators in (LFC.21).
To justify closure, the image of a closed subgroup of a
compact group in a Hausdorff quotient is closed.
The closed normal subgroup generated by the displayed
images has inverse image containing every generator
of \(\ker\pi\); conversely its generators all lie in
\(\ker p\).  These two inclusions prove (LFC.21).
\end{proof}
% END LFC13

% BEGIN LFC14
\paragraph{LFC14. The exact local correction at a discarded prime.}
\label{claim:LFC14}
Here \(\ell\) may be any prime, including two, and the
base number field is arbitrary.
Let \(v\nmid\ell\), let \(D=D_A\), \(I=I_A\), and
\(J=I\cap K\).  Let \(W\) be a finite-dimensional
representation over a finite field \(E\) of characteristic
\(\ell\), with continuous action and trivial \(J\)-action.
Define
\[
 \mathcal E(D,J;W)=
 \{f:J\to W:\ f\text{ continuous and additive},\
                     f(g t g^{-1})=\rho(g)f(t)\}.
 \]
Under the local fields of \(P\), this is finite-dimensional and
\[
              \dim_E\mathcal E(D,J;W)=\dim_E H^2(D,W).
                                                        \tag{LFC.22}
\]
Moreover every class inflated to \(D\) from a continuous
quotient killing \(J\) is zero in \(H^2(D,W)\).

\begin{proof}
Here \([I:J]\) divides \([G_F:K]\), so LFC4 applies.
Restriction of a one-cocycle to \(J\) is an additive
equivariant map and kills coboundaries.
Its kernel in \(H^1(D,W)\) is \(H^1(D/J,W)\):
if the restriction is zero, the cocycle descends
continuously through the compact-to-Hausdorff quotient;
the converse is immediate.

We verify surjectivity of this restriction without
assuming an extension property for arbitrary functions.
For \(f\in\mathcal E(D,J;W)\), choose a normal open
subgroup \(U_0\subset D\) killing the action and satisfying
\(U_0\cap J\subset\ker f\).  Such a subgroup exists
by continuity and compactness.  On the finite quotient
by \(U_0\), choose representatives modulo the image of \(J\)
and define a function \(a:D\to W\) by
\[
                 a(r t)=\rho(r)f(t),\qquad t\in J.
 \]
The condition on \(U_0\cap J\) makes this definition
well defined on the quotient; choose value zero on
the representatives.  It is continuous and satisfies
\(a(g t)=a(g)+\rho(g)f(t)\).
Equivariance of \(f\) shows by substitution that
\(d a(g,h)\) is unchanged on replacing either \(g\) or \(h\)
by a right translate by \(J\).  It is therefore a
two-cocycle on \(D/J\).
By LFC4 it is \(d b\) for a continuous cochain \(b\)
on \(D/J\).  The cocycle \(a-b\) restricts to \(f\).
This proves the exact sequence
\[
 0\longrightarrow H^1(D/J,W)\longrightarrow H^1(D,W)
       \longrightarrow\mathcal E(D,J;W)\longrightarrow0.
 \]
LFC4 gives dimension \(\dim H^0(D,W)\) for its first term.
The local Euler characteristic away from \(\ell\) is
\(\dim H^1(D,W)=\dim H^0(D,W)+\dim H^2(D,W)\).
This proves (LFC.22), with finiteness supplied by \(P\).
For the final assertion, the quotient map from \(D\)
factors through \(D/J\), whose \(H^2\) is zero by LFC4.
The pullback of every two-cocycle class is therefore zero.
\end{proof}
% END LFC14

% BEGIN LFC15
\paragraph{LFC15. The inflation defect is bounded by the same corrections.}
\label{claim:LFC15}
Let \(p:\Gamma'\twoheadrightarrow\Gamma\) be a continuous
surjection of profinite groups whose kernel is the closed
normal subgroup generated by a finite family \(J_v\).
Let \(D_v\) be arbitrary subgroups of \(\Gamma'\), and let
\(W\) be a finite representation of \(\Gamma\) over a finite
field of any prime characteristic.
No containment \(J_v\subset D_v\) or normalization of \(J_v\)
by \(D_v\) is assumed.  Here \(\mathcal E(D_v,J_v;W)\)
means the continuous additive functions on \(J_v\)
satisfying
\(f(g t g^{-1})=\rho(p g)f(t)\) whenever \(g\in D_v\),
\(t\in J_v\), and \(g t g^{-1}\in J_v\).
This is the source's conditional equivariance convention.
Put
\[
 T=\ker\bigl(H^2(\Gamma,W)\to H^2(\Gamma',W)\bigr).
 \]
If \(H^1(\Gamma',W)\) and the spaces
\(\mathcal E(D_v,J_v;W)\) are finite-dimensional, then
\(H^1(\Gamma,W)\) and \(T\) are finite-dimensional and
\[
 \dim T+\dim H^1(\Gamma',W)
   \le\dim H^1(\Gamma,W)+
                          \sum_v\dim\mathcal E(D_v,J_v;W).
                                                        \tag{LFC.23}
\]

\begin{proof}
Inflation on \(H^1\) is injective: a pulled-back
cocycle which is \(g\mapsto\rho(p g)X-X\) has the
same formula downstairs, by surjectivity of \(p\).
Restriction to the \(J_v\)'s defines
\[
 R:H^1(\Gamma',W)\longrightarrow
                  \bigoplus_v\mathcal E(D_v,J_v;W).
 \]
Its kernel is the image of inflation.
Indeed a cocycle vanishing on every \(J_v\) vanishes
on their conjugates by the cocycle formula, and its zero
set inside \(\ker p\) is a closed normal subgroup.
It thus vanishes on \(\ker p\) and descends continuously.

For the second-degree defect, consider the vector space
\(\mathcal T\) of pairs \((z,\beta)\), where \(z\) is a
normalized continuous two-cocycle on \(\Gamma\),
\(\beta(1)=0\), and \(d\beta=p^*z\).
It maps onto \(T\) by \((z,\beta)\mapsto[z]\).
Every element of \(T\) has such a representative:
subtract the coboundary of the constant value \(\beta(1)\)
to normalize a primitive and the corresponding cocycle.
For \(n\in\ker p\), the identity \(d\beta=p^*z\) gives
\[
 \beta(ng)=\beta(n)+\beta(g),\quad
 \beta(gn)=\beta(g)+\rho(p g)\beta(n),\quad
 \beta(gng^{-1})=\rho(p g)\beta(n).                       \tag{LFC.24}
\]
Hence restrictions of \(\beta\) define a linear map
\(r:\mathcal T\to\bigoplus_v\mathcal E(D_v,J_v;W)\).
If \(r(z,\beta)=0\), (LFC.24) and closed normal generation
force \(\beta|_{\ker p}=0\).
The first two identities then make \(\beta\) constant
on fibers, so it descends to a continuous \(b:\Gamma\to W\).
Surjectivity gives \(z=db\), and \([z]=0\).
Thus \(\ker r\) is contained in the kernel of
\(\mathcal T\to T\).
The restriction image \(R(H^1(\Gamma',W))\) is contained
in \(r(\ker(\mathcal T\to T))\), using the pairs \((0,\beta)\).
It follows that \(T\) is a quotient of a subspace of
\[
 \left(\bigoplus_v\mathcal E(D_v,J_v;W)\right)
                    \big/R(H^1(\Gamma',W)).
 \]
In particular it is finite-dimensional.
Since
\(\dim R(H^1(\Gamma',W))=
 \dim H^1(\Gamma',W)-\dim H^1(\Gamma,W)\),
the displayed quotient bound is exactly (LFC.23).
This argument uses actual continuous primitives and
requires no global continuous section of \(p\).
\end{proof}
% END LFC15

% BEGIN LFC16
\paragraph{LFC16. The finite-level Euler and Poitou--Tate identities.}
\label{claim:LFC16}
Let \(F\) be totally real, \(\ell\) odd,
\(S'\) finite containing every \(\ell\)-place, and \(W\)
a finite-dimensional continuous representation of
\(G_{F,S'}\) over a finite field \(E\) of characteristic
\(\ell\).  The retained fields of \(P\) give finite
\(H^1,H^2\) and the identities
\[
 \begin{split}
 \dim H^1+\sum_{w\mid\infty}\dim W^{c_w=1}
   &=\dim H^0+\dim H^2+[F:\mathbb Q]\dim W,\\
 \dim\operatorname{Sha}^2+\sum_{v\in S'}\dim H^2(D_v,W)
   &=\dim H^2+\dim H^0(G_F,W^\vee(1)).
 \end{split}                                           \tag{LFC.25}
\]
Here \(\operatorname{Sha}^2\) is precisely the kernel of localization
to the displayed \(H^2(D_v,W)\)'s.

\begin{proof}
Apply \(P\)'s local and global finiteness and its global
Euler characteristic at exponent \(\ell^1\), after
restricting scalars to \(\mathbb F_\ell\).
The cocycle and coboundary sets do not change under this
restriction: their defining equations and continuity
conditions are identical.  If \(q=|E|\), each finite
cohomology space therefore has cardinality \(q^{\dim_E}\).
The Euler characteristic's cardinal identity becomes
the first equality after taking exponents of \(q>1\).

The last two maps in the normalized Poitou--Tate sequence
provided by \(P\) are exact, and the last is onto:
\[
 H^2(G_{F,S'},W)\longrightarrow
 \bigoplus_{v\in S'}H^2(D_v,W)
       \longrightarrow H^0(G_F,W^\vee(1))^\vee
       \longrightarrow0.
 \]
Its first kernel is the stated \(\operatorname{Sha}^2\), by the
literal restriction maps on two-cocycles.
Taking cardinalities gives the second identity.
The coefficient-field dual is correctly identified as
follows.  Composition with
\(\operatorname{Tr}_{E/\mathbb F_\ell}\) identifies the
underlying additive space of \(\operatorname{Hom}_E(W,E)\)
with \(\operatorname{Hom}_{\mathbb F_\ell}(W,\mathbb F_\ell)\).
Injectivity follows from nondegeneracy of the finite-field
trace pairing, and equal cardinalities give surjectivity.
The map commutes with the representation and the
cyclotomic twist, so it identifies the invariant spaces.
This supplies the exact dual term in (LFC.25), without
assuming an unintended \(E\)-linearity of the original
prime-field duality normalization.
\end{proof}
% END LFC16

% BEGIN LFC17
\paragraph{LFC17. The global estimate for the actual quotient.}
\label{claim:LFC17}
For a \texttt{GLink}, a totally real \(F\), and a
finite-dimensional finite-field representation \(W\)
of \(\Gamma\) with continuous action, \(H^1(\Gamma,W)\)
and \(H^2(\Gamma,W)\) are finite-dimensional, and
\[
 \begin{split}
 &\dim\operatorname{Sha}^2_S(\Gamma,W)+
       \sum_{v\in S}\dim H^2(D_v,W)+\dim H^0(\Gamma,W)
                          +[F:\mathbb Q]\dim W\\
 &\hspace{12mm}\le
       \dim H^1(\Gamma,W)+
       \sum_{w\mid\infty}\dim W^{\pi(c_w)=1}
       +\dim H^0(G_F,W^\vee(1)).
 \end{split}                                           \tag{LFC.26}
\]
The local actions on the left are through \(\pi\), and
\(\operatorname{Sha}^2_S\) is their exact localization kernel.

\begin{proof}
Choose \(S',p\) by LFC13.
Let \(a_i=\dim H^i(\Gamma,W)\),
\(a_i'=\dim H^i(G_{F,S'},W)\),
\(t=\dim\ker(H^2(\Gamma,W)\to H^2(G_{F,S'},W))\),
and \(b_v=\dim H^2(D_v,W)\).
At a discarded place \(v\in S'\setminus S\), pullback
injects the equivariant homomorphisms on
\((qD_v,q(I_v\cap K))\) into those on \((D_v,I_v\cap K)\).
LFC14 bounds their dimension by \(b_v\).
LFC15 consequently gives
\[
                     t+a_1'\le a_1+\sum_{S'\setminus S}b_v.
                                                        \tag{LFC.27}
\]
It also gives finiteness of \(a_1\) and \(t\).
The target \(H^2(G_{F,S'},W)\) is finite by LFC16;
its finite kernel \(T\) then proves finiteness of \(a_2\).

Inflation sends \(\operatorname{Sha}^2_S(\Gamma,W)\) into
\(\operatorname{Sha}^2_{S'}(G_{F,S'},W)\).
At places in \(S\) this is the definition; at the other
places LFC14 says that all inflated classes vanish
locally, because \(\pi\) kills \(I_v\cap K\).
Thus
\[
             \dim\operatorname{Sha}^2_S\le t+\dim\operatorname{Sha}^2_{S'}.
 \]
Use the second equality in (LFC.25) to replace
\(\dim\operatorname{Sha}^2_{S'}\), and then the first to replace
\(a_2'\).  The invariants and the real fixed spaces
are unchanged by the surjection \(p\).
The resulting discarded-place sum is cancelled by
(LFC.27).  Explicitly the left side of (LFC.26) is at most
\[
 t+a_2'+\dim H^0(G_F,W^\vee(1))
     -\sum_{S'\setminus S}b_v+a_0+[F:\mathbb Q]\dim W,
 \]
which is at most its right side by these two identities.
This proves the estimate without erasing any discarded
prime or assuming inflation on \(H^2\) is injective.
\end{proof}
% END LFC17

% BEGIN LFC18
\paragraph{LFC18. Local coefficient quotients preserve the required bound.}
\label{claim:LFC18}
In LFC17 suppose \(H^0(G_F,W^\vee(1))=0\).
For every family of equivariant surjections
\(q_v:W\twoheadrightarrow Q_v\) for the local groups,
let
\[
 X_Q=\ker\left(H^2(\Gamma,W)\longrightarrow
                              \bigoplus_{v\in S}H^2(D_v,Q_v)\right).
 \]
Then
\[
 \begin{split}
 \dim X_Q+\sum_{v\in S}\dim H^2(D_v,Q_v)
       +\dim H^0(\Gamma,W)+[F:\mathbb Q]\dim W\\
 \le \dim H^1(\Gamma,W)
       +\sum_{w\mid\infty}\dim W^{\pi(c_w)=1}.
 \end{split}                                           \tag{LFC.28}
\]
No flatness, choice of a local complement preserved by
the action, or global stability of \(\ker q_v\) is required.

\begin{proof}
Local Tate duality identifies \(H^2(D_v,W)\) with the
dual of \(\operatorname{Hom}_{D_v}(W,E(1))\), and similarly
for \(Q_v\).  These perfect finite pairings are natural
in the coefficients: on cochains this is the equality
\(f(q_v x)=(f\circ q_v)(x)\) in the evaluation pairing.
Thus the map on \(H^2\) induced by \(q_v\) is dual to the
injection
\(\operatorname{Hom}_{D_v}(Q_v,E(1))
 \hookrightarrow\operatorname{Hom}_{D_v}(W,E(1))\).
It is therefore onto.  The same conclusion follows after
the trace identification in LFC16, so it holds with the
precise prime-field duality contained in \(P\).

Here is the needed linear algebra.  For maps
\(r_v:A\to B_v\) and surjections \(f_v:B_v\to C_v\),
put \(X=\bigcap\ker(f_vr_v)\) and \(Y=\bigcap\ker r_v\).
Restriction of \((r_v)_v\) gives an injection
\(X/Y\hookrightarrow\bigoplus_v\ker f_v\).
Hence
\[
          \dim X+\sum_v\dim C_v
                   \le\dim Y+\sum_v\dim B_v.
 \]
Apply this to \(A=H^2(\Gamma,W)\), the local
\(H^2(D_v,W)\), and their quotient-coefficient images.
LFC17 with its zero dual term gives (LFC.28).
\end{proof}
% END LFC18

\subsection{From ordinary cohomology to the literal cone}

% BEGIN LFC19
\paragraph{LFC19. The full cone dimension identity.}
\label{claim:LFC19}
This claim is independent of the arithmetic linked data.
Let \(\Gamma\) and \(G_v\) be arbitrary profinite groups,
with finitely many \(v\), continuous maps
\(\operatorname{loc}_v:G_v\to\Gamma\), a finite field \(E\),
and a finite-dimensional continuous representation \(W\) of
\(\Gamma\) over \(E\).
Let \(q_v:W\twoheadrightarrow Q_v=W/b_v\) be equivariant
for each local group, and suppose the global \(H^1,H^2\)
and local \(H^1(G_v,Q_v)\) are finite-dimensional.
Equivalently one may start with arbitrary equivariant
surjections \(q_v\), taking \(b_v=\ker q_v\) and identifying
the targets with \(W/b_v\); all the maps below respect
that canonical identification.
Write
\[
 \begin{split}
 H^0_C&=\{x\in W^\Gamma:x\in b_v\text{ for every }v\},\\
 X_Q&=\bigcap_v\ker\bigl(H^2(\Gamma,W)\to H^2(G_v,Q_v)\bigr).
 \end{split}
 \]
Then the literal \(H^1_C,H^2_C\) are finite-dimensional and
\[
 \begin{split}
 &\dim H^2_C+\dim H^0_C+\dim H^1(\Gamma,W)
                       +\sum_v\dim H^0(G_v,Q_v)\\
 &\qquad=\dim H^1_C+\dim H^0(\Gamma,W)
                       +\sum_v\dim H^1(G_v,Q_v)+\dim X_Q.
 \end{split}                                           \tag{LFC.29}
\]

\begin{proof}
We use the cochain model, so that the local primitives
are accounted for rather than discarded.
A degree-two cone representative is a global continuous
two-cocycle \(z\), together with local continuous
one-cochains \(a_v:G_v\to W\), satisfying
\[
 q_v z(\operatorname{loc}_v s,\operatorname{loc}_v t)
                         =d(q_v a_v)(s,t).
 \]
Its coboundaries have the form
\[
 z=d\beta,\qquad
 a_v(s)\equiv
   \beta(\operatorname{loc}_v s)
        -\bigl(\rho(\operatorname{loc}_v s)Y_v-Y_v\bigr)
                                    \pmod {b_v}.        \tag{LFC.30}
\]
These are exactly the source's
\texttt{continuous\allowbreak Cocycles2} and
\texttt{continuous\allowbreak Coboundaries2}.
Choose arbitrary linear sections of the \(q_v\)'s;
they are continuous for the discrete topologies and
are not required to intertwine the action.

The definitions give the following exact sequence, ending
with the indicated image:
\[
 \begin{gathered}
 0\longrightarrow H^0_C\longrightarrow H^0(\Gamma,W)
 \longrightarrow \bigoplus_v H^0(G_v,Q_v)
 \longrightarrow H^1_C\longrightarrow H^1(\Gamma,W)\\
 \longrightarrow\bigoplus_v H^1(G_v,Q_v)
 \longrightarrow H^2_C\longrightarrow X_Q
 \longrightarrow0.                                    \tag{LFC.31}
 \end{gathered}
\]
Here an invariant local tuple maps to the class of
\((0,(x_v))\) after lifting it; a local one-cocycle
tuple maps to the cone class \((0,(a_v))\) after lifting.
The maps from the cone groups forget the local data.
Changing a lift changes its value by \(b_v\), hence
does not change these classes.

For clarity, exactness can be checked directly at every
nontrivial step.  A local invariant tuple is a cone
coboundary precisely when it is the image of one global
invariant \(X\), by (LFC.3).
If the global cocycle of a cone-one class is \(dX\),
subtract its coboundary; the remaining local primitives
are invariant in \(Q_v\), giving the preceding image.
A global one-class has zero local classes precisely
when its local coboundary primitives can be lifted to
form a cone-one representative.
A local one-cocycle tuple gives zero in \(H^2_C\)
precisely when (LFC.30) supplies a global one-cocycle
\(\beta\) whose local cohomology classes equal that tuple.
If a cone-two class has zero global class, subtract
the coboundary of a global primitive \(\beta\);
the remaining \(q_v(a_v-\beta\operatorname{loc}_v)\)
are local one-cocycles.
Finally a global two-class lies in \(X_Q\) precisely
when local continuous primitives exist in each \(Q_v\);
lift them to obtain a cone-two representative.
This proves (LFC.31).

Every term except possibly the cone terms is
finite-dimensional; exactness proves their finiteness
as well.  The alternating dimension sum of (LFC.31)
is zero, and rearranging it is exactly (LFC.29).
\end{proof}
% END LFC19

% BEGIN LFC20
\paragraph{LFC20. The complete \texttt{EC} field for every finite module.}
\label{claim:LFC20}
Assume \(P\) and a full \texttt{FLink}.
For every finite field \(E\) of characteristic \(\ell\),
finite-dimensional \(E\)-space \(W\), representation
\(\rho:\Gamma\to\operatorname{End}_E(W)\) with open kernel,
and locally stable subspaces \(b_v\subset W\), suppose
\[
 \bigl[f(\rho(g)x)=\omega_E(g)f(x)
           \text{ for all }g,x\bigr]\Longrightarrow f=0
                   \quad(f:W\to E\text{ linear}).        \tag{LFC.32}
\]
Then \(H^2_C\) is finite-dimensional and
\[
 \dim_E H^2_C+\dim_E H^0_C+\sum_v d_v\dim_E b_v
 \le \dim_E H^1_C+
                 \sum_{w\in\mathrm{Inf}}\dim_E W^{c_w=1}.
                                                        \tag{LFC.33}
\]
This is the full source \(\Phi.\texttt{EC}\), with its
quantifiers over coefficient fields, representations,
and all stable local conditions.

\begin{proof}
An open kernel makes the action continuous: for fixed
\(x\), its value is constant on cosets of that kernel;
the finite discrete space \(W\) then gives joint
continuity.  Condition (LFC.32), surjectivity of \(\pi\),
and \(\epsilon\pi=\epsilon_F\) say exactly that
\(H^0(G_F,W^\vee(1))=0\).
Put \(Q_v=W/b_v\).
The maps \(\iota_v:G_v\simeq D_{A_v}\) are homeomorphisms,
since they are continuous bijections from compact
groups to Hausdorff groups.  Pullback and its inverse
identify their continuous cocycles, coboundaries,
invariants and cohomology.  Thus LFC18 applies to
the actual local family of the frame.

Local Euler characteristic at \(v\mid\ell\) gives
\[
 \dim H^1(G_v,Q_v)+d_v\dim b_v
 =\dim H^0(G_v,Q_v)+\dim H^2(G_v,Q_v)+d_v\dim W.
                                                        \tag{LFC.34}
\]
All local groups in this equation are finite-dimensional.
The local degree formula, with every place above \(\ell\)
present exactly once, gives
\(\sum_v d_v=[F:\mathbb Q]\).
Use (LFC.34) in (LFC.29).
The combination on its right involving \(X_Q\), local
\(H^2\), global \(H^0\), and \([F:\mathbb Q]\dim W\)
is bounded by (LFC.28).
Cancel the remaining global \(H^1\) and local \(H^0\)
terms.  This is precisely (LFC.33), with the real
fixed spaces identified by the given bijection
\(\mathrm{Inf}\simeq\{\text{real places of }F\}\).
Finiteness of \(H^2_C\) follows from LFC17 and LFC19.
In particular no chosen local condition was replaced
by \(0\) or by \(W\) in proving this field.
\end{proof}
% END LFC20

\subsection{Assembly without changing the field, marking, or bound}

% BEGIN LFC21
\paragraph{LFC21. All twenty-eight \texttt{Ok} fields on the MFC frame.}
\label{claim:LFC21}
Fix the full original \(\texttt{Gen}\), \(P\), and the
single \(L,g_1,\Phi,C_0,w\) constructed in MFC14.
Then this very frame satisfies
\[
                         \Phi.\texttt{Ok}(C_0)
                    \quad\hbox{and}\quad
                         \Phi.\texttt{Big}(C_0).          \tag{LFC.35}
\]

\begin{proof}
Let \(F=L^+\), \(U=G_F\subset G_{\mathbb Q}\), and
\(\alpha=\overline\Psi_{\mathbb Q}\).
The \texttt{FLink} supplied by the \texttt{Linked} witness
has
\[
 K=(\ker\eta\cap\ker\alpha)|_{G_F},\qquad
 J=K\cap\ker\bar\epsilon_F.
 \]
Under the actual field-embedding map, \(J\) is precisely
\[
     U\cap\ker\eta\cap\ker\alpha\cap\ker\bar\epsilon_{\mathbb Q}.
                                                        \tag{LFC.36}
\]
This is an equality of subgroups after the natural
topological isomorphism, not just an index comparison.
Inertia groups are carried to intersections with the
same valuation inertia, and the cyclotomic characters
agree on every root of unity.  Precomposition therefore
identifies the literal \(\texttt{XDual}(1,-,\mathbb F_\ell,
\mathbb F_\ell)\) spaces, including continuity and
every inertia condition.
Their ranks are equal.  Thus MFC5's uniform bound,
at MFC's already chosen layer \(m\), is exactly
\(\operatorname{clRk}_\ell(J)\le C_0\) required above.

MFC4 and MFC8 give a nontrivial local
\(\bar\epsilon_F\), at every local group, and MFC9
gives the nonquadratic residual local witness
\(\Psi(\operatorname{loc}_v s)^2\ne1\).
The conversion of cyclotomic reduction to
\(\omega_k\) uses the injective prime-field map,
as in LFC6, and does not strengthen a residue condition.
Apply LFC2 for \texttt{phi}, LFC6 for \texttt{gen},
LFC3 for \texttt{cl}, LFC20 for \texttt{ec},
LFC9 for \texttt{b1p}, LFC10 for \texttt{b1},
and LFC12 for \texttt{h1}, using this one \(C_0\).

The twenty-one already proved fields are
\[
 \begin{gathered}
 \texttt{l\_\allowbreak odd},\ \texttt{fin},\ \texttt{del\_\allowbreak open},\
 \texttt{c\_\allowbreak sq},\ \texttt{del\_\allowbreak c},\ \texttt{loc\_\allowbreak cont},\
 \texttt{eps\_\allowbreak cont},\ \texttt{Psi\_\allowbreak open},\ \texttt{Psi\_\allowbreak c},\\
 \texttt{base\allowbreak Subgroup\_\allowbreak spec},\ \texttt{card\_\allowbreak inf},\
 \texttt{cw\_\allowbreak sq},\ \texttt{del\_\allowbreak cw},\
 \texttt{f2a},\ \texttt{f2b},\ \texttt{f2c},\
 \texttt{f3},\ \texttt{f4},\ \texttt{f4'},\
 \texttt{f5},\ \texttt{Psi\_\allowbreak N}.
 \end{gathered}
 \]
Here the displayed \texttt{Psi} is the source's
\(\Psi\)-named field.  MFC13 supplies both nonzero
strict cone classes for \texttt{f4}, \texttt{f4'} on
this same quotient.  Together with the seven just
proved, these are exactly all twenty-eight fields.
MFC9 and MFC14 already proved
\(3C_0+7<\sum_v d_v\) and \(d_v>5\) for every \(v\),
which are the two \texttt{Big} requirements.
They are unchanged.  This proves (LFC.35).
\end{proof}
% END LFC21

% BEGIN LFC22
\paragraph{LFC22. The full original \texttt{Has\allowbreak Frame\allowbreak W} constructor.}
\label{claim:LFC22}
For every commutative local ring \(O\) whose residue
field has characteristic \(\ell\), the full original
\(\texttt{Gen}(\ell,\eta,\Psi,\sigma,B,M,H_b)\) and \(P\)
imply, for every \(Q>0\),
\[
                \texttt{Has\allowbreak Frame\allowbreak W}
                     (\ell,\eta,\Psi,\sigma,B,M,Q).
                                                        \tag{LFC.37}
\]
The output is one CM field \(L\subset\Omega\), one
marked \(g_1\in G_L\cap D_B\), one frame \(\Phi\),
and \(C_0,w\), with all the literal output clauses.

\begin{proof}
Fix \(Q\) first.  MFC1 chooses the real quadratic
subgroup \(H_0\).  CAF5 chooses a witness prime \(w\)
with \(w\equiv1\pmod Q\); its witness conditions are
the complete \texttt{Frobenius\allowbreak Witness\allowbreak Prime}.
MFC5 chooses \(C_0\) uniformly for every later layer,
using TW1 at the exact fixed abelian field.
Only then set \(m=3C_0+7\), so
\(3C_0+7<\ell^m\).
MFC4, MFC6 and CAF8 construct the actual CM fixed
field with
\[
  U=H_0\cap H_\theta(w)\cap H_b\cap H_m,\qquad
  G_L=U\cap\ker\eta,\qquad G_{L^+}=U.
 \]
The nonquadratic local element provided there is the
marked \(g_1\in D_B\cap G_L\).
MFC8 constructs \(\Phi\) and the full \texttt{Linked}
witness for this \(L,g_1\); MFC12--MFC13 supply both
strict classes, and LFC21 gives its full \texttt{Ok}
and \texttt{Big} at the already chosen \(C_0\).
Finally \(G_{L^+}=U\subset H_\theta(w)\) follows directly
from the displayed intersection.

Thus the same tuple has, in the order of the definition:
\(\texttt{Linked}\), \(\Phi.\texttt{Ok}(C_0)\),
\(\Phi.\texttt{Big}(C_0)\), the full witness-prime record,
\(w\equiv1\pmod Q\), and
\(\operatorname{upl}(\Omega,L)\subset H_\theta(w)\).
These are exactly all the fields of \texttt{Has\allowbreak Frame\allowbreak W}.
The finite residue-field requirement was supplied by
\texttt{Gen.\allowbreak fin}; no DVR or characteristic-zero condition
on \(O\) was used for (LFC.37).
\end{proof}
% END LFC22

% BEGIN LFC23
\paragraph{LFC23. The original coefficient-DVR \texttt{Has\allowbreak Frame} join.}
\label{claim:LFC23}
Specialize to the original coefficient DVR \(O=c.O\)
and retain the full \texttt{Gen} and \(P\).
The construction gives \(\texttt{Has\allowbreak Frame}\), with
the character \(\texttt{Seed\allowbreak Clause}\) on the same
\(L,\Phi,g_1\) as its complete \texttt{Ok} and
\texttt{Big} fields.

\begin{proof}
CAF7, applied to the full \texttt{Gen}'s seed-relevant
fields in this coefficient-DVR setting, chooses one
positive integer \(Q\) and the core finite characters
before the later compatible \(L\) and \(w\).
Apply LFC22 to this \(Q\).  The produced \(w\) has
the stipulated congruence, and the resulting \(L\)
has the required embedded fixed-field description
and \(\operatorname{upl}(\Omega,L)\subset H_\theta(w)\).
CAF7 therefore supplies the complete original
\texttt{Seed\allowbreak Clause} for that very field.
CAF11's literal installation then returns the
\texttt{Gal\allowbreak Link}, the complete \(\texttt{Ok}\),
\(\texttt{Big}\), and this \texttt{Seed\allowbreak Clause} with
the same marked element and frame.
This is the defining conjunction of \texttt{Has\allowbreak Frame}.
The statement is restricted to \(O=c.O\) because that
is CAF7's original coefficient-DVR scope.
Its character seed is not an automorphic rank-two
or rank-three representation.
\end{proof}
% END LFC23

\paragraph{Exact remaining arithmetic boundary.}
The seven missing \texttt{Ok} fields, the full
\texttt{frame\_\allowbreak ex}/\texttt{Has\allowbreak Frame\allowbreak W} theorem, and
the coefficient-DVR \texttt{Has\allowbreak Frame} join have now
been proved at their original premises.
To obtain \texttt{Gd\allowbreak Frame\allowbreak Ex} and \texttt{Host\allowbreak Frame\allowbreak Ex}
from their named data, one must still supply the full
entry records \texttt{gen\_\allowbreak gd} and \texttt{gen\_\allowbreak host}.
The latter requires the actual
\(H_b=\ker(\operatorname{dirG}(\kappa_B))\), its
\(\ell\)-power index, its inertia containment and exact
Teichmuller lifting, as well as the host's local
nonquadratic witness.  The host case continues to
allow \(\ell=3\); no stronger lower bound on \(\ell\)
has been introduced in this chapter.
Both automorphic seed producers, the shared coefficient
extension, the original/alternative endpoint matching,
and the full common \texttt{Gd\allowbreak End} remain separate work.
No whole selected compound obligation or complete
human proof of RH or Dirichlet GRH is closed by this
chapter alone.


\section{All tensor-map derivatives and the complementary Capelli operators}
\label{sec:bessel-shift-suppliers-v51}

This chapter supplies the whole original \texttt{TMap\allowbreak Statement} and
\texttt{TMap\allowbreak DStatement}, the two complementary Capelli identities, and
\texttt{Eig\allowbreak Statement}.  In particular, the estimates apply to every
tuple from the generated module and to every derivative word, with no
loss in the original exponent.  The final claim isolates the remaining
period-evaluation and shifted-quotient inputs to the two unit shifts.
Their existence is not inferred from the algebra below.

We use the reviewed DC, BNV and BSP notation.  Thus
$G=U(J_\epsilon)$, $K=U(p)\times U(q)$, $\rho(g)=\|g\|_{\rm F}$,
$q_0=\min(p,q)$, and
\[
 a=\frac{h}{8\sqrt{2q_0}},\qquad \beta=a-1,\qquad
 \Lambda=\lambda+\rho_n-\rho_c.
 \tag{BST.1}
\]
Statements about arbitrary dimension or arbitrary algebra below do not
impose the Bessel signature.  In the tensor-map application
$n=2r+k+1$, $p=k+r$, $q=r+1$, and the original \texttt{Step\allowbreak Gap}
is assumed.  BSP1 then gives $q_0\ge1$, $h>h_{\rm need}$, and
$\beta>4\max(2nq_0,2rn^2)$.

Our inner product is conjugate-linear in its first slot.  Write
$R_XF(g)=\frac{d}{dt}|_{t=0}F(g\exp(tX))$ for a real Lie direction.
The source's $E_{ab}$ on functions is its complex-linear extension to
the matrix unit.  For a word $\mathbf w=(w_1,\ldots,w_\ell)$,
$E_{\mathbf w}=E_{w_1}\cdots E_{w_\ell}$, with the rightmost
operator applied first.  Smoothness on $G$ always means the original
local extension to a smooth real function on ambient matrices.
No claim requires the originally given off-group values to be smooth.

\subsection{Compact smoothing without loss of decay}

% BEGIN BST1
\paragraph{BST1. A compact smooth convolution fixing the entire lowest type.}
\label{claim:BST1}
For arbitrary $n,\epsilon$ and a kernel datum
$(\lambda,h,T,\mu)$ at the original
Harish--Chandra, gap, $q_0\ge1$, $h>h_{\rm need}$ and
$\mu=\lambda+\rho$ hypotheses, there exists one smooth compactly
supported function $b$ on ambient complex matrices such that
\[
 \Psi_{T,\mu}(g)
    =\int_G b(y)\Psi_{T,\mu}(gy)\,dy\qquad(g\in G).
 \tag{BST.2}
\]
The Haar normalization is arbitrary.  The same $b$ works for every
pair of vectors in $T$.

\begin{proof}
Take the normalized Hilbert model of DC1--DC11, writing $j:T\to
\mathcal H$ for its isometric lowest-type embedding.  Its coefficients
are exactly $\Psi$ and the full $T$-isotypic subspace is $j(T)$.
The latter assertion is the multiplicity and closure conclusion of
DC9--DC10, not an additional admissibility assumption.

Choose a nonnegative smooth approximate identity $b_0$ on $G$, of
compact support and integral one.  Average it under conjugation by
$K$:
\[
 b_1(y)=\int_K b_0(k^{-1}yk)\,dk ,
 \tag{BST.3}
\]
where compact Haar has mass one.  This is smooth, and its support is
contained in the compact set $K(\operatorname{supp}b_0)K^{-1}$.
Conjugation by compact elements preserves group Haar measure: its
modulus is a continuous homomorphism from a compact group to
$\mathbb R_{>0}$ and hence equals one.
The bounded operator $A=\int_G b_1(y)\sigma(y)\,dy$ commutes with
$K$.  It consequently carries $j(T)$ into that same full isotypic
subspace, and Schur's lemma makes its restriction a scalar $c$.

By strong continuity, the operators obtained from a sufficiently small
$b_0$ are arbitrarily close to the identity on this finite-dimensional
subspace, uniformly on its unit sphere.  Averaging the conjugates
preserves that estimate, because $K$ preserves $j(T)$ isometrically.
Choose the support small enough that $|c-1|<1/2$.  Then $c\ne0$,
and $b=b_1/c$ satisfies
$\int_G b(y)\sigma(y)j(v)\,dy=j(v)$ for every $v\in T$.
Pairing after $\sigma(g)$ with $j(v')$ proves (BST.2), first as
scalar coefficients and hence as an equality of finite-dimensional
endomorphisms.  Every integral here is absolutely convergent because
its integrand has compact support and the representation is unitary.

Finally, a compactly supported smooth function on the closed embedded
submanifold $G$ has a compactly supported smooth ambient extension.
Indeed cover its compact support by finitely many submanifold charts,
use a smooth partition of unity, extend each chart function constantly
in the normal coordinates, and multiply by a compact bump function
equal to one near its support.  Sum these finitely many extensions.
Replacing $b$ by this extension changes no integral in (BST.2).
A change of Haar normalization is absorbed by the inverse scaling
of $b$.
\end{proof}
% END BST1

% BEGIN BST2
\paragraph{BST2. The whole original \texttt{Dec\allowbreak Deriv\allowbreak Statement}.}
\label{claim:BST2}
At every datum of BST1, for every matrix-unit word $\mathbf w$
and every $v_1,v_2\in T$ there exists a constant $C_{\mathbf w,v_1,v_2}$
such that
\[
 \left|E_{\mathbf w}
       \langle v_2,\Psi_{T,\mu}(g)v_1\rangle\right|
       \le C_{\mathbf w,v_1,v_2}\rho(g)^{-a}
       \qquad(g\in G).
 \tag{BST.4}
\]
The constant may depend on the word and vectors; the exponent does not.

\begin{proof}
We first recall explicitly how DC6 gives the required undifferentiated
Frobenius bound at this arbitrary signature.  Write $g=ke^X$,
$s=\|X\|_{\rm F}$ and $R=\sqrt{2q_0}\ge1$.  DC6 gives
\[
 \|\Psi(g)\|\le
  \exp\left[-\frac h8\min(s^2,s/R)\right],
 \qquad \rho(g)\le\sqrt n\,e^s.
\]
The elementary inequality
$\min(s^2,s/R)\ge(s-1)/R$ holds: if $sR\ge1$ the minimum
is $s/R$, and otherwise $s<1$ and the right side is nonpositive.
Thus
\[
           \|\Psi(g)\|\le e^a n^{a/2}\rho(g)^{-a}.
 \tag{BST.5}
\]
Here $a>0$, since the displayed large-gap hypothesis implies $h>0$.

Let $\phi(g)=\langle v_2,\Psi(g)v_1\rangle$.
By BST1, $\phi(g)=\int_G b(y)\phi(gy)\,dy$.
For a real Lie direction $X$, left Haar invariance gives
\[
 \phi(g\exp(tX))
    =\int_G b(\exp(-tX)y)\phi(gy)\,dy.
 \tag{BST.6}
\]
For $t$ in a fixed small interval the supports on the right lie in
one compact set.  All derivatives of the smooth factor $b$ are
bounded there.  For $g$ in a fixed compact neighborhood, continuity
of $\phi$ bounds the other factor as well.  Dominated differentiation
therefore applies to every order, locally uniformly in $g$.
In particular, set
$D_Xb(y)=\frac{d}{dt}|_{0}b(\exp(tX)y)$; the kernel for $R_X\phi$
is $-D_Xb$.  Extend this operation complex linearly to each
matrix unit and iterate in the specified word order.  We obtain a
smooth function $b_{\mathbf w}$ with support contained in
$\operatorname{supp}(b|_G)$ and
\[
 E_{\mathbf w}\phi(g)=
           \int_G b_{\mathbf w}(y)\phi(gy)\,dy .
 \tag{BST.7}
\]
The same construction may be performed on the ambient extension of
$b$; its directional differential operators have smooth coefficients
and preserve compact support.  Thus it also has the source's precise
ambient meaning.

On that fixed compact support, let $B$ bound $\rho(y^{-1})$.
Frobenius submultiplicativity gives
$\rho(g)\le\rho(gy)\rho(y^{-1})$, hence
$\rho(gy)^{-a}\le B^a\rho(g)^{-a}$.
Apply (BST.5) inside (BST.7) and integrate
$|b_{\mathbf w}|$:
\[
 |E_{\mathbf w}\phi(g)|
 \le e^a n^{a/2}\|v_2\|\|v_1\|\,
       B^a\|b_{\mathbf w}\|_{L^1}\rho(g)^{-a}.
\]
This is (BST.4).  It proves the complete source statement for all
words, including the empty word.  No estimate for a differentiated
kernel was assumed.
\end{proof}
% END BST2

% BEGIN BST3
\paragraph{BST3. Every generated-module coefficient, with all words and translates.}
\label{claim:BST3}
Take any of the original complete unitary derived data
$(\mathcal H,\sigma,D,E,V,T,\mu,\iota)$ in
\texttt{TMap\allowbreak Statement}.  In particular its lowest-type coefficients
equal the specified $\Psi$ kernel, its derived domain is invariant,
and all the original gap hypotheses hold.
Put $S=U(\mathfrak{gl}_n)V\subset D$.
For every $x_0\in G$, every $w\in V$ and every $m\in S$, the
function
\[
 F_m(g)=\langle w,\sigma(x_0^{-1}g)m\rangle
 \tag{BST.8}
\]
is smooth in the original local sense.  For every word $\mathbf w$,
\[
 E_{\mathbf w}F_m=F_{E_{\mathbf w}m},\qquad
 |E_{\mathbf w}F_m(g)|\le C_{m,\mathbf w,x_0,w}\rho(g)^{-a}.
 \tag{BST.9}
\]
No irreducibility of the ambient representation is required.

\begin{proof}
The algebraic subspace spanned by all finite words $E_{\boldsymbol\alpha}v$,
$v\in V$, contains $V$ and is stable under every $E_{ab}$.
Conversely every such stable subspace contains every word and their
finite linear combinations.  It is therefore exactly the source's
generated module $S$.

For $m=v\in V$, the original coefficient identity makes (BST.8)
equal on $G$ to
$\langle\iota w,\Psi(x_0^{-1}g)\iota v\rangle$.
Its smoothness follows from DC2.  Right differentiation commutes with
the fixed left translate $x_0^{-1}$.
The defining derived-representation identity gives
$E_{ab}F_m=F_{E_{ab}m}$, first for real directions by differentiation
of $\sigma(x_0^{-1}g)\sigma(\exp(tX))m$, and then by complexification.
Iterating gives the first formula of (BST.9).
For $m=E_{\boldsymbol\alpha}v$, the resulting coefficient is a
derivative of the preceding smooth kernel coefficient; its further
$\mathbf w$ derivative is the kernel derivative for the concatenated
word $\mathbf w\boldsymbol\alpha$.
BST2 bounds it by a constant times $\rho(x_0^{-1}g)^{-a}$.
Since $\rho(g)\le\rho(x_0)\rho(x_0^{-1}g)$, this is at most a
constant times $\rho(g)^{-a}$.  The constant $\rho(x_0)^a$ is
independent of $g$.
Taking finite linear combinations proves both assertions for every
$m\in S$.  All the local smooth extensions used here may be chosen
after the fixed translate; on the group they are the actual functions
(BST.8), so equality of tangential derivatives is unaffected by
off-group values.
\end{proof}
% END BST3

\subsection{The paired tensor maps at their original scope}

% BEGIN BST4
\paragraph{BST4. Exact standard and dual tensor intertwining.}
\label{claim:BST4}
For every $x_1\in\mathbb C^n$, every $x_0\in G$, $w\in V$, and
every tuple $f=(f_j)_{j=0}^{n-1}\in S^n$, define the original maps
\[
 \begin{aligned}
 Q_f(g)&=\sum_j q_j(g)F_{f_j}(g),
       &q_j(g)&=x_1^*J_\epsilon ge_j,\\
 Q_f^\vee(g)&=\sum_j r_j(g)F_{f_j}(g),
       &r_j(g)&=e_j^{\,t}g^{-1}x_1.
 \end{aligned}
 \tag{BST.10}
\]
In particular $x_1$ is arbitrary; it need not come from a Bessel frame.
The two functions are smooth, and the exact tensor actions on tuples are
\[
 \begin{split}
 (L_{ab}^{+}f)_c&=E_{ab}f_c+\delta_{ca}f_b,\\
 (L_{ab}^{-}f)_c&=E_{ab}f_c-\delta_{cb}f_a.
 \end{split}
 \tag{BST.11}
\]
They satisfy on $G$
\[
                  E_{ab}Q_f=Q_{L_{ab}^{+}f},
 \qquad E_{ab}Q_f^\vee=Q_{L_{ab}^{-}f}^\vee .
 \tag{BST.12}
\]
The actions in (BST.11) are the literal \texttt{Lem\allowbreak K.\allowbreak tens\allowbreak V}
and \texttt{Lem\allowbreak K.\allowbreak tens\allowbreak Vd}.

\begin{proof}
The first multiplier is linear in the matrix entries; the second is
holomorphic on invertible matrices.  Both have the required local
real smoothness.  Differentiating $g\exp(tX)$ and its inverse gives
\[
 E_{ab}q_j=\delta_{jb}q_a,\qquad
 E_{ab}r_j=-\delta_{ja}r_b.
 \tag{BST.13}
\]
These are first obtained along real Lie directions and then extended
complex linearly, so their conventions agree with the original
$E_{ab}$.  By BST3 the coefficient factors are smooth and their
derivatives act on $f_j$.
The ordinary product rule now gives
\[
 E_{ab}Q_f=\sum_jq_jF_{E_{ab}f_j}+q_aF_{f_b},\qquad
 E_{ab}Q_f^\vee=\sum_jr_jF_{E_{ab}f_j}-r_bF_{f_a}.
 \tag{BST.14}
\]
Linearity of $F_m$ in $m$ turns these formulas exactly into
(BST.11)--(BST.12).  All tuples on the right still lie in $S^n$,
since $S$ is a submodule stable under every $E_{ab}$.
\end{proof}
% END BST4

% BEGIN BST5
\paragraph{BST5. The whole original \texttt{TMap\allowbreak Statement} and \texttt{TMap\allowbreak DStatement}.}
\label{claim:BST5}
Under every original datum and hypothesis of those two statements,
for every $x_1,x_0,w,f$ as in BST4 and every derivative word
$\mathbf w$ there is a real constant $C_{\mathbf w}$ such that
\[
 \begin{aligned}
 E_{\mathbf w}Q_f&=Q_{L_{\mathbf w}^{+}f},
 &|E_{\mathbf w}Q_f(g)|&\le C_{\mathbf w}\rho(g)^{1-a},\\
 E_{\mathbf w}Q_f^\vee&=Q_{L_{\mathbf w}^{-}f}^\vee,
 &|E_{\mathbf w}Q_f^\vee(g)|&\le C_{\mathbf w}\rho(g)^{1-a}.
 \end{aligned}
 \tag{BST.15}
\]
The two constants may be enlarged to one as displayed.
Both functions and every derivative word lie in the actual
$L^2(G,\mu_G)$ of the source.  Thus both original
\texttt{Is\allowbreak Smooth\allowbreak L2} conclusions, both exact first-derivative
intertwinings, and all original word estimates hold.

\begin{proof}
Iterate BST4.  Define
$L_\varnothing^\pm f=f$ and
$L_{(a,b)::\mathbf w}^\pm f=L_{ab}^\pm L_{\mathbf w}^\pm f$.
Induction using (BST.11) keeps every component of every such tuple
in $S$ and proves the equalities in (BST.15), including their order.

There is one constant $B_{x_1}$ such that, for all $j$ and $g\in G$,
\[
                    |q_j(g)|,\ |r_j(g)|\le B_{x_1}\rho(g).
 \tag{BST.16}
\]
For $q_j$ this follows directly from Cauchy--Schwarz and
$\|J_\epsilon\|_{\rm op}=1$.  For $r_j$ use
$g^{-1}=J_\epsilon g^*J_\epsilon$, so
$\rho(g^{-1})=\rho(g)$.
Fix a word and apply BST3 to each component of
$L_{\mathbf w}^\pm f$.  Sum the finitely many coefficient bounds
and multiply by (BST.16).  This proves the exact exponent $1-a$,
with no extra factor depending on word length in the exponent.

Finally $\beta=a-1>0$ and $2\beta>2nq_0$ by BSP1.
The square of either bound in (BST.15) is a constant times
$\rho^{-2\beta}$, which is integrable by BSP6 for the actual Haar
normalization.  Smoothness gives measurability, so every derivative
word is genuinely in $L^2$, rather than merely having a formal
\texttt{to\allowbreak L2} value.  Smoothness itself was proved in BST4 and is
preserved by every tangential derivative.  These are precisely all
fields of the source's \texttt{Is\allowbreak Smooth\allowbreak L2}.
The argument retains arbitrary generated-module tuples, arbitrary
free $x_0$, all $x_1$, every lowest-type $w$, and all original
signature and gap assumptions.
\end{proof}
% END BST5

\subsection{Complementary minors over an arbitrary complex algebra}

For this subsection let $\mathcal U$ be any associative unital complex
algebra and let $E_{ij}\in\mathcal U$ obey
\[
 [E_{ij},E_{kl}]=\delta_{jk}E_{il}-\delta_{li}E_{kj}.
 \tag{BST.17}
\]
Put $A_{ij}(t)=E_{ij}+t\delta_{ij}$.  For lists
$R,C:\{0,\ldots,p-1\}\to\{0,\ldots,n-1\}$, define
\[
 \begin{split}
 C_{R,C}(x)
  &=\sum_{\sigma\in S_p}\operatorname{sgn}(\sigma)
       \prod_{j=0}^{p-1}A_{R_{\sigma(j)},C_j}(x-j),\\
 D_{R,C}(x)
  &=\sum_{\sigma\in S_p}\operatorname{sgn}(\sigma)
       \prod_{i=0}^{p-1}A_{R_i,C_{\sigma(i)}}(x-p+1+i).
 \end{split}
 \tag{BST.18}
\]
Products have the displayed left-to-right order.  These are exactly
the source's column and row minors; empty products equal one.
A hat over an index deletes that position and retains the increasing
order of all remaining positions.

% BEGIN BST6
\paragraph{BST6. Alternation and the exact minor commutator.}
\label{claim:BST6}
Both minors in (BST.18) are alternating in both index lists, including
vanishing if either list repeats an index.  Moreover
\[
 [E_{ab},C_{R,C}(x)]
 =\sum_{i:R_i=b}C_{R[i\leftarrow a],C}(x)
  -\sum_{j:C_j=a}C_{R,C[j\leftarrow b]}(x).
 \tag{BST.19}
\]
All assertions hold in every algebra satisfying (BST.17).

\begin{proof}
Row alternation of $C$ and column alternation of $D$ are immediate
by reindexing their defining permutation sums.  To prove the other
two alternations without commuting noncommuting entries, adjoin
exterior variables $e_0,\ldots,e_{n-1}$ commuting with $\mathcal U$.
Set
$\eta_j=\sum_i e_iE_{ij}$ and $\eta_j(t)=\eta_j+te_j$.
Directly from (BST.17),
\[
 \eta_i\eta_j+\eta_j\eta_i=-e_i\eta_j-e_j\eta_i,
 \qquad
 \eta_i(t)\eta_j(t-1)=-\eta_j(t)\eta_i(t-1).
 \tag{BST.20}
\]
For example the first identity follows by writing its left side as
$\sum_{a,b}e_ae_b[E_{ai},E_{bj}]$; the two Kronecker sums give
the two displayed terms.  Expanding the scalar shifts proves the
second identity.  The coefficient of
$e_{R_0}\cdots e_{R_{p-1}}$ in
$\eta_{C_0}(x)\cdots\eta_{C_{p-1}}(x-p+1)$ is
$C_{R,C}(x)$ when the rows are distinct, with the indicated ordering.
Adjacent column interchange negates the whole product by (BST.20).
Together with the already proved row alternation, this gives column
alternation for every row list as well.

For the row version set $\xi_i=\sum_j e_jE_{ij}$.
The same calculation gives
$\xi_i\xi_j+\xi_j\xi_i=e_i\xi_j+e_j\xi_i$, and therefore
\[
       \xi_i(t)\xi_j(t+1)=-\xi_j(t)\xi_i(t+1).
 \tag{BST.21}
\]
Coefficients in
$\xi_{R_0}(x-p+1)\cdots\xi_{R_{p-1}}(x)$ prove row alternation of
$D$.  A repeated list entry gives its own negative under a
transposition; characteristic zero makes the minor zero.

Finally
$[E_{ab},A_{cd}(t)]=
\delta_{bc}A_{ad}(t)-\delta_{da}A_{cb}(t)$.
The scalar correction terms on the right cancel each other.
Apply the Leibniz rule for a commutator to each ordered product
in $C$.  In the first set of terms, replace its unique row position
by $a$ and sum over that position; in the second set replace the
specified column position by $b$.  Summing over permutations is
exactly (BST.19).  No reordering of the product factors occurs.
\end{proof}
% END BST6

% BEGIN BST7
\paragraph{BST7. Equality of the shifted row and column minors.}
\label{claim:BST7}
For every $p,R,C,x$,
\[
                         C_{R,C}(x)=D_{R,C}(x).
 \tag{BST.22}
\]
In particular the first-row expansion of the column minor is
\[
 C_{R,C}(x)=
 \sum_{j=0}^{p}(-1)^j A_{R_0,C_j}(x-p)
           C_{R\widehat{0},C\widehat{j}}(x)
 \tag{BST.23}
\]
when $R,C$ have length $p+1$.

\begin{proof}
Induct on the common list length $N$.  The empty case is one.
Assume the assertion for length $N-1$, and abbreviate
$C_{ij}=C_{R\widehat{i},C\widehat{j}}(x-1)$.
Ordinary first-column expansion and column alternation from BST6
allow any chosen column to be moved to the first position; hence
\[
 N C_{R,C}(x)=
       \sum_{i,j}(-1)^{i+j}A_{R_i,C_j}(x)C_{ij}.
 \tag{BST.24}
\]
Similarly, use the last-row expansion of the row minor, moving each
chosen row to the last position.  Its entry at that position has
parameter $x$.  The induction hypothesis on the remaining minor
gives
\[
 N D_{R,C}(x)=
       \sum_{i,j}(-1)^{i+j}C_{ij}A_{R_i,C_j}(x).
 \tag{BST.25}
\]
These expansions are valid also for repeated indices, by alternation.

We check the entire commutator difference of the two double sums.
Define
\[
 T_{sj}=\delta_{R_s,C_j}(-1)^{s+j+1}
                  C_{R\widehat{s},C\widehat{j}}(x-1).
 \tag{BST.26}
\]
Formula (BST.19) gives exactly
\[
 (-1)^{i+j}[E_{R_i,C_j},C_{ij}]
               =\sum_{s\ne i}T_{sj}-\sum_{t\ne j}T_{it}.
 \tag{BST.27}
\]
To verify its signs, a term in the first commutator sum requires
$R_s=C_j$, $s\ne i$.  Replacing row $R_s$ in the list with $i$
deleted by $R_i$, and putting the result into the order with $s$
deleted, has sign $(-1)^{i+s+1}$.  Multiplication by
$(-1)^{i+j}$ leaves $(-1)^{s+j+1}$, the sign of $T_{sj}$.
The column replacements similarly leave $(-1)^{i+t+1}$ and occur
with the minus sign in (BST.19).  This proves (BST.27).
The scalar part of $A_{R_i,C_j}(x)$ commutes and contributes nothing.

On summing (BST.27) over $i,j$, every $T_{sj}$ occurs $N-1$
times in the first sum and $N-1$ times in the second.  The
difference is zero.  Equations (BST.24)--(BST.25) therefore give
$N(C_{R,C}-D_{R,C})=0$; division by the nonzero complex integer
$N$ completes the induction.  Finally apply ordinary first-row
expansion to $D$ and (BST.22) to its remaining minors.  Its first
parameter is $x-p$, which proves the exact shift in (BST.23).
\end{proof}
% END BST7

% BEGIN BST8
\paragraph{BST8. The full original complementary Capelli expansions.}
\label{claim:BST8}
Let $n=m+1$, $\mathcal C(x)=C_{\mathrm{id},\mathrm{id}}(x)$,
and define
\[
 U_{dc}(x)=(-1)^{d+c}C_{\widehat d,\widehat c}(x-1),
 \qquad
 V_{dc}(x)=(-1)^{d+c}C_{\widehat c,\widehat d}(x).
 \tag{BST.28}
\]
Then for every complex $x$,
\[
 \begin{split}
 \sum_d(E_{db}+x\delta_{db})U_{dc}(x)
                         &=\delta_{bc}\mathcal C(x),\\
 \sum_d(E_{ad}+(x-m)\delta_{ad})V_{dc}(x)
                         &=\delta_{ac}\mathcal C(x).
 \end{split}
 \tag{BST.29}
\]
This is the complete original \texttt{Capelli\allowbreak Expand\allowbreak Statement}.

\begin{proof}
For the first identity take the column minor with full row list and
column list $(b,0,\ldots,\widehat c,\ldots,m)$.
Expand its first column.  The terms are
$(-1)^d(E_{db}+x\delta_{db})
C_{\widehat d,\widehat c}(x-1)$.
Multiply their sum by $(-1)^c$.  If $b=c$, the column list is
the identity list with $c$ moved to the front, so its sign is
$(-1)^c$.  If $b\ne c$, the list repeats $b$ and the minor
vanishes by BST6.  This proves the first identity, including the
order of every factor.

For the second take row list $(a,0,\ldots,\widehat c,\ldots,m)$
and full column list.  Use (BST.23), whose first-row entries
have parameter $x-m$, and again multiply by $(-1)^c$.
Row alternation gives the displayed right side.  Empty complementary
minors when $m=0$ are one, so no positive-rank exception was used.
\end{proof}
% END BST8

% BEGIN BST9
\paragraph{BST9. The whole original \texttt{Eig\allowbreak Statement}.}
\label{claim:BST9}
Let $M$ be any complex module, $E_{ij}\in\operatorname{End}_{\mathbb C}M$
satisfy (BST.17), and $S\subset M$ be stable under every $E_{ij}$.
Let $m\ge0$ and suppose, for one real $\omega$, that
$\mathcal C(\omega)s=0$ for every $s\in S$.
For arbitrary $y,\alpha\in\mathbb C^{m+1}$ and $v\in S$, set
\[
 (n^+)_d=\sum_c y_cU_{dc}(\omega)v,\qquad
 (n^-)_d=\sum_c \alpha_cV_{dc}(\omega)v.
 \tag{BST.30}
\]
These are exactly \texttt{nn\allowbreak Plus} and \texttt{nn\allowbreak Minus}.
Writing
\[
 (\Omega^+f)_b=\sum_d E_{db}f_d,\qquad
 (\Omega^-f)_a=-\sum_d E_{ad}f_d,
 \tag{BST.31}
\]
one has
\[
 n^+\in S^{m+1}\cap\ker(\Omega^++\omega),\qquad
 n^-\in S^{m+1}\cap\ker(\Omega^--(\omega-m)).
 \tag{BST.32}
\]
Thus the eigenvalues are the original $-\omega$ and $\omega-m$.

\begin{proof}
Every entry in a complementary minor is a matrix-unit operator plus
a scalar.  Stability of $S$, followed by finite products and sums,
puts every component of (BST.30) in $S$.
Apply the first identity of (BST.29) to $v$, multiply by $y_c$,
and sum over $c$.  Its right side is zero.  The left side is
$(\Omega^+n^+)_b+\omega(n^+)_b$, giving the first kernel condition.
The second identity, multiplied by $\alpha_c$ and summed, gives
$\sum_dE_{ad}(n^-)_d+(\omega-m)(n^-)_a=0$.
The minus sign in the definition of $\Omega^-$ gives the second
condition exactly.  These proofs work over an arbitrary module;
they assume no Hilbert structure, nonzero vector, or irreducibility.
The complex version of the two eigenvalue identities also holds;
real $\omega$ is retained to match the source's real eigenvalue
parameter in \texttt{Xs}.
\end{proof}
% END BST9

% BEGIN BST10
\paragraph{BST10. The full tensor Casimir identities and both exact shifts.}
\label{claim:BST10}
For every $n,M,E$ as above, the two actions $L^\pm$ of (BST.11)
are representations of $\mathfrak{gl}_n$ and obey
\[
 C_{L^\pm}f=(C_Ef_j)_j+n f+2\Omega^\pm f,
 \qquad C_E=\sum_{a,b}E_{ab}E_{ba}.
 \tag{BST.33}
\]
This is the complete original \texttt{Capelli\allowbreak Tensor\allowbreak Statement}.
In the original kernel application, put
$\theta_s=\lambda_s-(n-1)/2$ and $\omega=-\theta_s$.
BSP3 and BST9 then give the prescribed standard and dual
eigenspace vectors.  On those eigenspaces their tensor Casimirs are
respectively
\[
                 c_{\lambda+e_s},\qquad c_{\lambda-e_s}.
 \tag{BST.34}
\]
The vectors are not asserted nonzero.

\begin{proof}
Write $L^\pm_{ab}$ as the componentwise operator $E_{ab}$ plus
the standard, respectively contragredient, matrix-unit action on
the tuple index.  The two summands commute.  Each satisfies the
matrix-unit bracket relations, so their sum does also.
Expand $\sum_{a,b}L^\pm_{ab}L^\pm_{ba}$.
The componentwise terms give $(C_Ef_j)_j$.
On either the standard or dual index representation,
$\sum_{a,b}e_{ab}e_{ba}=nI$; in the dual action the two minus
signs cancel.  The two cross terms agree after exchanging $a,b$.
For the standard action each has component
$\sum_dE_{db}f_d$; for the dual each has component
$-\sum_dE_{ad}f_d$.  They are precisely the two terms
$2\Omega^\pm$ in (BST.33).

In the application $C_E=c_\lambda$ on $S$.
Indeed it has that scalar on $V$ by the original datum and commutes
with every $E_{ab}$ by the matrix-unit relations, so the identity
extends to the generated module.
The eigenvalues from BST9 are
$\theta_s$ and $-\theta_s-(n-1)$, respectively.
Since $c_\lambda=\|\lambda\|^2-\|\rho\|^2$,
\[
 c_\lambda+n+2\theta_s=c_\lambda+2\lambda_s+1
        =c_{\lambda+e_s},\qquad
 c_\lambda+n+2(-\theta_s-(n-1))
        =c_\lambda-2\lambda_s+1=c_{\lambda-e_s}.
 \tag{BST.35}
\]
This proves the literal shifts, including the contragredient sign.
It does not identify a lowest compact type in an image or quotient.
\end{proof}
% END BST10

\subsection{The exact residual unit-shift problem}

% BEGIN BST11
\paragraph{BST11. Reduction of the whole original \texttt{Step\allowbreak Statement} to four precise suppliers.}
\label{claim:BST11}
Retain the full original \texttt{Eval\allowbreak PStatement},
\texttt{Eval\allowbreak MStatement}, \texttt{Quot\allowbreak PStatement} and
\texttt{Quot\allowbreak MStatement}, at the exact data and quantifiers in
\texttt{Laplace\allowbreak Estimates/\allowbreak Psi\allowbreak Datum\allowbreak Statement}.  Then the reviewed
DC/BNV/BSP results and BST1--BST10 prove the whole original
\texttt{Step\allowbreak Statement}.
Thus these four suppliers, rather than a derivative-decay or
complementary-eigenspace input, are the remaining assumptions in
this analytic reduction.

\begin{proof}
Fix arbitrary original $k,r,\epsilon,u,h$ with the signature and
\texttt{Step\allowbreak Gap} hypotheses.  The original subgroup-integrability
antecedent is supplied by BNV1.
Take a witness to $\texttt{Bessel\allowbreak NZ}(\epsilon,\lambda,u)$ and
one of the two allowed endpoints $\lambda\pm e_s$.
The definition of an allowed step retains its Harish--Chandra and
noncompact-gap hypotheses, as well as the appropriate exclusion
$u_j\ne\lambda_s\pm1/2$ for every $j$.

By BSP1 and BSP2 choose an exact kernel datum for the old
$T,\mu$.  If the old left translate is $x_0\in G_H$, put
$x'_0=Px_0P^{-1}\in G$.  Under the isomorphism $\iota$,
the two old kernel vectors become $v,w\in V$.
On the actual Haar support the coefficient identity is
\[
 \langle w,\sigma((x'_0)^{-1}P h'n'xP^{-1})v\rangle
 =
 \langle v_2,\Psi(P(x_0^{-1}h'n'x)P^{-1})v_1\rangle.
 \tag{BST.36}
\]
Hence its original \texttt{per} integral is the old nonzero
\texttt{aa\allowbreak Fun} integral.  The Cayley matrix, character, compact
representation, Haar measures, and the vectors and translates
are exactly the old ones.

For the positive shift, the retained \texttt{Eval\allowbreak P} supplies $y$
such that the period of the exact tensor map of
$\texttt{nn\allowbreak Plus}(y,-\theta_s,v)$ is nonzero.
BSP3 kills $\mathcal C(-\theta_s)$ on all $S$; BST9 places this
tuple in the original standard $\theta_s$ eigenspace.
Apply the retained \texttt{Quot\allowbreak P} to this tuple and nonzero period.
It supplies the actual spaces $W_f,W'$ with their exact image
relation, shifted lowest type and Casimir, actual compact stability,
and period injectivity.
Every $F\in W_f$ is a tensor map of a tuple in this eigenspace.
BST5 gives its continuity and the required
$C_F\rho^{1-a}$ bound.  Thus every original hypothesis of
BSP10 is supplied, and BSP10 gives
$\texttt{Bessel\allowbreak NZ}(\epsilon,\lambda+e_s,u)$.

For the negative shift use \texttt{Eval\allowbreak M}, the dual tuple
$\texttt{nn\allowbreak Minus}(\alpha,-\theta_s,v)$ and its eigenvalue
$-\theta_s-(n-1)$ from BST9.  These are the exact inputs of
\texttt{Quot\allowbreak M}.  Its resulting functions obey the dual half of
BST5, and BSP10 gives
$\texttt{Bessel\allowbreak NZ}(\epsilon,\lambda-e_s,u)$.
Both arguments keep every original parameter quantifier and allow
the new existential kernel datum to be the one supplied by BSP10.
Together they say exactly that the original Bessel predicate is
step closed, which is \texttt{Step\allowbreak Statement}.
\end{proof}
% END BST11

\paragraph{What this chapter leaves open.}
The remaining evaluation statements require a nonzero period of the
specified complementary-minor tensor image, with the actual
character and Haar measures.  BST9 does not establish this
nonvanishing; its eigenspace vector may be zero.
The remaining quotient statements must construct $W_f,W'$,
their exact shifted lowest compact type, and period injectivity.
The Casimir computation (BST.35) supplies one required scalar, not
that construction.
Consequently \texttt{Step\allowbreak Statement}, BNV23's unconditional
\texttt{Bessel\allowbreak NVGap}, and the subsequent setting remain open at this
stage.  The eventual global-setting join still retains the original
supplied \texttt{Forms0} scope.


\section{Exact evaluation of both complementary-minor Bessel periods}
\label{sec:bessel-period-evaluation-v51}

This chapter proves the whole original \texttt{Eval\allowbreak PStatement} and
\texttt{Eval\allowbreak MStatement}.  It identifies explicit choices of the
standard and dual input vectors and computes their actual periods as
nonzero scalar multiples of the old period.  The character, compact
representation, Haar measures and both free translates remain the
original ones.  The shifted-quotient constructions are still separate.

Fix arbitrary data in those two source statements.  Explicitly,
$n=2r+k+1$, $p=k+r$, $q=r+1$, the original
\texttt{Step\allowbreak Gap} holds for $h$, $\lambda$ is Harish--Chandra with
noncompact gaps at least $h$, and
$(\mathcal H,\sigma,D,E,V,T,\mu,\iota)$ is any original complete
unitary derived datum.  It has the prescribed compact irreducible
lowest type, Casimir, actual compact-group stability,
$\mu=\lambda+\rho$, and the exact $\Psi$ coefficient identity.
We impose no irreducibility on the ambient representation.
Let $S=U(\mathfrak{gl}_n)V$, $G=U(J_\epsilon)$,
$G_H=U(H_{pk})$ and $d=n-1=2r+k$.

The original Bessel datum consists of a negative middle vector $z$,
$T_z=-z^*H_{pk}z>0$, nonzero coefficients $c_a$ for $0\le a<r$,
an arbitrary permitted $\pi$ of compact type $u$, and the original
weak Haar measures $\mu_W,\mu_N$.  Fix its Cayley matrix $P$,
any $x_0\in G$, $w\in V$, $\xi\in V_\pi$, and any $g\in G_H$.
For $m\in S$ put
\[
 \begin{split}
 F_m(y)&=\langle w,\sigma(x_0^{-1}y)m\rangle,\\
 \ell(m)&=\int_{H_z\times N}
       F_m(P h'n'gP^{-1})\,
       \overline{\psi_{z,c}(n')}\,\pi(h')^{-1}\xi\,
                    d\mu_W(h')d\mu_N(n').
 \end{split}
 \tag{BSE.1}
\]
This is the literal \texttt{coef\allowbreak L}/\texttt{per} functional.
BST3 and BSP5 prove its absolute convergence and linearity on $S$,
and justify all differentiations used below.  Those results also
apply to every matrix-unit word in $m$.
For the two tensor maps use BST's $Q_f,Q_f^\vee$ with
\[
 x_*= \begin{cases}e_0,&r>0,\\ z,&r=0,\end{cases}
                  \qquad x_1=P x_* .
 \tag{BSE.2}
\]
Their periods are absolutely convergent by BST5 and BSP5.
The arbitrary $g$ above is the point at which each period is evaluated.

\subsection{The actual period and its framed infinitesimal law}

% BEGIN BSE1
\paragraph{BSE1. Exact period equivariance with the original Haar measures.}
\label{claim:BSE1}
For every original Bessel datum, every complex matrix $P$,
every matrix function $F$, every $\xi\in V_\pi$ and every matrix
$x$, the source's period obeys
\[
 \mathcal P_F(h_0m_0x)
       =\psi(m_0)\pi(h_0)\mathcal P_F(x),
       \qquad h_0\in H_z,\quad m_0\in N.
 \tag{BSE.3}
\]
This holds in particular for every function and derivative used in
(BSE.1), and with no scalar change in either Haar measure.

\begin{proof}
By AMD8 the original supported weak Haar measures are ordinary Haar
measures on their closed subgroups.  The group $H_z$ is compact by
BWP3; $N$ is nilpotent and unimodular by the polynomial structure
proved in BWP1 and BNV11.  In exponential coordinates the adjoint
operators of its Lie algebra are nilpotent and have trace zero,
so the modulus of every inner translation is one.  Thus both
measures are also right invariant.
The action of $H_z$ on $N$ preserves the latter Haar measure:
its modulus is a continuous homomorphism from the compact group
$H_z$ into $\mathbb R_{>0}$ and hence equals one.
These arguments concern the given measures themselves, not a
replacement normalization.

The group definitions give $H_zNH_z^{-1}=N$ and
$\psi(h^{-1}nh)=\psi(n)$.  To check the character equality directly,
the first $r-1$ relevant matrix entries are unchanged because
$h$ fixes every outer flag vector, and the last relevant entry
is the coefficient of $nz$, unchanged because $hz=z$.
The same first-level entries add under multiplication in $N$,
so $\psi$ is its actual unitary character.

In the period at $h_0m_0x$ make the measure-preserving substitution
\[
             h''=h'h_0,\qquad
             n''=h_0^{-1}n'h_0m_0.
 \tag{BSE.4}
\]
Then $h'n'h_0m_0=h''n''$,
$\psi(n')^{-1}=\psi(m_0)\psi(n'')^{-1}$, and
$\pi(h')^{-1}=\pi(h_0)\pi(h'')^{-1}$.
Pull the fixed scalar and fixed unitary operator outside the
integral.  This gives (BSE.3).
For the functions used here every interchange is legitimate by
BST3/BST5 and BSP5.  The source states the covariance for arbitrary
matrix functions as well, with its zero convention for undefined
Bochner integrals: the same measurable invertible substitution and
fixed isometry preserve integrability and measurability, so if these
fail both sides are zero.  Thus its unrestricted convention causes
no change to the identity.
\end{proof}
% END BSE1

% BEGIN BSE2
\paragraph{BSE2. The exact compact Capelli polynomial.}
\label{claim:BSE2}
Let $\tau$ be any original $\texttt{Is\allowbreak HCType\allowbreak U}(k,u)$ representation.
Its complexified derived matrix units $D_{ij}$ satisfy the
matrix-unit relations and, on all of $V_\pi$,
\[
 \mathcal C_D(t)=p_u(t)I,\qquad
 p_u(t)=\prod_{j=0}^{k-1}
               \left(t+u_j-\frac{k-1}{2}\right).
 \tag{BSE.5}
\]
For every skew-Hermitian $A$, the derivative of
$\tau(\exp(tA))$ is its derived action.  The statement includes $k=0$.

\begin{proof}
The original compact-type predicate is a nonzero finite-dimensional
unitary representation, irreducible of the specified highest weight.
A continuous finite-dimensional Lie-group representation is smooth.
One may see this without an additional classification theorem:
convolution by a smooth approximate identity gives smooth orbit
vectors and an operator arbitrarily close to the identity on the
entire finite-dimensional space; that operator is invertible,
so all its vectors are smooth.  Differentiate the representation
law and complexify $\mathfrak u(k)$ to obtain $D$ and its
matrix-unit relations.

The original compact-type predicate supplies a nonzero torus
vector of weight $\nu_j=u_j-(k-1)/2+j$ and says that every weight
is dominated by $\nu$.  Differentiating torus conjugation shows
that $D_{ab}$ sends that vector to weight $\nu+e_a-e_b$.
If $a<b$, this weight is not dominated by $\nu$: the linear
functional $\sum_j(k-1-j)\mu_j$ strictly increases under adding
$e_a-e_b$ and weakly decreases under subtracting positive roots.
Thus every upper-triangular $D_{ab}$ kills the vector, and it is
the required highest vector.
In the original ordered column determinant
for $\mathcal C_D(t)$, apply the rightmost factor to this vector.
Any nonidentity permutation has a largest column not fixed by it;
the earlier fixed larger columns have already acted diagonally, and
that column's row is strictly smaller than the column.  Its
upper-triangular matrix unit kills the highest vector.  Thus only
the identity permutation survives, giving
$\prod_j(t+\nu_j-j)=p_u(t)$.

The operator $\mathcal C_D(t)$ is central.  This follows directly
from BST6's minor commutator with full identical row and column
lists: if $a\ne b$, each replacement creates a repeated index and
vanishes; if $a=b$, the equal row and column terms cancel.
Hence the scalar identity extends from the highest vector to its
generated module.  That module is the entire irreducible compact
type: its finite-dimensional Lie-stable subspace is invariant
under exponentials, and the connected group $U(k)$ is generated
by them.  Irreducibility completes the argument.
If $k=0$, both the determinant and product in (BSE.5) are the
empty product one, so the assertion remains literal.
\end{proof}
% END BSE2

% BEGIN BSE3
\paragraph{BSE3. The whole original framed period law, with explicit character values.}
\label{claim:BSE3}
Choose a frame $g_0$ furnished by the original
$\texttt{Is\allowbreak HCType\allowbreak W}(k,r,z,u,\pi)$.
It is $H_{pk}$-unitary, fixes the outer flag vectors, and sends
$\sqrt{T_z}e_{r+k}$ to $z$.  Put
\[
 Q=g_0^{-1}gP^{-1},\qquad
 E^Q_{ij}=\sum_{a,b}(Q^{-1})_{ai}Q_{jb}E_{ab}.
 \tag{BSE.6}
\]
Let $\tau(A)=\pi(g_0\operatorname{blockW}(A)g_0^{-1})$ on
$U(k)$, and let $D_{ij}$ be its derived units from BSE2.
For the levels
\[
 b(i)=
 \begin{cases}
 i,&i<r,\\
 r,&r\le i\le r+k,\\
 i-k,&i>r+k,
 \end{cases}
 \qquad W=\{r,\ldots,r+k-1\},
 \tag{BSE.7}
\]
there are scalars $\chi_{ij}$ such that, on every $m\in S$,
\[
 \ell(E^Q_{ij}m)=\chi_{ij}\ell(m)\quad(b(i)<b(j)),\qquad
 \ell(E^Q_{r+i,r+j}m)=D_{ij}\ell(m)\quad(i,j<k).
 \tag{BSE.8}
\]
The scalars vanish except on the consecutive edges of
\[
                   0,1,\ldots,r-1,r+k,r+k+1,\ldots,d.
 \tag{BSE.9}
\]
For $r>0$ their product on these edges is exactly
\[
              \kappa=-\frac{T_z}{4^r}
                               \prod_{a=0}^{r-1}|c_a|^2\ne0.
 \tag{BSE.10}
\]
Together with (BSE.5), these are every conclusion of the original
\texttt{Frame\allowbreak Law\allowbreak Statement}.

\begin{proof}
The frame and compact representation are the actual existential
witnesses in the given compact-type predicate.
For a matrix $Z$ write $E(Z)=\sum Z_{ab}E_{ab}$.
The operators $E^Q$ are $E(Q^{-1}e_{ij}Q)$ and preserve $S$;
conjugation of matrices and the original bracket identity show that
they satisfy the same matrix-unit relations.
For a real $H_{pk}$-Lie direction $Y$, the curve through the period
point is $g_0\exp(tY)g_0^{-1}g$.
Its scalar coefficient in the integrand is
\[
 F_m\bigl(P h'n'gP^{-1}\exp(tQ^{-1}YQ)\bigr).
 \tag{BSE.11}
\]
Here $Q^{-1}YQ$ is in the real Lie algebra of $G$, since
$Q^*H_{pk}Q=J_\epsilon$.
BST3 and BSP5 differentiate under the actual integral and identify
its derivative as $\ell(E(Q^{-1}YQ)m)$.

For the nilpotent directions define
\[
 \begin{split}
 \mathcal L(Y)&=\sum_{a=0}^{r-1}c_a\gamma_aY_{a,\operatorname{next}(a)},\\
 \gamma_a&=\begin{cases}1,&a+1<r,\\\sqrt{T_z},&a=r-1,\end{cases}
 \quad
 \operatorname{next}(a)=
       \begin{cases}a+1,&a+1<r,\\r+k,&a=r-1.\end{cases}
 \end{split}
 \tag{BSE.12}
\]
The Lie algebra of $N$ consists of the $H_{pk}$-skew-Hermitian
matrices strictly above the levels (BSE.7), by BWP1/BNV11.
Their consecutive-level exponential entries are their linear
entries, since every higher power raises level by at least two.
The defining frame fixes the outer coordinates and sends
$\sqrt{T_z}e_{r+k}$ to $z$.  Consequently the character of the
framed one-parameter subgroup is exactly
\[
 \psi\bigl(g_0\exp(tY)g_0^{-1}\bigr)
                         =\exp\bigl(it\,\operatorname{Re}\mathcal L(Y)\bigr).
 \tag{BSE.13}
\]
Differentiating BSE1 along this subgroup and comparing with
(BSE.11) gives
$\ell(E(Q^{-1}YQ)m)=i\operatorname{Re}\mathcal L(Y)\ell(m)$.
For the compact block, BSE1 similarly gives
$\ell(E(Q^{-1}\operatorname{emb}(A)Q)m)=d\tau(A)\ell(m)$
for skew-Hermitian $A$.  Complexification yields the second
identity of (BSE.8).

We make the nilpotent complexification and every scalar explicit.
Let $Z^\sharp=H_{pk}Z^*H_{pk}$.  Strictly increasing levels are
preserved by this involution.  Decompose
$Z=(Z-Z^\sharp)/2+i[-i(Z+Z^\sharp)/2]$ into two real
nilpotent Lie directions.  The complexified character is
\[
 \chi(Z)=\frac i2
       \left(\mathcal L(Z)-\overline{\mathcal L(Z^\sharp)}\right).
 \tag{BSE.14}
\]
Thus $\chi_{ij}=\chi(e_{ij})$ proves the first identity of (BSE.8).
The form $H_{pk}$ pairs $e_i$ with $e_{d-i}$ outside the middle
block, is $+1$ on $W$, and is $-1$ on $e_{r+k}$.
If $\iota$ fixes middle indices and interchanges the paired outer
indices, and $s_i=-1$ only at $i=r+k$, then
$e_{ij}^\sharp=s_is_j e_{\iota(j),\iota(i)}$.
Substituting in (BSE.12)--(BSE.14) shows that no edge outside
(BSE.9) contributes.  For $a<r-1$ the paired values are
$ic_a/2$ and $-i\overline{c_a}/2$.
For the last left edge $(r-1,r+k)$ and its right partner
$(r+k,r+k+1)$ they are
$i\sqrt{T_z}c_{r-1}/2$ and
$i\sqrt{T_z}\overline{c_{r-1}}/2$.
Their product is (BSE.10).
This also proves the source's nonzero \texttt{kappa} condition
with all complex $c_a\ne0$, without choosing their phases.
For $r=0$ there are no strictly increasing levels, so the nilpotent
conditions are empty.  BSE2 supplies (BSE.5) in all cases.
\end{proof}
% END BSE3

\subsection{Restriction and change of frame for the actual minors}

% BEGIN BSE4
\paragraph{BSE4. The whole Capelli--Whittaker restriction, including rank zero.}
\label{claim:BSE4}
Let $M,R$ be arbitrary complex modules, $E$ a matrix-unit action
on $M$ indexed by $0,\ldots,k+2r$, and
$l:M\to R$ a linear map.  Suppose
$lE_{ij}=\chi_{ij}l$ whenever $b(i)<b(j)$, and on these
strictly increasing-level pairs suppose $\chi_{ij}=0$ off the
edges (BSE.9).  No condition is imposed on the remaining values
of $\chi$.  Suppose also
$lE_{r+i,r+j}=D_{ij}l$ for $i,j<k$.
For $r>0$ and every $t\in\mathbb C$,
\[
 l\,C^E_{\widehat d,\widehat0}(t)
       =(-1)^k\kappa\,\mathcal C_D(t-r)\,l,
       \qquad \kappa=\prod_{\text{edges in (BSE.9)}}\chi_{ij}.
 \tag{BSE.15}
\]
This is the whole original \texttt{Capelli\allowbreak BStatement}; no nonzero
assumption on $\kappa$ is needed for the identity.
For $r=0$ one has, under the compact-block intertwining alone,
\[
                   l\,C^E_{\widehat k,\widehat k}(t)
                                =\mathcal C_D(t)\,l.
 \tag{BSE.16}
\]
The latter is the entire original \texttt{Whittaker\allowbreak Zero\allowbreak Statement}.

\begin{proof}
Use the exact shifted minors and expansions proved in BST6--BST8.
First take the minor with rows $0,\ldots,d-1$ and columns
$1,\ldots,d$, and expand at its last column, moving that column
to the front before the ordinary column expansion.
Every row has smaller level than the last column, so applying $l$
replaces each first operator by its character scalar.
All terms except the edge $(d-1,d)$ vanish.
Its row and column positions are both the last position, so its
sign is $+1$.  Removing them lowers the minor parameter by one.
Repeat for the $r$ successive right edges.  The result is their
product times the minor with rows $0,\ldots,r+k-1$, columns
$1,\ldots,r+k$, and parameter $t-r$.

Next expand at the first row using BST7's row expansion.
Until row $r-1$ is reached, only its next left edge survives.
It occupies the first column, so its sign is $+1$.
A first-row expansion leaves the remaining minor's parameter
unchanged.  After these $r-1$ removals the row list is
$r-1,r,\ldots,r+k-1$ and the column list is
$r,r+1,\ldots,r+k$, still at parameter $t-r$.
The only allowed edge from its first row is $(r-1,r+k)$,
which now occupies column position $k$.  Its sign is $(-1)^k$.
Removing it leaves the identical row and column lists
$r,\ldots,r+k-1$.

All discarded entries were off-diagonal, so their scalar shifts
were zero; hence the preceding substitutions did not alter any
Capelli parameter.  The collected scalar is exactly
$(-1)^k\kappa$.  On the remaining block,
$lE_{r+i,r+j}=D_{ij}l$ transfers every ordered factor through $l$,
and hence transfers its full determinant to
$\mathcal C_D(t-r)l$.  This proves (BSE.15), including $k=0$,
when the last determinant is the empty product.
For $r=0$ the minor in (BSE.16) is already the compact block;
the same factor-by-factor substitution proves it directly.
That argument does not even require bracket relations for the
rank-zero restriction identity.
\end{proof}
% END BSE4

% BEGIN BSE5
\paragraph{BSE5. Exact change of frame for every complementary minor.}
\label{claim:BSE5}
For an arbitrary associative complex algebra, elements $E_{ij}$
satisfying
$[E_{ab},E_{cd}]=\delta_{bc}E_{ad}-\delta_{da}E_{cb}$,
and an invertible complex matrix $Q$, define $E^Q$ by
(BSE.6).  Then for all $b,c$ and all $t$,
\[
 (-1)^{b+c}C^{E^Q}_{\widehat b,\widehat c}(t)
 =\sum_{a,j}Q_{ba}(Q^{-1})_{jc}(-1)^{a+j}
                                      C^E_{\widehat a,\widehat j}(t).
 \tag{BSE.17}
\]
In particular, with the literal $\texttt{nn\allowbreak Plus}$ and
$\texttt{nn\allowbreak Minus}$ of BST9,
\[
 \begin{split}
 \sum_a Q_{ba}\,
       \texttt{nn\allowbreak Plus}_E(Q^{-1}e_c,\omega,v)_a
      &=(-1)^{b+c}C^{E^Q}_{\widehat b,\widehat c}(\omega-1)v,\\
 \sum_j(Q^{-1})_{jc}\,
       \texttt{nn\allowbreak Minus}_E(e_b^tQ,\omega,v)_j
      &=(-1)^{b+c}C^{E^Q}_{\widehat b,\widehat c}(\omega)v.
 \end{split}
 \tag{BSE.18}
\]
The identities retain the source's ordered factors and complement
signs.
This is the whole original \texttt{Capelli\allowbreak Cov\allowbreak Statement}:
its two square scalar matrices with product one are mutual
inverses, so take them to be $Q^{-1}$ and $Q$.

\begin{proof}
The entry relation, including its scalar shift, is
\[
 E^Q_{ij}+t\delta_{ij}
    =\sum_{a,b}(Q^{-1})_{ai}Q_{jb}
                                      (E_{ab}+t\delta_{ab}).
 \tag{BSE.19}
\]
The scalar part is correct because $QQ^{-1}=I$.
Expanding a shifted column minor in these entries and grouping
row choices gives
\[
 C^{E^Q}_{R,C}(t)
 =\sum_{\substack{I\text{ increasing}\\ |I|=p}}
   \ \sum_{j_0,\ldots,j_{p-1}}
       \det(Q^{-1}[I,R])
       \prod_{\ell=0}^{p-1}Q_{C_\ell,j_\ell}\,
                             C^E_{I,(j_0,\ldots,j_{p-1})}(t).
 \tag{BSE.20}
\]
Here is why this expansion is valid despite noncommutativity.
For a fixed ordered row choice, the original alternating
permutation sum gives the ordinary determinant of its scalar
$Q^{-1}$ entries.  Repeated row choices have determinant zero.
For distinct choices, group them by their increasing list $I$;
their permutation signs give exactly the row-alternating sum
defining $C^E_{I,(j_0,\ldots,j_{p-1})}$.
Only scalar factors were moved past operator factors.
Thus the operator products still have the original order and
the original descending shifts.

By the column alternation proved in BST6, repeated $j$'s now
vanish.  Grouping their permutations gives a second ordinary
determinant and yields the precise Cauchy--Binet formula
\[
 C^{E^Q}_{R,C}(t)
    =\sum_{\substack{I,J\text{ increasing}\\|I|=|J|=p}}
       \det(Q^{-1}[I,R])\,\det(Q[C,J])\,C^E_{I,J}(t).
 \tag{BSE.21}
\]
Apply this with $p=n-1$, $R=\widehat b$, $C=\widehat c$.
The ordinary cofactor identities give
\[
 \det(Q^{-1}[\widehat a,\widehat b])
      =(-1)^{a+b}\det(Q)^{-1}Q_{ba},\qquad
 \det(Q[\widehat c,\widehat j])
      =(-1)^{c+j}\det(Q)(Q^{-1})_{jc}.
 \tag{BSE.22}
\]
Substitute them into (BSE.21); the determinants cancel and give
(BSE.17).  Expanding the definitions of $\texttt{nn\allowbreak Plus}$ and
$\texttt{nn\allowbreak Minus}$ then gives (BSE.18), at parameters
$\omega-1$ and $\omega$ respectively.
No commutative adjugate identity was applied to the operator entries.
\end{proof}
% END BSE5

% BEGIN BSE6
\paragraph{BSE6. Both tensor periods in the same frame.}
\label{claim:BSE6}
At the fixed original period point $g$, put
\[
 j_*=\begin{cases}0,&r>0,\\d,&r=0,\end{cases}
 \quad
 b_+=\begin{cases}1,&r>0,\\-\sqrt{T_z},&r=0,\end{cases}
 \quad
 b_-=\begin{cases}1,&r>0,\\\sqrt{T_z},&r=0.\end{cases}
 \tag{BSE.23}
\]
For every tuple $f\in S^n$ the actual standard and dual periods are
\[
 \mathcal P_{Q_f}(g)=b_+\sum_jQ_{dj}\ell(f_j),\qquad
 \mathcal P_{Q_f^\vee}(g)=b_-\sum_j(Q^{-1})_{jj_*}\ell(f_j).
 \tag{BSE.24}
\]

\begin{proof}
Every element of $H_zN$ fixes $x_*$ from (BSE.2), by BSP4.
Because it is also $H_{pk}$-unitary, it fixes the row $x_*^*H_{pk}$.
Consequently the scalar multipliers in the standard and dual
tensor maps are constant along the left $H_zN$ factor of the
period.  Pulling them out and using the finite linear sum yields
\[
 \begin{split}
 \mathcal P_{Q_f}(g)
     &=\sum_j(x_*^*H_{pk}gP^{-1})_j\,\ell(f_j),\\
 \mathcal P_{Q_f^\vee}(g)
     &=\sum_j(Pg^{-1}x_*)_j\,\ell(f_j).
 \end{split}
 \tag{BSE.25}
\]
The individual summands are integrable by BST5 and BSP5, so this
linearity uses the actual integrals.

Now $gP^{-1}=g_0Q$.
If $r>0$, $x_*=g_0e_0$,
$e_0^*H_{pk}=e_d^t$, and $Pg^{-1}g_0=Q^{-1}$.
Thus the two multiplier lists are row $d$ of $Q$ and column
$0$ of $Q^{-1}$.
If $r=0$, $x_*=\sqrt{T_z}g_0e_d$ and
$e_d^*H_{pk}=-e_d^t$.  The lists are
$-\sqrt{T_z}$ times row $d$ of $Q$ and
$\sqrt{T_z}$ times column $d$ of $Q^{-1}$.
These are exactly (BSE.23)--(BSE.24).
\end{proof}
% END BSE6

\subsection{The two nonzero evaluated periods}

% BEGIN BSE7
\paragraph{BSE7. The whole original \texttt{Eval\allowbreak PStatement}, with an explicit witness.}
\label{claim:BSE7}
For every original datum fixed above, every $v\in V$, and every
index $s$, suppose $\ell(v)\ne0$ and
$u_j\ne\lambda_s+1/2$ for every $j<k$.
Then the exact original standard tensor period is nonzero with
\[
       y=Q^{-1}e_{j_*},\qquad
       \omega=-\lambda_s+\frac d2 .
 \tag{BSE.26}
\]
More precisely it equals
\[
 \mathcal P_{Q_{\texttt{nn\allowbreak Plus}(y,\omega,v)}}(g)
   =
 \begin{cases}
 \displaystyle \kappa\prod_{j<k}(u_j-\lambda_s-1/2)\,\ell(v),
                                                   &r>0,\\[2mm]
 \displaystyle -\sqrt{T_z}\prod_{j<k}(u_j-\lambda_s-1/2)\,\ell(v),
                                                   &r=0 .
 \end{cases}
 \tag{BSE.27}
\]
This proves the original statement for arbitrary original
$x_0,w,v,\xi,g$, not just for a special Bessel character or frame.

\begin{proof}
Every complementary-minor component of the tuple lies in $S$,
since it is a finite polynomial in its stable matrix-unit
operators applied to $v$.  Thus BSE6 applies.
Its first formula and the first identity of (BSE.18), with
$b=d,c=j_*$, give
\[
 \mathcal P_{Q_{\texttt{nn\allowbreak Plus}}}(g)
       =b_+(-1)^{d+j_*}
                 \ell\!\left(C^{E^Q}_{\widehat d,\widehat{j_*}}
                                           (\omega-1)v\right).
 \tag{BSE.28}
\]
The restriction of $E^Q$ to $S$ satisfies every hypothesis of
BSE4 by the actual functional law BSE3.
If $r>0$, $j_*=0$, so the right side of (BSE.28) is
\[
       (-1)^{d+k}\kappa\,p_u(\omega-1-r)\ell(v)
                              =\kappa\,p_u(\omega-1-r)\ell(v).
 \tag{BSE.29}
\]
The sign is one because $d+k=2r+2k$ is even.
If $r=0$, $j_*=d=k$, the complement is the compact principal
minor and $(-1)^{d+j_*}=1$.
BSE4 instead gives $-\sqrt{T_z}p_u(\omega-1)\ell(v)$.
In either case
\[
 \omega-1-r+u_j-\frac{k-1}{2}
                                  =u_j-\lambda_s-\frac12,
 \tag{BSE.30}
\]
where $r=0$ is included.  This proves the exact formula (BSE.27).
The scalar $\kappa$ is nonzero by (BSE.10), and $\sqrt{T_z}>0$.
Each displayed linear factor is nonzero by the original exclusion.
The empty product when $k=0$ is one.
Since $\ell(v)\ne0$, the period is nonzero.
All factors came from the prescribed integral; no abstract
branching functional was substituted.
\end{proof}
% END BSE7

% BEGIN BSE8
\paragraph{BSE8. The whole original \texttt{Eval\allowbreak MStatement}, with an explicit witness.}
\label{claim:BSE8}
For the same arbitrary original data, suppose $\ell(v)\ne0$ and
$u_j\ne\lambda_s-1/2$ for every $j<k$.
Take the row
\[
                    \alpha=e_d^tQ,\qquad
                    \omega=-\lambda_s+\frac d2 .
 \tag{BSE.31}
\]
Then the exact original dual tensor period is nonzero, and equals
\[
 \mathcal P_{Q^\vee_{\texttt{nn\allowbreak Minus}(\alpha,\omega,v)}}(g)
   =
 \begin{cases}
 \displaystyle \kappa\prod_{j<k}(u_j-\lambda_s+1/2)\,\ell(v),
                                                    &r>0,\\[2mm]
 \displaystyle \sqrt{T_z}\prod_{j<k}(u_j-\lambda_s+1/2)\,\ell(v),
                                                    &r=0 .
 \end{cases}
 \tag{BSE.32}
\]

\begin{proof}
As before the tuple lies in $S^n$.
Use the dual formula of BSE6 and the second identity of (BSE.18):
\[
 \mathcal P_{Q^\vee_{\texttt{nn\allowbreak Minus}}}(g)
       =b_-(-1)^{d+j_*}
               \ell\!\left(C^{E^Q}_{\widehat d,\widehat{j_*}}
                                                   (\omega)v\right).
 \tag{BSE.33}
\]
There is no minus-one shift in this complementary minor.
For $r>0$, BSE4 gives
$\kappa p_u(\omega-r)\ell(v)$ after the same sign cancellation.
For $r=0$ it gives
$\sqrt{T_z}p_u(\omega)\ell(v)$.
Now
\[
 \omega-r+u_j-\frac{k-1}{2}
                                  =u_j-\lambda_s+\frac12.
 \tag{BSE.34}
\]
This proves (BSE.32).  Its scalar is nonzero by the exact
negative-step exclusion, (BSE.10), and $T_z>0$.
The coefficient slot, character and inverse compact action are
unchanged from (BSE.1); this is precisely the original
\texttt{Eval\allowbreak MStatement}.
\end{proof}
% END BSE8

% BEGIN BSE9
\paragraph{BSE9. The exact remaining quotient input to the full Bessel theorem.}
\label{claim:BSE9}
Retain the complete original \texttt{Quot\allowbreak PStatement} and
\texttt{Quot\allowbreak MStatement}.  Then the reviewed DC/BNV/BSP results,
BST1--BST11 and BSE1--BSE8 imply the whole original
\texttt{Step\allowbreak Statement} and therefore the whole original
\texttt{Bessel\allowbreak NVGap}.
On the inherited supplied \texttt{Forms0} scope, they also imply
the original \texttt{Concrete\allowbreak Setting\allowbreak Exists}.

\begin{proof}
BSE7 and BSE8 supply the two evaluation premises retained in BST11,
with all their original quantifiers, actual periods and exclusions.
Thus only its two quotient premises remain, and the assumed
\texttt{Quot\allowbreak P}/\texttt{Quot\allowbreak M} give both directions of
\texttt{Step\allowbreak Statement} by BST11.
BNV23 then supplies the arbitrary-data endomorphism-valued Bessel
nonvanishing theorem and the inherited BWP24 global-setting join.
The latter still applies on precisely the supplied
\texttt{Forms0} scope.
No new ambient form supplier or orbital/fiber comparison is used.
\end{proof}
% END BSE9

\paragraph{Remaining construction.}
The whole \texttt{Quot\allowbreak P} and \texttt{Quot\allowbreak M} statements are not
proved here.  They must still construct the literal spaces
$W_f,W'$, the exact image relation, the shifted lowest compact type
and Casimir, actual compact-group stability, and injectivity of
the period on the chosen quotient.
The now proved nonzero evaluations and scalar Casimir identities
supply their source premises; they do not by themselves construct
these spaces.  BSE9 is therefore an explicitly conditional
implication, and no whole selected global obligation is declared
closed by this chapter.


\section{Compatible representation families and canonical denominator data}
\label{sec:denominator-family-v051}

\status{DFD1--DFD9 construct the actual continuous integral complements,
the original \texttt{Pt\allowbreak Fam}, and the whole original
\texttt{exists\_\allowbreak hub\allowbreak Data} output. DFD10--DFD14 supply the good-row
selection, the exact twist-killing implication, the finite-coefficient
reduction of the retained integral crystalline theorem, and uniform
parameter and small-prime bounds. DFD15 assembles the original family
setup conditional on two explicitly stated remaining suppliers: the
coherent Wach comparison and the Selmer signature comparison. The full
\texttt{exists\_\allowbreak setup\_\allowbreak large} producer and \texttt{Concl\allowbreak F} remain open in
the human reconstruction.}

Fix the original datum $D:\texttt{D7}$. Its host is the same CM field
$E$, its absolute Galois group is $G_E$, and its rank parameters are
$N=n_V=r+t$. At a location $i$, write $O_i$ for its coefficient DVR,
$K_i=\operatorname{Frac}(O_i)$, $\ell_i=D.\mathrm{nd}.\ell(i)$,
$X_i$ for the point set and $Z_i\subseteq X_i$ for the specified
automorphic points. The characters $\chi_{i,c}:G_E\to O_i^\times$,
for $1\le c\le t$, and both coefficient embeddings are those of $D$.
An integral representation will mean a homomorphism into a general
linear group whose matrix entries are continuous. Its rational
semisimplicity is always stated separately.

The original \texttt{Cheb} premise is retained wherever used. It
asserts the degree-one, unramified Frobenius approximation with finite
avoidance in every finite normal quotient of $G_E$. The present
argument does not replace that premise by a density statement for an
unspecified set of places. Likewise \texttt{adm2},
\texttt{Conclusion41}, \texttt{Facts\allowbreak F}, and the local comparison data
below retain their original source meanings.

\subsection{From trace identities to actual integral complements}

% BEGIN DFD1
\paragraph{DFD1: a character summand descends to the original field.}
\label{claim:DFD1}
Let $K\hookrightarrow L$ be an embedding of fields with $L$ of
characteristic zero, $G$ an arbitrary group, and $V$ a finite-dimensional
semisimple $K[G]$-module. Suppose $T:G\to L$ is the trace of an actual
finite-dimensional $L$-representation and
\[
 \operatorname{tr}(g\mid V)_L=T(g)+\sum_{j=1}^{t}\chi_j(g),
 \qquad \chi_j:G\longrightarrow K^\times.                 \tag{DFD.1}
\]
Then there are independent vectors $v_1,\ldots,v_t\in V$ with
$gv_j=\chi_j(g)v_j$, and a $G$-stable complement $W$ such that
\[
 V=W\oplus\bigoplus_{j=1}^{t}Kv_j.                       \tag{DFD.2}
\]
The characters may repeat; no continuity of $K\hookrightarrow L$ is
required for this algebraic assertion.

Here is the trace argument, including the descent. Extend scalars to
an algebraic closure of $L$. Replace each representation by its
semisimplification when computing its trace: a matrix preserving a
finite composition series is block triangular, so its trace is the
sum of its traces on the composition factors. For the finite direct
sum of all the resulting representations, the image of the group
algebra is a finite-dimensional semisimple algebra. Indeed it acts
faithfully on a semisimple module, so its radical acts by zero and
is zero. Over the algebraically closed field it is a product of
matrix algebras. The trace on the central idempotent of one factor
is the dimension of its simple module times its multiplicity.
The equality (DFD.1), first extended linearly from group elements to
the group algebra and then evaluated on a lift of each idempotent,
therefore identifies the multiplicities. Characteristic zero makes
the positive dimension invertible. Every occurrence of a $\chi_j$
on the right consequently occurs on the left.

Scalar extension preserves semisimplicity here: the image algebra
on $V$ is semisimple; its simple factors have separable centers over
$K$, since $K$ has characteristic zero, and their scalar extensions
are semisimple. In particular the occurrence just found supplies
a nonzero $\chi_1$-eigenvector in the scalar extension of $V$.
The simultaneous equations
\[
 (\rho(g)-\chi_1(g))v=0\quad(g\in G)
\]
have coefficients in $K$. Their row space is finite-dimensional,
so finitely many rows already span all the rows. The rank of this
finite matrix is unchanged by field extension. Its kernel is
therefore nonzero over $K$ as well. Choose $v_1$ there.

Semisimplicity gives a stable complement to $Kv_1$. In a basis
adapted to this decomposition the trace is
$\chi_1+\operatorname{tr}(W_1)$. Subtract $\chi_1$ from (DFD.1)
and repeat on $W_1$. At the $j$th stage the new vector belongs to
the complement of the previously chosen lines, so independence is
automatic, including for repeated characters. After $t$ stages
the remaining stable summand is the required $W$. This proves the
full eigenvector and rational splitting assertions used by the
source's \texttt{exists\_\allowbreak eigenvectors} and \texttt{exists\_\allowbreak complement}.
% END DFD1

% BEGIN DFD2
\paragraph{DFD2: a saturated lattice gives a continuous integral complement.}
\label{claim:DFD2}
Let $O$ be a topological principal ideal domain with fraction field
$K$, and let $R:G\to\operatorname{GL}_N(O)$ have continuous matrix
entries. Suppose $R_K$ is semisimple and its trace has the form
(DFD.1), with $N=a+t$ and the characters valued in $O^\times$.
There is a continuous $\rho:G\to\operatorname{GL}_a(O)$ and an
intertwining $K$-linear isomorphism
\[
 K^N\simeq K^a\oplus K^t,
 \qquad R_K\simeq\rho_K\oplus
      \operatorname{diag}(\chi_1,\ldots,\chi_t),           \tag{DFD.3}
\]
whose block representation is semisimple.

Take the stable summand $W\subset K^N$ from DFD1 and set
$\Lambda=W\cap O^N$. It is finitely generated because $O$ is
Noetherian. It is saturated: if $0\ne c\in O$, $v\in O^N$ and
$cv\in\Lambda$, then $v\in W$ because $W$ is a $K$-space, hence
$v\in\Lambda$. Thus $O^N/\Lambda$ is torsion-free and finitely
generated, and so free over the PID. The exact sequence splits
as $O$-modules. Choose a basis of $\Lambda$ and an $O$-linear
retraction $q:O^N\to O^a$ of its inclusion
$b:O^a\to O^N$. Its rank is $a$: clear the finitely many
denominators of a $K$-basis of $W$ to see that $K\Lambda=W$.

Both $R(g)$ and $R(g^{-1})$ preserve $\Lambda$. Define
\[
             \rho(g)=qR(g)b.                             \tag{DFD.4}
\]
On $\Lambda$, $bq$ is the identity. This proves
$\rho(gh)=\rho(g)\rho(h)$ and
$\rho(g^{-1})=\rho(g)^{-1}$. Every matrix entry of (DFD.4) is a
finite $O$-linear combination of entries of $R(g)$, so it is
continuous. This argument uses continuity only inside $O$;
there is no assumption that a chosen rational basis change is
integral or continuous in an unspecified coefficient topology.

The map $b_K$ identifies $K^a$ with $W$. Adjoin the character
vectors of DFD1 to obtain (DFD.3). Semisimplicity transfers along
this isomorphism from $R_K$. These are exactly the three fields
\texttt{rho}, \texttt{ss}, and \texttt{cont} of the original
\texttt{Is\allowbreak Complement}. The original lattice $O^N$ need not be
the direct sum of the character lattices and $\Lambda$; the
asserted block decomposition is rational, as in the source.
% END DFD2

% BEGIN DFD3
\paragraph{DFD3: equality on the permitted Frobenius elements is enough.}
\label{claim:DFD3}
Assume the original \texttt{Cheb}. If two continuous maps
$f,f':G_E\to Y$, with $Y$ Hausdorff, agree on every Frobenius lift
at every split place outside a finite exceptional set, then $f=f'$.
Here ``split'' is the literal host condition
$\texttt{Is\allowbreak Split\allowbreak Pl}$, not a new density convention.

First a place of $E$ unramified over $\mathbb Q$ and of residue
degree one is an allowed split place. The host extension
$E/\mathbb Q$ is Galois. Its decomposition group at that place
has order $ef=1$, so it is trivial. If complex conjugation fixed
the place, it would belong to this group. But the host's
conjugation sends its nonzero square root of $-d$ to its negative,
so it is not the identity. This proves the required split
condition.

If $f(g)\ne f'(g)$, continuity and the Hausdorff property give an
open neighborhood $U$ of $g$ throughout which they differ.
The Krull topology supplies an open normal subgroup $H$ with
$Hg\subseteq U$. Apply \texttt{Cheb} to $H$, $g$, and the finite
exceptional set. It produces an unramified degree-one place and
a Frobenius lift in $Hg$. The preceding paragraph makes that
place split, so equality there contradicts its membership in
$U$. This proves the source's full
\texttt{AE.\allowbreak eq\_\allowbreak of\_\allowbreak continuous} assertion.

For the applications below enlarge the finite exceptional set
by the places over $\ell_i$, the prescribed sets $S_i,Q_i$, and
the finitely many ramification sets of the automorphic parts.
There are finitely many places over $\ell_i$ because they divide
the nonzero ideal $(\ell_i)$ in $\mathcal O_E$. Thus all such
enlargements are legitimate. Equality is extended in
$\overline{\mathbb Q}_{\ell_i}$, with its given Hausdorff
topology. The complex embedding is used only to check the
Frobenius values and is then canceled by injectivity; it is
never assumed continuous.
% END DFD3

% BEGIN DFD4
\paragraph{DFD4: the cuspidal traces supplied by the original admissibility.}
\label{claim:DFD4}
The original $\texttt{adm2}(D)$ implies
$\texttt{Galois\allowbreak Reps\allowbreak Cusp}(D)$ at every location, with the
original embeddings and normalizations. In particular, for a
regular algebraic conjugate-self-dual cuspidal part $\pi$ of
rank $a\ge1$ there is a continuous actual character $T_\pi$
whose cofinite split Frobenius values satisfy
\[
 \iota_{\ell,\mathbb C}T_\pi(\operatorname{Frob}_w)
       =s_\pi(w)q_w^{(1-a)/2},                            \tag{DFD.5}
\]
where $s_\pi(w)$ is the sum of its Satake parameters and
$q_w=\mathrm Nw$.

The first field of \texttt{adm2} supplies the continuous
semisimple Galois representation with its local
Weil--Deligne compatibility away from $\ell_i$. Remove the
finitely many places over $\ell_i$ and the cofinite automorphic
ramification exceptions. At any remaining place the supplied
compatibility identifies the Frobenius characteristic
polynomial with the normalized unramified Satake polynomial.
The coefficient of degree $a-1$ in
$\prod_j(X-\alpha_j)$ is $-\sum_j\alpha_j$, while that of a
matrix characteristic polynomial is minus its trace. Comparing
these coefficients gives (DFD.5). The trace is continuous since
it is a finite sum of diagonal matrix entries; it is an actual
character because it is the trace of this same representation.
Continuity and semisimplicity are unchanged. This proves all
fields of \texttt{Galois\allowbreak Reps\allowbreak Cusp}, and also the exact
\texttt{has\allowbreak Trace\_\allowbreak cusp} consequence. It assumes the Galois
representation already included in \texttt{adm2}; it does not
assert a new automorphy theorem.
% END DFD4

% BEGIN DFD5
\paragraph{DFD5: integral Tate twists in both parity cases.}
\label{claim:DFD5}
Assume \texttt{Galois\allowbreak Reps\allowbreak Cusp}. Let $(a,\pi,d)$ be a part in
the actual isobaric parameter list of a point of $Z_i$. The
normalized trace of its Speh block,
\[
       \sum_{k=0}^{d-1}s_\pi(w)q_w^{(d-N)/2-k},           \tag{DFD.6}
\]
is the cofinite split Frobenius trace of an actual continuous
representation over $\overline{\mathbb Q}_{\ell_i}$.

The infinitesimal parameter of the whole list has distinct
entries in $\mathbb Z+(N-1)/2$. The $k=0$ submultiset of the
chosen part is its infinitesimal parameter translated by
$(d-1)/2$. As a submultiset of the whole parameter it too has
distinct entries. Therefore the infinitesimal weights of
$\pi$ lie in $\mathbb Z+(N-d)/2$ and are distinct. The
distinguished character $\mu=\texttt{mu\allowbreak Eta}(i)$ has weights in
$\mathbb Z+(N-t)/2$. Since $r=k+m-1$ is odd, write
\[
                    N=t+2u+1.                           \tag{DFD.7}
\]
The fixed character $\chi_1$ has Frobenius value
$s_\mu(w)q_w^{(t-N)/2}$; the cyclotomic character has value
$q_w$, and $\chi_{c+1}=\chi_1\operatorname{cyc}^{-c}$.

If $a+d-N=2e+1$, then
$(N-d)/2=(a-1)/2-e$, so $\pi$ is regular algebraic in the
normalization of DFD4. Twist its Galois representation by the
integer Tate power $\operatorname{cyc}^{e-k}$. Its trace has
exponent
\[
                 (1-a)/2+e-k=(d-N)/2-k.
\]
If instead $a+d-N=2e$, twist the automorphic part first by
$\mu$. Its infinitesimal weights differ from $(a-1)/2$ by
integers, since the sum of the two shifts differs by
$u+1-e$. Translation preserves distinctness, and the original
conjugate-self-duality law preserves that condition. Hence
$\pi\otimes\mu$ is an allowed input to DFD4. Twist its
Galois representation by
\[
                \chi_1^{-1}\operatorname{cyc}^{e-u-1-k}.
\]
The factor $s_\mu(w)$ cancels and the exponent becomes
\[
 (1-a)/2+(N-t)/2+e-u-1-k
          =e-a/2-k=(d-N)/2-k.                            \tag{DFD.8}
\]
Every Tate exponent used in both cases is an integer.
The nonzero complex Satake parameter of the character makes
the cancellation legitimate. Powers of the positive real
$q_w$ have the usual additive exponent law, agreeing with the
source complex-power convention.

Finite direct sums over $k$ give (DFD.6), and finite direct
sums over parts give sums of these characters. Tensoring by a
continuous one-dimensional character preserves continuity.
All finite exceptional sets may be united as in DFD3. This is
the full \texttt{has\allowbreak Trace\_\allowbreak part} assertion, including the even
parity case; it requires no half-integral Tate character.
% END DFD5

% BEGIN DFD6
\paragraph{DFD6: the actual block complement at every specified point.}
\label{claim:DFD6}
Under \texttt{Cheb} and \texttt{Galois\allowbreak Reps\allowbreak Cusp}, the original
\texttt{Complement\_\allowbreak at} holds. Namely, for every $i$, every
$\pi\in Z_i$, and every supplied Hecke Galois representation
$R_\pi:G_E\to\operatorname{GL}_N(O_i)$, there is an actual
$\rho_\pi:G_E\to\operatorname{GL}_r(O_i)$ satisfying all
three fields of \texttt{Is\allowbreak Complement}.

Use the literal parameter list provided by the $Z_i$ membership
dictionary. It contains the distinguished part $(1,\mu,t)$.
Remove one occurrence of that part, retaining any other equal
occurrences, and let $T$ be the sum of the representations of
DFD5 for the remaining list. At a permitted Frobenius place the
original Hecke polynomial identity, the unramified sum law for
the isobaric parameter, and the coefficient-of-degree-$N-1$
calculation of DFD4 give
\[
 \iota_{K_i,\overline{\mathbb Q}_{\ell_i}}
            \operatorname{tr}R_\pi(g)
       =T(g)+\sum_{c=1}^{t}\chi_{i,c}(g).                \tag{DFD.9}
\]
Indeed the removed part contributes
$s_\mu(w)\sum_{k=0}^{t-1}q_w^{(t-N)/2-k}$, which is exactly
the displayed character sum. This verifies both its indexing
from $1$ to $t$ and its negative cyclotomic powers.

The left side of (DFD.9) is continuous under the supplied
$\ell_i$-adic coefficient embedding; the right side is
continuous by DFD5 and the given characters. DFD3 extends the
identity to all $g\in G_E$. Apply DFD1 and DFD2 with
$a=r$, $O=O_i$, and $K=K_i$. Their rational isomorphism is
precisely the source \texttt{Iso\allowbreak Block}; semisimplicity is that
of its complete block representation, and continuity is that
of the actual integral complement. Neither a trace-only
surrogate nor a representation over a different DVR is used.

The original \texttt{Complement\allowbreak F} and
\texttt{Complement\allowbreak Closed} follow with exactly their stated
quantifiers: apply this construction at each later member of
the given family, and then at every member of the given family
class. Their premises already supply \texttt{Galois\allowbreak Reps\allowbreak Cusp}
on that scope; no uniform choice of bases across distinct
coefficient rings is claimed or needed.
% END DFD6

% BEGIN DFD7
\paragraph{DFD7: one continuous point family, including points outside its support.}
\label{claim:DFD7}
Assume \texttt{Cheb}, \texttt{Galois\allowbreak Reps\allowbreak Cusp}, and
$\forall i\,\exists R_i\,\texttt{Good\allowbreak R}(D,i,R_i)$.
Choose these $R_i$ once. For each $\pi\in Z_i$ choose the
complement of DFD6 to $R_i(\pi)$. For $\pi\notin Z_i$ put
$\rho_i(\pi)(g)=1_r$ for every $g$.

Every $\rho_i(\pi)$ is continuous: this is DFD2 on the support
and continuity of a constant off the support. On $Z_i$ it is
the actual complement to the already chosen $R_i(\pi)$.
Thus the same two families $(R,\rho)$ satisfy the three
literal fields of \texttt{Pt\allowbreak Fam}:
\[
 \begin{gathered}
   \forall i\quad\texttt{Good\allowbreak R}(D,i,R_i),\\
   \forall i,\pi\quad\rho_i(\pi)\text{ has continuous matrix entries},\\
   \forall i,\pi\in Z_i\quad
       \texttt{Is\allowbreak Complement}(R_i(\pi),\rho_i(\pi)).
 \end{gathered}                                         \tag{DFD.10}
\]
This proves the whole original
\texttt{exists\_\allowbreak continuous\_\allowbreak complement\allowbreak Family} and
\texttt{exists\_\allowbreak pt\allowbreak Fam}. Classical choice is only over the
specified location and point sets. Continuity here is in the
Galois variable, so choosing the families does not create an
unproved assertion of continuous dependence on $i$ or $\pi$.
% END DFD7

\subsection{The literal canonical data record}

% BEGIN DFD8
\paragraph{DFD8: all nine canonical row-data laws.}
\label{claim:DFD8}
Fix $L:\texttt{Size\allowbreak Data}_2(D)$, a supplied
$HW:\texttt{HWOf}(D.\mathrm{loc})$, the row-constant record
$R=(C_2,C_B,C_C,C_A,C'_B,C_K)$, and a \texttt{Pt\allowbreak Fam}
$(R_i,\rho_i)$. This paragraph's letter $R_i$ denotes the
representation; $R$ alone denotes the row-constant record.
Put
\[
 \begin{split}
 K_1&=C_B+48,\quad K_2=6C_B+12+6C_C,\\
 K_3&=3C_AC_B+12C_A+9,\qquad h=\lfloor2b\rfloor+1,\\
 c_B(i)&=\texttt{c\allowbreak Bnum}(\ell_i,\lfloor2b\rfloor,t,4nn_V),\\
 c_0(i,v)&=5+\lfloor\log_{\ell_i}(2h)\rfloor
       +4nn_V\bigl(4+2\lfloor\log_{\ell_i}(2h)\rfloor\bigr).
 \end{split}                                           \tag{DFD.11}
\]
In the last line the floors denote the source natural-number
logarithm, so the formula has its prescribed integer meaning.
Choose $\epsilon_J$ to be the witness of
\texttt{Conclusion\allowbreak Ja} when that proposition holds, and choose
the Boolean value true otherwise. Define $dc_S(i)$ to be a
witness of $\texttt{hyp\allowbreak Decl}(i,dc,\epsilon_J)$ when both
\texttt{Split\allowbreak Row\allowbreak Prime} and such a witness hold; otherwise take
all places over $\ell_i$ as rows and no columns. This is exactly
$\texttt{dc\allowbreak Of}$ and is defined without assuming these
propositions true.

For arbitrary $e'(i,v)$, finite $\mathrm{exc}\subset\mathbb N$,
and $I:\texttt{Tw\allowbreak In}(D)$ construct $X$ as follows. Keep
$X.r=(R_i)_i$, $X.\rho=\rho$, $X.e'=e'$, $X.c_B=c_B$,
$X.c_0=c_0$, $X.\mathrm{exc}=\mathrm{exc}$,
$X.\eta=\lfloor L.\mathrm{etaSz}\rfloor$, and $X.I=I$.
The band is the actual predicate
$\texttt{band\allowbreak Of}(L,HW,C_B,i,R_i,\rho_i,v)$, and
$X.e_S(i)=\texttt{e\allowbreak A}(L,C_A,\ell_i)$. Write
\[
 \begin{split}
 B_i&:\quad\ell_i<\texttt{l\allowbreak A}(L)
                    \ \text{and}\ 4L_w<\texttt{ord\allowbreak H}(D,\ell_i),\\
 U_i&=\begin{cases}
 \texttt{len\allowbreak Row\allowbreak B}(L,C_2,\ell_i),&B_i,\\
 \texttt{len\allowbreak Row\allowbreak CD}(L,X.\mathrm{band}_i,K_1,K_2,K_3,\ell_i),&\neg B_i.
 \end{cases}
 \end{split}                                           \tag{DFD.12}
\]
Set $X.\mathrm{lenRow}_i=U_i$. If $\ell_i$ is inert in the
source sense, use the all-row declaration, put
$X.\mathrm{len0}(i,\pi,v)=U_i(v)$, and use
$\texttt{local\allowbreak Error\allowbreak Bound\allowbreak Decomposition}(i,X.dc_i,c_B(i))$
for the error. Otherwise use $X.dc_i=dc_S(i)$,
$\texttt{len\allowbreak A}(i,X.dc_i,X.e_S(i),v)$ for $\mathrm{len0}$,
and the actual
$\texttt{local\allowbreak Error\allowbreak Bound\allowbreak Places}(i,\rho_i,\pi,X.dc_i,c_0(i),c_B(i))$
for the error. Rational row lengths are cast to real numbers
only in the two real-valued fields.

These formulas verify each of the nine \texttt{Data\allowbreak AD} laws.
The inert declaration has empty column set. Its $B_i$ and
$\neg B_i$ branches give \texttt{row\allowbreak B} and \texttt{row\allowbreak CD}.
The inert length and error definitions give \texttt{len0I}
and \texttt{error\_\allowbreak bound\_\allowbreak inert\_\allowbreak eq}. The formula for $e_S$
gives \texttt{e\allowbreak S}. A split prime cannot be inert: the
original prime-class definition assigns disjoint classes to
these predicates. Hence the noninert length and error branches
give \texttt{len0S} and \texttt{error\_\allowbreak bound\_\allowbreak split\_\allowbreak eq}.
The actual band formula is the ninth law. On every noninert
location $X.dc=dc_S$, as required separately by
\texttt{Canon}. Together with (DFD.10) and the three exact
equalities for $c_B,c_0,\eta$, this proves all six fields of
the original \texttt{Canon}. No row bound or crystalline
comparison is inferred merely from these definitions.
% END DFD8

% BEGIN DFD9
\paragraph{DFD9: the whole original canonical-data existence theorem.}
\label{claim:DFD9}
Assume \texttt{Cheb}, $\texttt{adm2}(D)$, and
$\texttt{Conclusion41}(D)$. For every original $L,HW,R$,
every function
\[
 e':\texttt{RFam}(D)\to\texttt{Rho\allowbreak Fam}(D)\to
          \prod_i\bigl(\operatorname{Places}(E)\to\mathbb N\bigr),
 \quad
 \mathrm{exc}:\texttt{Rho\allowbreak Fam}(D)\to\operatorname{Finset}(\mathbb N),
                                                               \tag{DFD.13}
\]
and every $I:\texttt{Tw\allowbreak In}(D)$, there is an actual $X$ with
\[
 \begin{gathered}
 \texttt{Canon}(L,HW,R,X),\qquad X.e'=e'(X.r,X.\rho),\\
 X.\mathrm{exc}=\mathrm{exc}(X.\rho),\qquad X.I=I,
 \qquad\forall i\ \texttt{Good\allowbreak R}(D,i,X.r_i),\\
 0\le L.\mathrm{etaSz}\ \Longrightarrow\
          X.\eta\le2L.\mathrm{etaSz},\\
 \forall i\quad
 \texttt{c\allowbreak Bnum}(\ell_i,\lfloor2b\rfloor,t,4nn_V)\le X.c_B(i).
 \end{gathered}                                         \tag{DFD.14}
\]
DFD4 supplies \texttt{Galois\allowbreak Reps\allowbreak Cusp}. For every $i$, the
first existence clause of \texttt{Conclusion41} supplies
some $R_i$ with \texttt{Good\allowbreak R}. DFD7 chooses one pair of
families $(R,\rho)$. Only after that choice evaluate the
functions in (DFD.13), obtaining $e'(R,\rho)$ and
$\mathrm{exc}(\rho)$. Apply DFD8 with these values.
Its formulas prove the first four assertions in (DFD.14),
and its \texttt{Pt\allowbreak Fam.\allowbreak good} proves the fifth. For the sixth,
$\lfloor a\rfloor_{\mathbb N}\le a\le2a$ when $a\ge0$.
The final inequality is equality by (DFD.11).

This order of construction preserves the original functional
dependence of the error and exceptional set on the chosen
representations. No prime condition is added to the finite
set of natural numbers. This is the full source theorem
$\texttt{Global\allowbreak Theory.\allowbreak Hub\allowbreak X.\allowbreak exists\_\allowbreak hub\allowbreak Data}$, with every
original input and output clause.
% END DFD9

\subsection{The same rows, twists, and coefficient lattices}

% BEGIN DFD10
\paragraph{DFD10: an actual good row for each good point.}
\label{claim:DFD10}
The original \texttt{tie0}, \texttt{Conclusion42}, and
\texttt{Conclusion\allowbreak Z} imply \texttt{Good\allowbreak Pt\allowbreak Row}. In fact
only the last input and the built-in $D$ dictionaries are
needed for this implication.

For $\pi\in Z_i^{\mathrm{good}}$, the literal good-point
dictionary supplies a parameter list $\Sigma$ with
$\texttt{Has\allowbreak Param\allowbreak X}(i,\pi,\Sigma)$ and
$\texttt{good\allowbreak Sig}(\Sigma)$. Its actual Hecke eigenspace
$\mathrm{eigC}(i,\pi)$ is nonzero: the witnessing
automorphic representation has nonzero type space, and that
type space lies in the supremum defining this eigenspace.
The \texttt{Conclusion\allowbreak Z} equality identifies this eigenspace
with the image of
\[
 M_{i,\Sigma}
    =\bigvee_{\sigma:\,\texttt{Has\allowbreak Param\allowbreak W}(i,\sigma,\Sigma)}
                      \mathrm{isot}(j_i,\sigma).         \tag{DFD.15}
\]
Therefore $M_{i,\Sigma}\ne0$. If every displayed isotypic
space were zero, its supremum would be zero. Choose one
$\sigma$ with \texttt{Has\allowbreak Param\allowbreak W} and nonzero isotypic space.

At each designated type place $w_{k,v}$ the original
nonconjugacy condition implies it is split. Explicitly, if
conjugation carried its prime ideal into itself, maximality
would make the ideal equal to its conjugate, contrary to
nonconjugacy; this is the split predicate. The split local
identity in \texttt{Has\allowbreak Param\allowbreak W} consequently identifies the
label of $\sigma$ at every such place with the label computed
from $\Sigma$. Unfolding the two goodness predicates now
gives
$\texttt{good\allowbreak Row}(\sigma)\Longleftrightarrow
\texttt{good\allowbreak Sig}(\Sigma)$, one label condition at a time.
Thus the chosen $\sigma$ is good. Together with its nonzero
isotypic space and the same $\Sigma$, this proves every field
of \texttt{Good\allowbreak Pt\allowbreak Row}. If \texttt{Triples\allowbreak Good} is also
supplied, apply it to this very pair $(\sigma,\Sigma)$ to
obtain a good signature attached to $\pi$, with all the
original shared parameters unchanged.
% END DFD10

% BEGIN DFD11
\paragraph{DFD11: the exact chosen-twist killing implication.}
\label{claim:DFD11}
For the original $D$, $C:\texttt{NChoice}(D)$,
$I:\texttt{Tw\allowbreak In}(D)$, and $\rho:\texttt{Rho\allowbreak Fam}(D)$,
\[
 \texttt{NCh\allowbreak Tie}(D,C)\ \wedge\ \texttt{Tw\allowbreak Good\allowbreak At}(D,C,I)
 \ \wedge\ \texttt{Signature\allowbreak Choice\allowbreak Condition}(\rho,C)
       \ \Longrightarrow\ \texttt{Kill\allowbreak At}(\rho).        \tag{DFD.16}
\]
Fix a twisted location $i$ and put $\ell=\ell_i$.
The tied twist index has prime $\ell$, and its error budget
is the original $\texttt{tw\allowbreak Eps}=\texttt{eps\allowbreak L}(\ell)=1/\ell^2$.
The \texttt{Tw\allowbreak Good\allowbreak At} data give nonnegative real numbers
$\epsilon_s$ on the finite good-signature set, with
\[
 \sum_{s\text{ good}}\epsilon_s\le1/\ell^2,
 \qquad
 \mathrm{sLen}(s)\le
          \operatorname{ofReal}(\nu+\epsilon_s),         \tag{DFD.17}
\]
using the exact oriented character in the source. More
explicitly, put
\[
 \psi_W=(\texttt{psi\allowbreak At}(C.\mathrm{sd},C.\psi'))
                       \circ\mathrm{art}.
\]
Both the bound in \texttt{Good\allowbreak Tw.\allowbreak part\allowbreak AT} and the bound in
\texttt{Signature\allowbreak Choice\allowbreak Condition} use
$\texttt{psi\allowbreak At}(\mathrm{side}(s),\psi_W)$.
They therefore compare the identical character directly.
The two Boolean sides need not agree, and no cancellation
of these orientations is used.

For $\pi\in Z_i^{\mathrm{good}}$, the supplied signature
condition chooses a good $s$ with
\[
 \operatorname{ofReal}(\mathrm{sPi}(i,\pi))
    \le\mathrm{sLen}(s),
 \qquad\mathrm{sLen}(s)\ne\infty
                \Longrightarrow\mathrm{sFin}(i,\pi).
\]
The right side of (DFD.17) is finite, so $\mathrm{sFin}$
holds. Every $\epsilon_s$ is bounded by the sum in
(DFD.17), because all summands are nonnegative. Since
$\nu+\epsilon_s\ge0$, comparison of the finite extended
nonnegative reals gives
\[
       \mathrm{sPi}(i,\pi)
          \le\nu+\epsilon_s\le\nu_L(D,\ell)+1/\ell^2.     \tag{DFD.18}
\]
Here $\nu=\nu_L(D,\ell)$ is the exact host cyclotomic
identity in the original \texttt{Nu\allowbreak Host} definitions, not
the weaker logarithmic bound used for DNE2. Explicitly,
unfold \texttt{nu\allowbreak L}, select its prime branch using the
primality of the tied twist prime, and the remaining
expression is definitionally \texttt{Gen.\allowbreak nu}. Thus this
step uses no separate numerical estimate.

For the second field put
\[
 d_Q=\sum_{v\in Q_i}
       \mathrm{lenUnit}_{O_i,\ell}
                       (O_i/(\zeta_{i,v}-1)).
\]
The signature condition gives
$d_Q\le\sum_{P\in C.Q(w_i)}\mathrm{tameVal}(\psi_W,P)$.
The last \texttt{Tw\allowbreak Good\allowbreak At} inequality bounds $t$ times
this same tame sum by $1/\ell^2$. Multiplication by the
nonnegative integer $t$ gives $td_Q\le1/\ell^2$, exactly
the second field of \texttt{Kill\allowbreak At}. This proves the full
original \texttt{twist\_\allowbreak kills\_\allowbreak selmer} theorem. Producing
\texttt{Signature\allowbreak Choice\allowbreak Condition} for the chosen $\rho$
is a separate comparison problem and is not inferred from
the existence of an attached signature alone.
% END DFD11

% BEGIN DFD12
\paragraph{DFD12: finite coefficients from the original integral crystalline input.}
\label{claim:DFD12}
The source predicate \texttt{Bhatt\allowbreak Scholze\allowbreak Integral} implies
the complete predicate \texttt{Bhatt\allowbreak Scholze\allowbreak Crystalline\allowbreak Lattices}.
The retained input is the former predicate, including its
fully faithful assertion; no integral comparison theorem
is newly assumed in the conclusion.

Fix the original prime $\ell$, local base field, prism,
$A_{\inf}$ identification and standard period data. Let
$L/\mathbb Q_\ell$ be any finite coefficient extension and
$T\subset V$ an $\mathcal O_L$-lattice in the given
continuous crystalline $L$-representation. Restrict the
same subset $T$ to a $\mathbb Z_\ell$-module. The ring
$\mathcal O_L$ is finite over $\mathbb Z_\ell$, so $T$ is
finite over $\mathbb Z_\ell$. It is torsion-free, as it
lies in a $\mathbb Q_\ell$-space, and hence free over the
PID $\mathbb Z_\ell$.

It spans $V$ over $\mathbb Q_\ell$. To see the point
without changing the lattice, every $c\in L$ has the form
$c=q a$, $q\in\mathbb Q_\ell$ and $a\in\mathcal O_L$:
multiply by a sufficiently large power of $\ell$ to make
it integral. Therefore the $\mathbb Q_\ell$-span of $T$
is stable under $L$. Since its $L$-span is $V$, this
$\mathbb Q_\ell$-span is already $V$. Stability under the
Galois action is unchanged. The finite-dimensional
$L$-module topology is the same module topology after
restriction to $\mathbb Q_\ell$, so every original orbit
map remains continuous.

For completeness, crystalline admissibility also survives
this restriction. In the source period-space definition,
$D_{\mathrm{cris}}$ is cut out by the same equations
$(g-1)x=0$ before and after restriction. Thus
\[
 D_{\mathrm{cris}}(\operatorname{Res}V)
   =\operatorname{Res}D_{\mathrm{cris}}(V).              \tag{DFD.19}
\]
Each dimension over $\mathbb Q_\ell$ is
$[L:\mathbb Q_\ell]$ times its dimension over $L$.
Substitute this on both sides of the source crystalline
dimension equality (including its unchanged base-period
degree), then cancel the positive factor
$[L:\mathbb Q_\ell]$. This proves the exact
$\texttt{is\allowbreak Crys\_\allowbreak res\allowbreak Q\_\allowbreak iff}$, rather than assuming it.

Apply \texttt{Bhatt\allowbreak Scholze\allowbreak Integral} to this restricted
representation and this identical stable lattice. It
returns the required $F$-crystal, its semilinear realization
map, and its realization law with target
$W(C^\flat)\otimes_{\mathbb Z_\ell}T$. These are already
the output fields of the coefficient version. Its second,
unique-morphism field is literally the second field of
the retained integral predicate, so it passes through
unchanged. This completes the finite-coefficient theorem.
Constructing the compatible Wach functor over all layers
and its \texttt{Cmp\allowbreak Of} ties requires further work; it is
not a consequence of scalar restriction alone.
% END DFD12

\subsection{One eventual threshold and its exact remaining inputs}

% BEGIN DFD13
\paragraph{DFD13: all original parameter inequalities at one uniform cutoff.}
\label{claim:DFD13}
For every $\epsilon>0$ and real $A$ there is $x_0$ such
that every original $(D,L,P)$ with scale exponent
$\epsilon$, $P.C_{L_w}\le A$, and $x\ge x_0$ satisfies
both \texttt{Adm\allowbreak Par} and
\texttt{Parameters\allowbreak Polynomially\allowbreak Bounded}. All objects
$(D,L,P)$ are quantified after this one cutoff.

Write $b=e^{x^\epsilon}$,
$b'=xb/e^{x^{\epsilon/3}}$,
$w''=b'/(xn)$, $S=L.\mathrm{Sc}$, and
$\eta=L.\mathrm{etaSz}$. The original parameter fields
give $t\le\eta$, $t\ge7$, $n\ge1$,
$kS\le w''$, $xt\eta\le w''$ and
$k\ge m\sqrt{x}$ with $m\ge2$.
For $x\ge10000$, these imply
\[
 200S\le w'',\qquad70000\eta\le w'',\qquad
 100\le w'',\qquad100(\eta+S+t+2)\le w'',
 \qquad w''\le b'.                                    \tag{DFD.20}
\]
For the fourth inequality, $t\le\eta$ and $\eta\ge7$
give $\eta+t+2\le(16/7)\eta$; its contribution after
multiplication by $100$ is less than $w''/2$ by the
second inequality, while $100S\le w''/2$.

Here is also the top-weight inequality, with its actual
parameter inputs. Put
$T_0=\texttt{Nstar}(L)x^\epsilon\mathrm{LaG}/k$.
The field \texttt{rho\_\allowbreak le}, after dropping its other
nonnegative term, says $4T_0\le x^{-2}$. Nonnegativity
of that term follows from $\mathrm{LaG}>0$, which itself
follows from $4S\le\mathrm{LaG}\,b'/k$ and $S>0$.
The field \texttt{b\allowbreak T\_\allowbreak le} gives $b_a\le b'(1+T_0)$.
For $x\ge4$,
\[
        b_a+b'/x\le b'(1+1/64+1/4)<3b'/2.              \tag{DFD.21}
\]
Finally, uniformly in $P.C_{L_w}\le A$,
$L_w\le\max(A,0)\log x$. The elementary limit
\[
 x(\max(A,0)\log x+1)e^{-x^{\epsilon/3}}\longrightarrow0
\]
provides a cutoff for
\[
              100(L_w+1)b'\le b.                       \tag{DFD.22}
\]
To justify that limit directly, put $y=x^{\epsilon/3}$.
The expression is bounded for $x\ge1$ by a constant
times $y^{6/\epsilon}e^{-y}$; expand $e^{y/2}$ in its
power series to dominate any fixed power of $y$, leaving
a constant times $e^{-y/2}\to0$.

Equations (DFD.20)--(DFD.22) are all fields of the source
$\texttt{Roomy}(L,100)$. They give
$\eta+t+S<w''$, $3w''\le b$, every
$b_a+b'/x\le b$, and $t\le\eta$, exactly
\texttt{Adm\allowbreak Par}. They also give $S+\eta+t+2\le b$,
exactly \texttt{Parameters\allowbreak Polynomially\allowbreak Bounded}. Taking
the maximum of the numerical and exponential cutoffs
proves the full original
\texttt{admissible\allowbreak Parameters\_\allowbreak of\_\allowbreak large\_\allowbreak scale} theorem.
% END DFD13

% BEGIN DFD14
\paragraph{DFD14: the full family small-prime assertion.}
\label{claim:DFD14}
For an original family $F$ and $C\ge1$, there is one
$x_0$ such that, for every $x\ge x_0$, every original
$h_x:F.x_1<x$, and every $i\in F.I_x$, one has
\[
 \texttt{Small}(L,C),\qquad8\le L.\mathrm{sc}.L,
 \qquad60\le D.\mathrm{host}.\mathrm{splitBound},
 \qquad C\le b.                                        \tag{DFD.23}
\]
The proof below verifies all fields of \texttt{Small},
including the prime-three and two auxiliary-prime clauses.

Set $z=x^\epsilon$, $b=e^z$, $p_*=\texttt{l\allowbreak Star}(L)$,
$A_*=\texttt{l\allowbreak A}(L)$, and use the original rank symbols
$f,r,m,t$. The family size bounds give, for $x\ge3$,
\[
 k\le x^5,\quad f\le k,\quad r\le2k,
 \quad t\le x^{|C_t|+1},\quad
 frt\le2x^{11+|C_t|}.                                  \tag{DFD.24}
\]
The exponent $C_t$ is the one fixed by the family, before
its members. Choose a cutoff, depending only on the
family constants and $C$, such that
\[
 \begin{gathered}
 z\ge\max(16,C),\qquad x^{\epsilon_z/5}\ge60,\\
 4\cdot2^C x^{(11+|C_t|)C+\epsilon}\le e^{x^\epsilon},
 \qquad32x^{5+\epsilon}\le e^{x^\epsilon}.
 \end{gathered}                                         \tag{DFD.25}
\]
Every one of these conditions is eventually true since
$\epsilon,\epsilon_z>0$. For the exponential comparisons
use the same power-series argument as DFD13. They are
uniform in every family member.

The original scale and parameter definitions give
\[
 A_*\ge z,\quad S,\eta,t\le b',\quad
             rb'A_*\log x\le mb.                       \tag{DFD.26}
\]
For clarity these do not require a new asymptotic result:
$A_*=kz/(m\delta)$, $k\ge m$, $0<\delta\le1$ give the
first; $t\eta x\le w''\le b'$, $t\le\eta$, and
$kS\le w''$ give the middle bounds for $x\ge1$; and
the field \texttt{rb\allowbreak S}, after dropping its nonnegative
$S$ term and using $\delta\le1$, gives the last.
Since $\log x\ge1$ for $x\ge3$, (DFD.25)--(DFD.26)
give
\[
 16rb'\le mb,\quad16rz\le mb,\quad
 r(S+t+\eta+z)\le mb/4,\quad(frt)^Cz\le b/4.            \tag{DFD.27}
\]
The second follows from $r\le2x^5$ and the final
inequality in (DFD.25); the first uses $A_*\ge16$.

For either auxiliary prime $p=q,d$, the host inequality
$p>4tr+1$ gives $2r\le p$. Consequently $rb/p\le b/2$.
Using $f,m\ge1$, the actual expression
\[
 \texttt{mid\allowbreak G}(L,C,p)
   =Cf\{r(b/p+S+t+\eta+z)+mb\}+C(frt)^Cz
\]
is at most
$Cf(7mb/4)+C(b/4)\le2Cfmb$, by (DFD.27).
This is precisely the auxiliary-prime field of
\texttt{Small}.

It remains to locate the exceptional small primes. From
\texttt{lstar1} one has $b/p_*<p_*$, hence
$\sqrt b\le p_*$. The original conductor condition is
$M^{36}t^{150}\le\sqrt b$, where $M=Q_Sqd\ge1$.
Every prime in $D.\mathrm{host}.\mathrm{small}$ and
both $q,d$ divide the positive integer $M$. Thus each
is at most $M\le p_*$. The original \texttt{lstar2}
also gives $7<p_*$, in particular $3\le p_*$.

The split threshold in (DFD.25), together with the actual
floor tie for \texttt{split\allowbreak Bound}, gives
$60\le\texttt{split\allowbreak Bound}$. Hence the original host
small-splitting law applies to $3$. This prime differs
from $q$ by the host record and from $d$ because
$d>4tr+1>21$. If it belongs to the host's small set, its
assigned class is noninert; otherwise the just obtained
split law assigns its noninert split class. Thus
$\neg\texttt{Inert\allowbreak K}(D,3)$. We have proved every small,
auxiliary, prime-three, and \texttt{mid\allowbreak G} clause.
Finally $L.\mathrm{sc}.L=z\ge16$ and
$b=e^z\ge z\ge C$. This proves (DFD.23) and the full
original \texttt{special\_\allowbreak primes\_\allowbreak small} family theorem.
% END DFD14

% BEGIN DFD15
\paragraph{DFD15: the original family setup reduced to two actual suppliers.}
\label{claim:DFD15}
Fix the original $F:\texttt{Fam\allowbreak P}(\epsilon)$, constants
$C,C_2,CT$, and the original \texttt{Tw\allowbreak Good} and
\texttt{Facts\allowbreak F} premises. Retain \texttt{Cheb} and the
following two explicit supplier statements on these same
data:
\begin{enumerate}
\item For each relevant $D$ there exist actual $HW,F_W$
and $\texttt{Cmp\allowbreak Of}(D.\mathrm{loc},D.\mathrm{locA},HW,F_W)$.
\item For each actual pair $(R,\rho)$ satisfying
\texttt{Pt\allowbreak Fam} on this $D$, the original facts imply
$\texttt{Signature\allowbreak Choice\allowbreak Condition}(\rho,C_D)$ for the
same prescribed $C_D=\texttt{nch}$.
\end{enumerate}
Then the full output of \texttt{exists\_\allowbreak setup\_\allowbreak large}
follows, with one $x_0$ before $x,h_x,i$, before the
$\texttt{Sha\allowbreak Fin\allowbreak Odd}$ proof, and before the arbitrary
finite exceptional set. This is a conditional reduction;
the two displayed statements are remaining human-proof
obligations, not additional consequences claimed here.

The nonempty index at scale zero shows $\epsilon>0$
from the family's scale tie. DFD13 gives a cutoff $x_a$
with $A=F.PK[7]$, because that is exactly the family
bound for $C_{L_w}$. DFD14 at constant $1$ gives $x_b$.
The original \texttt{Tw\allowbreak Good} premise gives $x_c$.
Choose $x_0=\max\{x_a,x_b,x_c\}$ once.

For a later $x,h_x,i$, let $D,L,C_D$ be its original
data. The twist premise chooses $N_0,S$ and gives
\texttt{Tw\allowbreak Good\allowbreak At} for
$I=(\texttt{tw\allowbreak E}(F,x,h_x,i),N_0,S)$. Keep this $I$.
The first remaining supplier chooses $HW,F_W,\mathrm{Cmp}$.
The \texttt{Facts\allowbreak F} record contains \texttt{adm2} and
\texttt{Conclusion41}, so DFD9 constructs $X$ with
row constants
\[
            \texttt{row\allowbreak C}(C_2)=(C_2,320,25,23,32,32),
\]
$e'=0$, exceptional-set function constantly equal to
the given $\mathrm{exc}$, and exactly this $I$.
Thus $X.\mathrm{exc}=\mathrm{exc}$ and
$X.I=I$. The second supplier applies to its actual
\texttt{Pt\allowbreak Fam}; DFD11 and the original \texttt{NCh\allowbreak Tie}
give $\texttt{Kill\allowbreak At}(X.\rho)$.

Now assemble the eight \texttt{Setup} fields.
The facts field is the given \texttt{Facts\allowbreak F} member;
the canonical field is DFD9; its two parameter fields
are DFD13; its $8\le L$ and split-bound fields are
DFD14; its killing field was just proved; and its last
field is the supplied \texttt{Sha\allowbreak Fin\allowbreak Odd}. The
\texttt{Adm\allowbreak J} conclusion follows from the same member's
original \texttt{adm\allowbreak J} implication applied to
\texttt{Adm\allowbreak Par}. All fields refer to the one chosen
$X,HW,F_W,\mathrm{Cmp},D,L,C_D$.

This proves the reduction with the exact existential and
universal order. DFD12 supplies the finite-coefficient
layer of the first remaining supplier from the original
integral input; DFD10 supplies actual attached good rows
for the second. Neither yet supplies all their coherent
comparison laws. Later local row estimates and the full
common \texttt{Concl\allowbreak F} also remain. Accordingly this
chapter does not close \texttt{GLOBAL\_\allowbreak DENOMINATOR}.
% END DFD15

\subsection{Source alignment and remaining scope}

The exact pinned source ranges and full-file hashes for
each DFD claim are in this component's
\texttt{source\_\allowbreak manifest.\allowbreak json} and \texttt{provenance.\allowbreak json}.
DFD1--DFD7 follow the actual interfaces in
\texttt{Semisimple\allowbreak Of\allowbreak Injective},
\texttt{Block\allowbreak Complement\allowbreak Lattice},
\texttt{Galois\allowbreak Reps\allowbreak Cuspidal}, and
\texttt{Exists\allowbreak Continuous\allowbreak Complement\allowbreak Family}; the saturated
lattice proof supplies a direct simplification of the
matrix construction. DFD8--DFD9 prove the literal
\texttt{Canon} and \texttt{Exists\allowbreak Hub\allowbreak Data} output.
DFD10--DFD12 correspond to \texttt{Local\allowbreak Global\allowbreak Assembly},
\texttt{Chosen\allowbreak Twist\allowbreak Annihilates\allowbreak Selmer\allowbreak Groups}, and
\texttt{Bhatt\allowbreak Scholze\allowbreak Coefficients\allowbreak Of\allowbreak Padic\allowbreak Int\allowbreak Lattices}.
DFD13--DFD14 supply the original parameter and small-prime
theorems. DFD15 records precisely what is still needed
to replace the conditional use of
\texttt{Exists\allowbreak Hub\allowbreak Data\allowbreak Family} by a complete human proof.

The earlier DNE chapter remains the full
$\texttt{Concl\allowbreak F}\Rightarrow\texttt{Adm\allowbreak F}\Rightarrow
\texttt{Link\allowbreak N}$ consumer. The present actual integral
families feed its upstream construction without changing
the host, coefficient DVRs, embeddings, chosen twists,
or exceptional sets. No formal replay is claimed.

\endgroup
\end{document}
